
==================================================
Filename: 00_main.tex
==================================================
\documentclass{article}
\usepackage[utf8]{inputenc}

\makeatletter
\def\blfootnote{\gdef\@thefnmark{}\@footnotetext}
\makeatother

\setlength{\parindent}{0em}
\setlength{\parskip}{0.65em}

%
\input{math_commands.tex}

\usepackage{minitoc}
\usepackage[bottom]{footmisc}


\usepackage{cancel}
\usepackage{hyperref}
\usepackage{url}
\usepackage{fancyhdr}
\usepackage{microtype}

\usepackage{subcaption}
\usepackage{dirtytalk}
\usepackage{booktabs}
\usepackage{enumitem}
\usepackage{float}
\usepackage{multirow}
\usepackage{multicol}
\usepackage{amsmath}
\usepackage{amssymb}
\usepackage{epigraph}
\usepackage{graphicx}
\usepackage{amsthm}
\usepackage{amsfonts}
\usepackage{mathtools}
\usepackage{fullpage}
\usepackage[nice]{nicefrac}
\usepackage{algorithm}
\usepackage{algorithmic}
\usepackage{xcolor}
\usepackage{natbib}

\usepackage{xfrac}
\usepackage{framed}
\usepackage{bbm}
\usepackage{comment}
\usepackage{thmtools, thm-restate}

\usepackage[toc,page,header]{appendix}
\usepackage{minitoc}


%
\renewcommand \thepart{}
\renewcommand \partname{}

\makeatletter
\setlength{\@fptop}{0pt}
\setlength{\@fpbot}{0pt plus 1fil}
\makeatother

%
%
%
\usepackage{hyperref}
\hypersetup{colorlinks,linkcolor=blue}

\usepackage[capitalize,noabbrev,nameinlink]{cleveref}

%

\definecolor{amaranth}{rgb}{0.9, 0.17, 0.31}
\definecolor{green}{HTML}{549D54}

\makeatletter
\setlength{\@fptop}{0pt}
\setlength{\@fpbot}{0pt plus 1fil}
\makeatother

\makeatletter
\renewcommand\section{\@startsection{section}{1}{\z@}%
                                    {-2.5ex \@plus -1ex \@minus -.2ex} %
                                    {1.5ex \@plus.2ex} %
                                    {\normalfont\Large\bfseries}}
\makeatother

\makeatletter
\renewcommand\subsection{\@startsection{subsection}{2}{\z@}%
                                    {-1.5ex \@plus -1ex \@minus -.2ex} %
                                    {1ex \@plus .2ex} %
                                    {\normalfont\large\bfseries}}
\makeatother

\title{Positional Attention: Expressivity and Learnability of Algorithmic Computation}

\author{
    Artur Back de Luca$^{1*}$\qquad
    George Giapitzakis$^{1*}$\qquad
    Shenghao Yang$^{1*}$\\
    Petar Veli\v{c}ković$^{2}$\qquad
    Kimon Fountoulakis$^{1}$\\
    \vspace{-1mm}\\
    \normalsize{$^1$University of Waterloo \quad $^2$Google DeepMind}\\
    \vspace{-3mm}\\
    \normalsize{\texttt{\{\href{mailto:abackdel@uwaterloo.ca}{abackdel}, \href{mailto:ggiapitz@uwaterloo.ca}{ggiapitz}, \href{mailto:shenghao.yang@uwaterloo.ca}{shenghao.yang}, \href{mailto:kimon.fountoulakis@uwaterloo.ca}{kimon.fountoulakis}\}@uwaterloo.ca}}\\
    \normalsize{\texttt{\href{mailto:petarv@google.com}{petarv@google.com}}}
}

\date{}

\newcommand{\fix}{\marginpar{FIX}}
\newcommand{\new}{\marginpar{NEW}}

%
\theoremstyle{plain}

\newtheorem{theorem}{Theorem}[section]
\newtheorem{proposition}{Proposition}[section]
\newtheorem{claim}{Claim}[section]
\newtheorem{lemma}{Lemma}[section]
\newtheorem{fact}{Fact}[section]
\newtheorem{corollary}{Corollary}[section]
\theoremstyle{definition}
\newtheorem{definition}{Definition}[section]
\newtheorem{assumption}{Assumption}[section]
\theoremstyle{plain}
\newtheorem{remark}{Remark}[section]

\newcommand{\FF}{\ensuremath{\mathbb{F}}}
\newcommand{\ZZ}{\ensuremath{\mathbb{Z}}}
\newcommand{\RR}{\ensuremath{\mathbb{R}}}
\newcommand{\CC}{\ensuremath{\mathbb{C}}}

\DeclareMathOperator{\spn}{span}
\DeclareMathOperator{\rank}{rank}
\DeclareMathOperator{\mlp}{\mathchoice{\mbox{MLP}}{\mbox{MLP}}{\mbox{{\tiny MLP}}}{\mbox{{\tiny MLP}}}}
%
\newcommand{\hardmax}{\mathrm{hardmax}}

\newcommand\hA{\hat{A}}
\newcommand\hX{\hat{X}}
\newcommand\hY{\hat{Y}}
\newcommand\hG{\hat{G}}
\newcommand\hw{\hat{w}}
\newcommand\mnorm[1]{\left\Vert#1\right\Vert_{2,\infty}}
\newcommand\norm[1]{\left\Vert#1\right\Vert_{2}}
\newcommand\normOne[1]{\left\Vert#1\right\Vert_{1}}
\newcommand\norminf[1]{\left\Vert#1\right\Vert_{\infty}}
\newcommand{\BigPhi}{\scalebox{1.25}[1]{$\Phi$}}

\usepackage[textsize=tiny]{todonotes}



%
\begin{document}
\maketitle

\def\thefootnote{*}
\footnotetext{Equal contribution.}
\def\thefootnote{\arabic{footnote}}

\vspace{-5mm}

\begin{abstract}
    There is a growing interest in the ability of neural networks to execute algorithmic tasks (e.g., arithmetic, summary statistics, and sorting).
    The goal of this work is to better understand the role of attention in Transformers for algorithmic execution.
    Its importance for algorithmic execution has been studied theoretically and empirically using parallel computational models. Notably, many parallel algorithms communicate between processors solely using positional information. Inspired by this observation, we investigate how Transformers can execute algorithms using positional attention, where attention weights depend exclusively on positional encodings. We prove that Transformers with positional attention (positional Transformers) maintain the same expressivity of parallel computational models, incurring a logarithmic depth cost relative to the input length. We analyze their in-distribution learnability and explore how parameter norms in positional attention affect sample complexity. Our results show that positional Transformers introduce a learning trade-off: while they exhibit better theoretical dependence on parameter norms, certain tasks may require more layers, which can, in turn, increase sample complexity. Finally, we empirically explore the out-of-distribution performance of positional Transformers and find that they perform well in tasks where their underlying algorithmic solution relies on positional information.
    %
    %
    %
\end{abstract}

%

\doparttoc %
\faketableofcontents %

%
\parttoc %

%

\input{01_intro}
\input{02_review}
\input{03_preliminaries}
\input{04_architecture}
\input{05_expressivity}
\input{06_learnability}
\input{07_experiments}
\input{08_limitations}
\input{10_acknowledgements}

\bibliography{references}
\bibliographystyle{plainnat}

\newpage
\appendix

\part{Appendix} %
\parttoc %

\input{991_mpc}
\input{992_expressivity}
\input{993_complexity}
\input{994_experiments}
\input{995_ood_probability}
\input{996_failure_attention}
\input{997_odd_bounds}



\end{document}



==================================================
Filename: 01_intro.tex
==================================================
\section{Introduction}

Neural network models, particularly Transformers \citep{vaswani2017attention}, have achieved remarkable success across various domains, especially in their ability to learn and execute algorithms \citep{velivckovic2021neural}. Notable applications 
range from basic arithmetic to solving problems on abstract data structures such as graphs and lists \citep{lee2024teaching, tang2020towards, tay2020sparse, yan2020neural, Veličković2020Neural, velivckovic2022clrs, dudzik2022graph, ibarz2022generalist, cappart2023combinatorial, diao2023relational, bevilacqua2023neural, rodionov2023neural, minder2023salsa, georgiev2023beyond, engelmayer2024parallel, georgiev2024deep, li2024openbook}. Due to this success, there is a growing interest in understanding neural networks as computational models that can solve complex algorithmic and computational tasks.

In this context, \citet{perez2021attention, wei2022statistically} show that Transformers are Turing Complete, and \citet{giannou23a, pmlr-v235-back-de-luca24a, yang2023looped} illustrate how Transformers can effectively encode instructions in their parameters to solve linear algebra and graphs problems. Additionally, it has been shown that Transformers can perform computational tasks using far fewer layers than the number of sequential steps \citep{liu2022transformers}, indicating a connection between Transformers and parallel algorithms. 
Building on this, \citet{pmlr-v235-sanford24a} demonstrates that Transformers can simulate the Massively Parallel Computation (MPC) model \citep{8555148}, which is based on the MapReduce framework for large-scale data processing \citep{10.1145/1327452.1327492}.

While parallel computational models like MPC do not restrict how communication is established, we observe that in many parallel algorithms, the communication between processors is independent of the processed data.
Instead, this communication relies solely on processor identification (see \Cref{app:illustration_algo} for examples). 
%
%
Motivated by this observation, we study a Transformer variant that reflects this data-agnostic communication property. Specifically, we investigate \emph{positional attention}, where the attention weights are determined exclusively by fixed positional encodings.\footnote{While positional encodings remain constant across layers, their attention weights can vary -- unlike Graph Neural Networks (GNNs), which take a fixed connectivity structure as input. Furthermore, the tasks in this paper lack a predefined graph structure, necessitating an artificial graph for application to GNNs. For experimental comparisons, see \Cref{appdx:experiments}.}


\textbf{Our contributions.} We examine Transformers with positional attention (positional Transformers) from
both a theoretical and empirical perspective in the context of algorithmic computation. Our contributions are given below:

\textit{Expressivity.} We prove that positional Transformers can simulate a single round of the Massively Parallel Computation (MPC) model with $O(\log n)$ layers, where $n$ is the input length. The expressivity result emerges from a connection between Transformers and parallel computational models. In particular, we utilize a proxy computational model, which only allows communication over a static network, bridging the simulation results between positional Transformers and MPC and, consequently, by \citet{pmlr-v235-sanford24a}, standard Transformers. Due to the restrictive nature of communication in positional Transformers and static networks, this equivalence with MPC is established with logarithmic cost. 
%

\textit{In-distribution generalization.} We provide a norm-based generalization bound for positional Transformers, which, combined with our expressivity result, implies that they can learn algorithmic tasks to any desired accuracy via empirical risk minimization. Our analysis reveals a trade-off between the number of layers and parameter norms: positional Transformers exhibit better dependence on parameter norms but may require more layers for certain tasks, potentially leading to higher overall theoretical complexity. This result is established using the Provably Approximately Correct (PAC) framework, where the complexity bounds of the hypothesis class are expressed as a function of the network parameter norms \cite{edelman2022inductive}. 
    
\textit{Out-of-distribution (OOD) generalization.} We empirically analyze the OOD performance of positional and standard Transformers. We observe that positional Transformers exhibit better OOD generalization for tasks where the underlying algorithmic solution relies solely on positional information. However, for tasks that do not meet this criterion, such as induction heads tasks, both models perform poorly. We \textit{hypothesize} that the inherently dynamic nature of communication in such tasks affects the loss landscape, making it challenging to discover good parameters for the positional Transformer during training.
    

%

%

\begin{figure*}[t]
    \centering
    \includegraphics[width=\linewidth]{plots/diagram.pdf}
    \caption{Diagram comparing the operations of self-attention in Transformers with positional attention. The figure illustrates a single attention head, but in multi-head attention, multiple sets of queries, keys, and values are processed in parallel and then combined. 
    In Transformers, the model's input, $X^{(0)}$, is a combination of input values $X$ and positional encodings $P$. In positional attention, however, these components are processed separately. At layer $\ell$, the query ($Q^{(\ell)}$) and key ($K^{(\ell)}$) are derived solely from the positional encodings $P$, where $P$ remains fixed across layers. These are multiplied (denoted by $\textsf{Mul}$) and passed through a softmax function to produce the attention matrix $A^{(\ell)}$. As in self-attention, the value $V^{(\ell)}$ in positional attention is computed from the previous layer's input, $X^{(\ell-1)}$. The attention matrix $A^{(\ell)}$ and the value $V^{(\ell)}$ are then multiplied to form the weighted representation, which is linearly transformed into the output $O^{(\ell)}$. This output is passed to a Multilayer Perceptron (MLP) for further processing, as detailed in \Cref{sec:arch}.}
    \label{fig:diagram}
\end{figure*}


%



==================================================
Filename: 02_review.tex
==================================================
\section{Related work}

\textbf{Empirical:} Several studies focus on the experimental aspects of training neural networks to execute algorithms.
Notable examples include algorithms operating on common data structures, such as lists (e.g., sorting, binary search) and graphs (e.g., shortest paths, connected components) \cite{yan2020neural, tang2020towards, velivckovic2022clrs, ibarz2022generalist, bevilacqua2023neural, engelmayer2023parallel, georgiev2023beyond, minder2023salsa, diao2023relational}. Additional research on combinatorial optimization tasks using graph neural networks and Transformers offers further insights into algorithmic learning in structured domains~\citep{vinyals2015pointer, prates2019learning, joshi2019efficient, bai2020learning, selsam2019guiding, karalias2020erdos, gasse2019exact, cappart2023combinatorial, velickovic2023everything}.
Recent studies have also explored the broader algorithmic capabilities of Transformer-based models, shedding light on their potential for general algorithmic reasoning~\citep{deletang2023neural, butoi2025training, barbero2025round}. A parallel stream of research has focused on performing arithmetic operations (e.g. addition, multiplication) \cite{kaiser2016neural, klindt2023controlling, shen2023positional, lee2024teaching, ruoss2023randomized}. Specifically, \citet{rodionov2023neural} leverages problem-specific information within a self-supervised training framework, whereas the other studies use intermediate labels to guide the learning process. For example, \citet{engelmayer2024parallel} shows that using intermediate labels derived from parallel algorithms leads to improved performance. While some research, such as \citet{kazemnejad2023the}, investigates the impact of different positional encodings across several tasks, to the best of our knowledge, no existing work examines the use of positional attention within the context of neural algorithmic reasoning. The closest formulation to the positional Transformer architecture presented in \Cref{sec:arch} is the position-based attention of \citet{schmidt2023bridging}, but it is used in the context of neural machine translation.

%
%
%

\textbf{Theoretical:} From a theoretical perspective, the most closely related work to ours is \citet{pmlr-v235-sanford24a}, which presents simulation results for standard Transformers and the Massively Parallel Computation (MPC) model. 
Our work also presents simulation results demonstrating that positional Transformers can simulate MPC. The simulation proof involves introducing a proxy computational model with an alternative communication paradigm that can be shown to simulate MPC. Subsequently, we establish that the positional Transformer architecture can simulate this proxy model and, by extension, MPC itself. %
Other relevant studies demonstrate the expressive power of neural networks through simulation results. For example, \citet{siegelman95comp} establishes the Turing completeness of recurrent neural networks (RNNs), while \citet{hertrich2023provably} presents specific RNN constructions that solve the shortest paths problem and provide approximate solutions to the knapsack problem. Additionally, other simulation results focused on Transformers have shown their Turing completeness \citep{perez2021attention, wei2022statistically} as well as demonstrated constructive solutions to linear algebra and graph-related problems \citep{giannou23a, pmlr-v235-back-de-luca24a, yang2023looped}. 
%
%
%

%

%
%
%

    
%
%
%



==================================================
Filename: 03_preliminaries.tex
==================================================
\section{Preliminaries and notation}\label{sec:prelim}

Let $\mathbb{N} = \{1,2,\dots\}$ denote natural numbers and $[n] = \{1,2,\dots, n\}$ for $n \in \mathbb{N}$. For a set $S$, let $\mathcal{P}(S)$ be its power set (i.e. the set containing all subsets of $S$). For a matrix $Z$, let $Z_{i,:}$ and $Z_{:,i}$ be the $i$-th row and column of $Z$, respectively. $\|\cdot\|_2$ denotes the spectral norm for matrices, and $\|\cdot\|_{p,q}$ denotes the $(p,q)$ matrix norm where the $p$-norm is over columns and $q$-norm over rows (e.g. $\|Z\|_{2,1} = \sum_{i}\|Z_{:,i}\|_2$). For vectors, $\|\cdot\|_p$ denotes the $\ell_p$ norm. For a sequence of vectors $X \in \mathbb{R}^{n \times d}$ we may use $X[i] = X_{i,:} \in \mathbb{R}^d$ for the $i$-th vector in the sequence.

\begin{definition}[Risk]
Let $\mathcal{D}$ be a distribution over $\mathcal{X}\times \mathbb{R}$, and let $\ell : \mathbb{R}\times\mathbb{R} \rightarrow \mathbb{R}$ be a loss function. For a given class of functions $\mathcal{F} = \{ f : \mathcal{X} \rightarrow \mathbb{R} \}$ and $f\in\mathcal{F}$, the {\em risk} (i.e. the generalization error)  is defined as $\mathsf{risk}(f;\mathcal{D}) = \mathbb{E}_{(x,y)\sim\mathcal{D}} \left[\ell(f(x),y)\right]$, and the {\em empirical risk} is $\widehat{\mathsf{risk}}(f;(x^{(i)},y^{(i)})_{i=1}^m) = \frac{1}{m}\sum_{i=1}^m \ell(f(x^{(i)}), y^{(i)})$.
\end{definition}



%
%
%
%
%
%
%
%
%
%
%


%

%
%
%

%

%


==================================================
Filename: 04_architecture.tex
==================================================
\section{The Positional Transformer architecture}
\label{sec:arch}
We now define the \emph{positional Transformer}, an adaptation of the Transformer model \citep{vaswani2017attention}.
This adaptation is motivated by the fact that communication between machines is independent of the data processed in many real-world parallel algorithms. In particular, we decouple the input values from the computation of attention weights.
%
For an input $X\in \RR^{n\times d_X}$, we define the $\ell^\text{th}$ layer as follows:
\begin{equation}
\label{eq:arch}
\text{F}^{(\ell)}(X) = \BigPhi^{(\ell)}\!\left(\!\left(\bigoplus_{h=1}^HA^{(\ell, h)}XW^{(\ell, h)}_V\right)W^{(\ell)}_O \!\oplus X\right).
\end{equation}

The input is processed by $H$ attention heads, each associated with an attention weight matrix $A^{(\ell, h)} \in (0,1)^{n \times n}$ and a value matrix $W_V^{(\ell, h)} \in \R^{d_X \times d_V}$. Here, $\ell$ denotes the layer index and $h$ the head index, allowing a specific attention head within a layer to be identified as $(\ell, h)$. The outputs of these attention heads are concatenated and then transformed by an output matrix $W_O^\ell \in \RR^{H \cdot d_V \times d_O}$. This result is concatenated with a residual connection of the input $X$ and then passed through a multilayer perceptron (MLP), represented as $\Phi^{(\ell)}: \RR^{d_O + d_X} \rightarrow \RR^{d_{\text{out}}}$.

Traditionally, in standard Transformers, the attention matrix $A^{(\ell, h)}$ is derived using self-attention. 
\begin{equation}
\label{eq:selfattn}
A^{(\ell, h)}(X) = \softmax\left(\!\!\left(XW_Q^{(\ell, h)}\right)\!\cdot\!\left(XW_K^{(\ell, h)}\right)^{\!\!\top}\right).
\end{equation}
Here, the attention weights are derived by the input matrix $X$, using query and key matrices $W^{(\ell, h)}_Q, W^{(\ell, h)}_K \in \RR^{d_X\times d_m}$, where $d_m$ is the embedding dimension.

In contrast, we utilize \emph{positional attention}, where attention weights are learned solely using positional encodings $P$, which are constant across all layers.
This distinction is also illustrated in \Cref{fig:diagram}.
\begin{equation}
\label{eq:attn}
    A^{(\ell, h)} = \softmax\left(\left(PW^{(\ell, h)}_Q\right)\cdot \left(PW^{(\ell, h)}_K\right)^\top\right).
\end{equation}
We utilize node positional encodings defined by a matrix $P\in\RR^{n\times d_P}$ and whose attention weights are computed similarly to traditional attention, using query and key matrices but with input dimension $d_P$. The positional encodings are fixed across layers, indicated by the absence of the $\ell$ index.

Theoretically, we evaluate whether these changes reduce the expressive power of positional Transformers. We prove that positional Transformers, like standard Transformers, can simulate the Massively Parallel Computation (MPC) model. This indicates that there is no loss in expressivity, though with an added logarithmic increase in depth relative to the input length, as discussed further in the next section.

%
%

\begin{comment}
\subsection{Embedding dimension requirements}
In this section, we further evaluate the computational complexity of the standard attention and our architecture from an expressivity perspective, focusing on the dependence of the embedding dimensions $d_T$ and $d_E$ on the size of the input $n$. Specifically, we show that both architectures require an embedding dimension of the order of $n$ to be expressive enough. We can therefore conclude that the computational complexity of our architecture is superior to that of standard attention, without suffering from limited expressive power. Note that expressivity here measures the number of relevant attention patterns that can be produced by an architecture and should not be confused with the class of problems it can solve. By imposing sparsity in the attention mechanism this class is naturally restricted.
\subsubsection{Standard attention}
Let's first investigate what would be a reasonable requirement for the embedding dimension $d_T$ in standard self-attention. The attention scores in this case are calculated by 
$$a(x_i, x_j)= \sigma((W_K x_i)^T(W_Q x_j))$$
where $x_i$ and $x_j$ are input tokens, $W_K$ and $W_Q$ are the key and query matrices respectively and $\sigma$ is the activation function. By setting $H = W_K^TW_Q$ we can rewrite 
$$a(x_i, x_j) = \sigma(x_i^T H x_j)$$
Notice that 
$$x_i^T H x_j = \sum_{k=1}^{d_T}\sum_{l=1}^{d_T}
(x_i)_k (x_j)_l H_{kl}$$
is a linear function of $\text{vec}(H) \in \RR^{{d_T}^2}$ 
By relaxing $\sigma$ to be the hardmax function the calculation of attention scores can be relaxed to a binary classification problem of $n^2$ points in a Euclidean space of dimension $d_T^2$. Expressivity is measured in terms of the number of possible binary patterns achievable for a given matrix $H$. Cover's theorem \cite{Cover1965GeometricalAS} requires the dimension of the space to be at least as large as the number of points for a classifier to be able to produce all patterns, which, in this case, gives a lower bound on the embedding dimension $d_T \geq n$.
\end{comment}


==================================================
Filename: 05_expressivity.tex
==================================================
\section{Expressivity of Positional Transformers}
\label{sec:expressivity}
A popular approach to illustrate the expressivity of an architecture is by establishing an equivalence with a computational model \cite{pmlr-v235-sanford24a, perez2021attention, loukas2019graph}. For positional Transformers, we will utilize the Massively Parallel Computation (MPC) model, which was first used by \citet{pmlr-v235-sanford24a} to analyze standard Transformers. We begin by stating our main result.

\begin{restatable}{theorem}{expressivity}
\label{thm:expressivity}
Consider an instance of MPC $\mathsf{T}$ with $R$ rounds, $N$ machines with local memory $s$. Let $\mathcal{M}$ be a model following the architecture in \Cref{eq:arch} with $n=N+1$ nodes, $2R\lceil\log N\rceil$ layers and $2s$ attention heads. Then, for any instance $\mathsf{T}$ with Borel measurable local functions, there exists a configuration of $\mathcal{M}$ that approximates $\mathsf{T}$ to any desired degree of accuracy.
\end{restatable}

\emph{Proof overview:} We establish this result using a two-step procedure.
This procedure involves introducing a proxy computational model which, in contrast to MPC, uses a static network for communication. We call this model Parallel Computation with Oracle Communication (PCOC)\footnote{We note that PCOC is only used to establish our expressivity result. It is not meant as a model to compete with established parallel models, such as MPC.}, and its formal definition is presented in \Cref{app:background}.
In the first step of the proof, we demonstrate that the computations over a static network model such as PCOC can effectively simulate the dynamic communication of MPC. Specifically, we show that $R$ rounds of an MPC protocol can be simulated using $2R\lceil \log n \rceil$ rounds of PCOC, where $n$ represents the number of processors in MPC. The additional computation rounds and communication coordinated by Bene\v{s} networks \cite{benevs1965mathematical} are sufficient to simulate MPC.

In the second step of the proof, we show that PCOC can be, in turn, simulated by positional Transformers. Specifically, the attention mechanism is shown to simulate communication between nodes, while the computation stage leverages universal approximation results for MLPs \cite{cybenko1989approximation, hornik1989multi}. By demonstrating our architecture’s capability to approximate both communication and local computation, we establish that it can simulate any MPC instance of $R$ rounds using $2R\log n$ layers.

By Theorem 3.4 in \citet{pmlr-v235-sanford24a}, which demonstrates the simulation of standard Transformers by MPC, we can derive the following remark regarding the simulation of standard Transformers by positional Transformers. %

\begin{restatable}{remark}{equivalence}
\label{rem:expressivity}
Consider a standard Transformer $\mathcal{M}^\prime$ with $N$ nodes, $L$ layers, and $d_V\!\cdot\!H$ sublinear in $N$. Then, there exists a positional Transformer $\mathcal{M}$ with $O(N^2)$ nodes, $O(L\log N)$ layers, and $d_V\!\cdot\!H$ sublinear in $N$ that can simulate the computation of $\mathcal{M}^\prime$.
\end{restatable}

%

Note that such a simulation result introduces an additional quadratic dependency on the number of nodes in positional Transformers. These dependencies are merely a byproduct of the simulation strategy employed in \citet{pmlr-v235-sanford24a}, and other, more efficient, simulation strategies may exist.
Additionally, this quadratic cost arises only when simulating the exact same sequence of operations of the standard Transformer. In practice, the model may converge to parameter configurations that do not correspond to the standard Transformer. When solving a given problem, positional Transformers may adopt a completely different configuration to achieve the solution without incurring this additional cost, as demonstrated by our experiments in \Cref{sec:experiments}.


\begin{comment}
We present two lemmas that bridge positional Transformers and MPC. To prove these results, we utilize a proxy model, which has a static communication network. We call this proxy model Parallel Computation with Oracle Communication (PCOC). \textit{We note that PCOC is only used to establish our expressivity result. It is not meant as a model to compete with established parallel models, such as MPC}. We begin by outlining its main features, followed by the two lemmas that lead to our final result. The formal definition of the PCOC model is presented in \Cref{app:mpc}.





In the second step, we show that these computations over the static network can, in turn, be simulated by positional Transformers. Specifically, the attention mechanism is shown to simulate communication between nodes, while the computation stage leverages universal approximation results for MLPs (Cybenko, 1989; Hornik et al., 1989). We prove that  $\Phi_\ell$ can approximate each machine’s local computations, with unique node identifiers ensuring injectivity in the MLP approximation. This approach ensures that local functions remain distinct, even when input rows share values. By demonstrating our architecture’s capability to approximate both communication and local computation, we establish that it can simulate any MPC instance of  R  rounds using  2R\log n  layers






%
%

%

%

%

%

%
%

\textbf{Parallel Computation with Oracle Communication (PCOC).}
The PCOC model operates on $n$ machines labeled from $0$ to $n-1$ over $R$ rounds, with each machine having a local memory of size $s$. The model consists of two fundamental steps per round: communication and computation. During communication, machines exchange data through an oracle that provides destination information for each machine and message, executing a communication pattern that is agnostic to the processed data and depends only on the input length. In the computation step, machines perform local computations on data stored in their memory. Critically, the oracle's communication pattern remains fixed for a given input length and task, meaning that if the values at each machine change, the communication pattern between machines remains the same, a property that is observed in many parallel algorithms such as computing basic statistics or sorting, where communication is determined by positions rather than values. 

This contrasts with other parallel models, such as MPC, where the communication pattern is more flexible and directly influenced by the data being processed. In the MPC model, each machine has a limited memory size of $s$, and the input data is distributed among the machines. The computation proceeds in rounds: during each round, a machine processes the data in its memory, exchanges messages with other machines (up to a total size of $s$), and retains only the messages it received for the next round. After all rounds, the final output is stored in the memory of one or more machines. Unlike PCOC, the communication in MPC is not predetermined by the input length or task alone; instead, it can dynamically depend on the data being processed, provided the total communication per machine per round does not exceed the memory limit $s$.

\subsection{PCOC can simulate MPC}
This section presents an important insight: even with limited communication capabilities, PCOC can perform the same functions as MPC, incurring only a logarithmic cost in rounds relative to the number of machines, $N$. Below, we present our results and a high-level overview of the proof, with the formal proof presented in \Cref{app:proof_expr}.

\begin{restatable}{lemma}{MPCPCOC}
\label{lem:pcoc_mpc}
Consider an instance of MPC $\mathsf{T}$ with $R$ rounds, $N$ machines with local memory $s$. 
For any instance $\mathsf{T}$, there exists an instance of PCOC $\mathsf{P}$ with $2R\lceil\log N\rceil$ rounds, $N$ machines with local memory $2s$ that simulates $\mathsf{T}$.
\end{restatable}

\emph{Proof overview:} without loss of generality, we consider a single round of MPC communication. Our primary goal is to demonstrate that any MPC communication instance can be simulated by PCOC, as their local computation capabilities are equivalent. While PCOC's communication remains data-agnostic, it achieves communication flexibility through the Bene\v{s} network - a communication network that guarantees data transmission between any two processors in $2\lceil\log N\rceil$ rounds ($N$ being the number of processors), with each processor communicating with only two others per round \cite{benevs1965mathematical}.
The PCOC protocol's intermediate communication comprises two phases: fragmentation and recombination. During fragmentation (first $\log N$ rounds), processors fragment and distribute messages to neighboring processors, with local computation managing message bit placement in designated slots for subsequent communication rounds. This ensures that after $\log N$ rounds, each processor has data from all its network neighbors. In the recombination phase (final $\log N$ rounds), local functions reorganize data in local memory based on intended destinations. This crucial step enables PCOC to replicate any valid MPC communication pattern by strategically rearranging received message bits according to their destinations, which the oracle then utilizes to transmit data in a fixed communication pattern based solely on memory bit positions.

%
An analogy to illustrate the distinction between the two approaches can be drawn from a mailing system. In the MPC model, each processor acts as an individual delivering letters directly to their recipients. Communication is flexible and determined by the specific needs of the task. In contrast, PCOC resembles a structured mailing system with predefined routes and sorting stations. Each individual sends their letters to the nearest sorting station, where the letters are organized into bins and passed through successive stages of sorting and delivery. After several rounds, all letters reach their destinations, following a fixed communication pattern regardless of the specific contents.


\subsection{Positional Transformers can simulate PCOC}
Having established the PCOC model, we now show that the positional Transformer model in \Cref{eq:arch} can simulate it.
More specifically, our results show that a $R$-round PCOC protocol can be simulated by a positional Transformer with $R$ layers. We first present the corresponding theorem, followed by a proof overview. The details of the proof are presented in \Cref{app:expressivity}.

\begin{restatable}{lemma}{PCOC}
\label{lem:pcoc_sim}
Consider a PCOC instance $\mathsf{P}$ with $R$ rounds, $N$ machines with local memory $s$, and data type $\mathbb{T}=\mathbb{R}$. Let $\mathcal{M}$ be a model following the architecture in \Cref{eq:arch} with $n=N+1$ nodes, $R$ layers and $s$ attention heads. Then, for any instance $\mathsf{P}$ with Borel measurable local functions, there exists a configuration of $\mathcal{M}$ that approximates $\mathsf{P}$ to any desired degree of accuracy.
\end{restatable}

\emph{Proof overview:} The proof demonstrates that a single layer of the positional Transformer can simulate each PCOC round through two key components: communication and computation. In the communication stage, the oracle specifies which machines receive information from which memory positions, represented through $s$ positional attention heads where each head corresponds to a local memory position. The model requires $N+1$ nodes (one more than the number of machines) to accommodate a sink node for cases where no information is sent, as attention matrices cannot contain zero rows. For the computation stage, we leverage universal approximation results of MLPs \citep{cybenko1989approximation, hornik1989multilayer} to show that the MLP in \Cref{eq:arch} can approximate each machine's local computations, with unique node identifiers ensuring injectivity of the MLP approximation and maintaining distinct local functions even when input rows have similar values. By establishing our architecture's ability to approximate any oracle and local functions, we prove it can simulate any PCOC instance.

%
\emph{Proof overview:} The proof starts by demonstrating that a single layer of the positional Transformer can simulate each individual round of PCOC. The constructive proof can be further divided into two main components: communication and computation.

\textbf{Communication:} The communication stage leverages the oracle $\mathsf{RCV}_{n,\mathtt{Data}}$ to specify the subsets of machines and local memory positions from which each machine receives information.
These subsets can be transformed into binary encodings, which are represented by distinct attention heads, one for each local memory position, for a total of $s$ positional attention heads. This part of the proof relies on the capability of attention to represent binary patterns, which is shown in \Cref{appdx:hardmax}. 
It is important to note that the number of nodes exceeds the number of machines by one, as an additional node is necessary to represent unsent messages. In the computation of attention, no attention matrix can contain a row of all zeros, implying that a machine is not receiving any information. Due to such cases, we introduce an additional sink node to account for information not directed to any machine. Therefore, in this framework, we can show that any communication pattern defined by the oracle can be effectively represented by $s$ attention heads across $N+1$ nodes.

\textbf{Computation:} This stage accounts for the local computations executed by each machine. To this end, we invoke the universal approximation results of multilayer perceptrons (MLPs) \citep{cybenko1989approximation, hornik1989multilayer} to establish that, in each round, the local computations initiated by each machine can be approximated by $\Phi^{\ell}$ as detailed in \Cref{eq:arch}. One important consideration is the inclusion of unique node identifiers in the input to ensure the injectivity of the MLP approximation. 
Even if two input rows may have the same input values, the unique identifiers guarantee that each row corresponds to a distinct local function.
Furthermore, the node identifiers must be preserved at every layer to maintain this injectivity. This is achieved by the MLP when processing both the output of the attention heads and the residual input $X$, ensuring that identifiers are consistently retained.

Therefore, by demonstrating that our architecture can approximate any oracle and local functions, we show its ability to simulate any algorithm in PCOC. In practice, finding an oracle and local functions for a specific task can be posed as a learning problem. Our proposed architecture adopts this approach and can learn to execute parallel algorithms using fixed positional encodings in the attention mechanism. As our experiments illustrate, this approach helps mitigate OOD generalization issues. In our experiments we do not assume access to explicit supervision for communication or computation, both are learned indirectly through ground truth, as all models are trained end-to-end.
%

\subsection{Positional Transformers can simulate MPC}
Having established the two important lemmas \Cref{lem:pcoc_mpc} and \Cref{lem:pcoc_sim}, we are ready to state our main result.

\begin{restatable}{theorem}{expressivity}
\label{thm:expressivity}
Consider an instance of MPC $\mathsf{T}$ with $R$ rounds, $N$ machines with local memory $s$. Let $\mathcal{M}$ be a model following the architecture in \Cref{eq:arch} with $n=N+1$ nodes, $2R\lceil\log N\rceil$ layers and $s$ attention heads \textcolor{red}{Should this be $2s$?}. Then, for any instance $\mathsf{T}$ with Borel measurable local functions, there exists a configuration of $\mathcal{M}$ that approximates $\mathsf{T}$ to any desired degree of accuracy.
\end{restatable}

The proof of this result follows directly from the combination of the two preceding results. Notably, the structure of MPC enforces that local memory can only grow sublinearly relative to the input size. This constraint prevents trivial simulation configurations where all data is sent to a single node, which would undermine the purpose of parallelism.

Moreover, by Theorem 3.4 in \cite{pmlr-v235-sanford24a}, which demonstrates the simulation of standard Transformers by MPC, we can derive the following simplified remark regarding the equivalence between standard and positional Transformers. A more formal version of this result is provided in \Cref{app:mpc}.

\begin{restatable}{remark}{equivalence}
\label{rem:expressivity}
Consider a standard Transformer $\mathcal{M}^\prime$ with $N$ nodes, $L$ layers, and memory and total embedding dimension (including the number of attention heads) that grow sublinearly with $N$. Then, there exists a positional Transformer $\mathcal{M}$ with $2L\lceil\log N\rceil$ layers and $O(N^2)$ nodes that can simulate the computation of $\mathcal{M}^\prime$.
\end{restatable}

Note that such a simulation result introduces an additional quadratic dependency on the number of nodes in positional Transformers. These dependencies are merely a byproduct of the simulation strategy employed in \cite{pmlr-v235-sanford24a}, and other, more efficient, simulation strategies may exist.
Additionally, this quadratic cost arises only when simulating the same sequence of operations across architectures. In practice, the model may converge to parameter configurations that do not correspond to interpretable algorithms. When solving a given problem, positional Transformers may adopt a completely different configuration to achieve the solution without incurring this additional cost, as demonstrated by our experiments in \Cref{sec:experiments}.

%

%


%
\end{comment}


==================================================
Filename: 06_learnability.tex
==================================================
\section{Learnability of Positional Transformers}
\label{sec:learnability}
In this section, we present a norm-based generalization bound for the positional Transformer architecture, discuss its implications, and compare it with the previously derived bound for the Transformer architecture that employs self-attention~\citep{edelman2022inductive}. In particular, by decoupling the attention weights from the input $X$ and requiring them to depend exclusively on the positional encodings $P$, we remove the exponential dependence (w.r.t. the depth $L$) on the norms of key and query weight matrices.

We consider multi-layer networks obtained by composing individual layers from \Cref{eq:arch}, denoted as
\begin{equation}\label{eq:multi_layer_arch}
  F^{1:L}(X) = F^{(L)} \circ \cdots \circ F^{(2)} \circ F^{(1)}(X).
\end{equation}
The resulting network $F^{1:L}$ has $L$ layers and each layer has $H$ attention heads. For simplicity we assume that the MLP $\Phi^{(\ell)}$ consists of two layers:
\begin{equation}\label{eq:phi}
  \Phi^{(\ell)}(Z) = \sigma(ZW_1^{(\ell)})W_2^{(\ell)}, ~ \forall \ell \in [L],
\end{equation}
where $\sigma$ is an $L_\sigma$-Lipschitz activation with $\sigma(0)=0$. Using MLPs with more layers would lead to a worse dependence on the spectral norm of weight matrices, but does not affect the results in this section.\footnote{More layers in the MLP would simply increase the constant $c$ in \Cref{thm:generalization} to $c' \ge c$. On the other hand, two-layer MLPs are already sufficient to invoke the universal approximation properties required by the expressivity result of \Cref{thm:expressivity}.} Following prior work, our generalization bound is derived for scalar-valued functions. We extract a scalar output from our architecture as $w^\top X^{\mathsf{out}}[n]$ with trainable weights $w \in \mathbb{R}^{d_\mathrm{out}}$, where $X^{\mathsf{out}} = F^{1:L}(X) \in \mathbb{R}^{n \times d_\mathrm{out}}$. The index $i$ at which we apply the linear function $x \mapsto w^\top x$ can be arbitrary, and here we simply set to the last index. This setting captures algorithmic computations whose output is a scalar. For comparison purposes, it aligns with the setup considered previously by \citet{edelman2022inductive} for the analysis of self-attention Transformers. In general, for computational tasks that require vector or matrix output, one can easily extend our result to parallel architectures using a union bound on the failure probability $\delta$ and incur an additional logarithmic cost with respect to the dimension of the output.

Given $B_{2,1},B_2 \ge 1$, the class of $L$-layer $H$-head bounded-norm scalar-output positional Transformer networks is:
\begin{align*}
  &\mathcal{F} = 
  \bigg\{ 
  w^\top \left(F^{1:L}(X)[n]\right): \|w\|_2 \le B_2, ~ \mbox{and} ~ \forall \ell,h:\\  &\|W{^{(\ell,h)}_K}^\top W{^{(\ell,h)}_Q}\|_{2,1}, \|W_V^{(\ell,h)}\|_{2,1}, \|W_1^{(\ell)}\|_{2,1} \le B_{2,1},\\
  &\|W_2^{(\ell)}\|_{2,1} \le B_{2,1},
  \|W_V^{(\ell,h)}\|_2, \|W_1^{(\ell)}\|_2, \|W_2^{(\ell)}\|_2 \le B_2
  \bigg\},
\end{align*}
with $F^{1:L}$ from \Cref{eq:multi_layer_arch}, $A^{(\ell,h)}$ from \Cref{eq:attn}, and $\Phi^{(\ell)}$ from \Cref{eq:phi}. In order to avoid overloading the result with unnecessary notation that do not add value to the discussion in any meaningful way, we assume that the input $X \in \mathbb{R}^{n \times d_X}$ and positional encodings $P \in \mathbb{R}^{n \times d_P}$ are normalized row-wise so that $\|X_{i,:}\|_2 = \|P_{i,:}\|_2 = 1$ for all $i\in[n]$. We denote $d = \max\{d_X, d_P, d_V, d_\mathrm{out}\}$, and therefore $d(H+1)$ is the largest possible size of any weight matrix along any axis.

Let $\mathcal{D}$ be a distribution on $\mathbb{R}^{n\times d_X} \times \mathbb{R}$ and $(X^{(1)},y^{(1)}),$ $\ldots,$ $(X^{(m)},y^{(m)}) \in \mathbb{R}^{n\times d_X} \times \mathbb{R}$ be a sequence of $m$ samples. 

\begin{theorem}[Generalization bound]\label{thm:generalization}
For any $\delta>0$, with probability at least $1-\delta$, simultaneously for all $f \in \mathcal{F}$, it holds that for some $c>0$,
\begin{align*}
 &\left|\mathsf{risk}(f;\mathcal{D}) - \widehat{\mathsf{risk}}(f;(X^{(i)},y^{(i)})_{i=1}^m)\right|\\
 &\le \tilde{O} \bigg((HL_\sigma B_2)^{cL}B_{2,1}\sqrt{\frac{\log(Hdmn)}{m}} + \sqrt{\frac{\log(1/\delta)}{m}}\bigg).
\end{align*}
\end{theorem}

In the context of learning algorithmic computations, \Cref{thm:generalization} means that the generalization gap $|\mathsf{risk}(f)-\widehat{\mathsf{risk}}(f)|$ goes to 0 as the number of samples tends to infinity. On the other hand, by \Cref{thm:expressivity}, for any computational task that is computable in an MPC model with Borel measurable local functions, there is $f \in \mathcal{F}$ (where the parameters $L, H, B_2, B_{2,1}$ that specify $\mathcal{F}$ might depend on $n$ but are independent of $m$) such that the empirical risk can be arbitrarily close to 0 for all samples of size $m$ and for all $m\ge0$, i.e. $\widehat{\mathsf{risk}}(f;(X^{(i)},y^{(i)})_{i=1}^m) \approx 0,$ $\forall m$. Consequently, in this case, \Cref{thm:generalization} implies that as the number of samples $m$ tends to infinity, we have $\mathsf{risk}(f;\mathcal{D}) \rightarrow 0$. This means that positional Transformer learns algorithmic computations through empirical risk minimization.

Compared with the norm-based risk bounds of self-attention Transformer networks in \citet{edelman2022inductive}, the bound in \Cref{thm:generalization} does not depend on the spectral norm of key and query weight matrices. Intuitively, this is due to restricting attention weights to be determined exclusively by positional encodings, as opposed to allowing them to be determined compositionally by the output from the previous layer. If we replace positional attention with self-attention in the definition of $\mathcal{F}$, then the bound in \Cref{thm:generalization} would incur an additional factor $B_{QK}^{O(L)}$ where $\|W{^{(\ell,h)}_{K}}^\top W{^{(\ell,h)}_{Q}}\|_2 \le B_{QK}$, $\forall \ell,h$. This shows a potential advantage of positional attention when all other model configurations are the same.

Note that achieving a low empirical risk with self-attention or positional attention can require very different model configurations, which may depend on the computational task. For example, in the simple task of computing the minimum of an $n$-number array, a standard Transformer can solve it approximately using just one attention layer with an error controlled by how well softmax approximates hardmax. In contrast, a positional Transformer can solve it exactly but requires $\log n$ layers with two attention heads per layer (e.g., via binary tree reduction). Here, the positional Transformer’s worst-case sample complexity depends on $n$, while the standard Transformer does not. However, the standard Transformer can have a much worse dependence on the magnitude of the input numbers due to the softmax-hardmax approximation. Therefore, one architecture might be better depending on the magnitude of $n$ and the input numbers. For more complex tasks, setting an ``optimal'' model configuration (e.g., choosing $L,H,d_V$) is difficult.
%

%



==================================================
Filename: 07_experiments.tex
==================================================
\section{Experiments}
\label{sec:experiments}
In this section, we evaluate the performance of positional and standard Transformers across various algorithmic tasks and regimes.
We first outline the tasks considered in this work and then describe the experimental setup in detail. Next, we present results for the in-distribution regime across different settings, followed by an analysis of the out-of-distribution performance.
Finally, we introduce a more complex task that combines textual and numerical data.
For additional analyses, including comparisons with other models and ablation studies, refer to \Cref{appdx:experiments}.

%

%
%
%
%
%
%

\begin{figure*}[h!]
    \centering
    \includegraphics[width=\linewidth]{plots/fixed_length/all_in_dist/sample.pdf}
    \caption{In-distribution loss across all five tasks for standard Transformers (red) and positional Transformers (blue) as a function of the number of training samples (indicated on the x-axis).}
    \label{fig:fix_sample}
\end{figure*}

\textbf{Tasks:} We consider the following 
algorithmic tasks:

\begin{enumerate}[leftmargin=.4cm]
    \item \emph{Cumulative sum}: Given $x \in \mathbb{R}^n$, output $y \in \mathbb{R}^n $ where each element $ y_i $ is the sum of the first $ i $ elements of $ x $, i.e. $y_i = \sum_{j=1}^{i} x_j$.
    \item \emph{Cumulative min}: Given $x \in \mathbb{R}^n$, output $y \in \mathbb{R}^n$ where each element $ y_i $ is the minimum value among the first $i$ elements of $x$, i.e. $y_i = \min \{ x_j \mid 1 \leq j \leq i \}$.
    \item \emph{Cumulative median}: Given $x \in \mathbb{R}^n$, output $y \in \mathbb{R}^n$ where each element $y_i$ is the median of the first $ i $ elements of $x$, i.e. $y_i = \text{median} \{x_j \mid 1 \leq j \leq i\}$.
    \item \emph{Sorting}: Given $x \in \mathbb{R}^n$, output $\text{sort}(x)$, a vector containing the entries of $x$ sorted in ascending order.
    \item \emph{Cumulative maximum sum subarray}: given $x \in \mathbb{R}^n$, output $y \in \mathbb{R}^n$ where each element $y_i$ is the sum of elements of a maximum sum subarray of the first $i$ elements of $x$, i.e. $y_i = \max_{1 \leq j \leq k \leq i} \left( \sum_{l=j}^{k} x_l \right)$. 
    %
\end{enumerate}

The tasks selected were chosen to represent varying levels of complexity and illustrate the strengths and limitations of each architecture. We adopt cumulative versions of algorithms when feasible for several reasons: They naturally provide $n$ to $n$ training settings and are more challenging than non-cumulative versions. %

To test the non-numeric capabilities of positional Transformers, we further present two additional tasks that utilize textual data. The first is the $k$-hop inductive heads of \citet{pmlr-v235-sanford24a},
which represents a higher-order version of the standard inductive heads task by recursively executing the completion of a bigram to determine the subsequent bigram. The second task employs the same computational tasks described earlier in the list, but incorporates additional textual data that serve as categories, which the model must learn to use for conditional reasoning. 

Although both tasks involve pattern matching, the nature of pattern matching differs. The $k$-hop induction heads task requires dynamic pattern matching, where each step depends on previous matches, whereas the mixed-type input task involves static pattern matching, which better aligns with our architecture. Despite our expressivity result guaranteeing that this task can be solved using a static network, we \textit{hypothesize} that the inherently dynamic nature of communication in this task is reflected in the loss landscape. Consequently, discovering good parameters for the positional Transformer is difficult. This is evident in our experiments, where we run multiple trials, all of which result in poor out-of-distribution performance.

\textbf{Experimental setting:} All tasks employ the same model structure of \Cref{eq:arch}, augmented with linear encoding and decoding layers. We compare the standard Transformer, which utilizes the attention mechanism in \Cref{eq:selfattn}, and the positional Transformer, which employs the attention defined in \Cref{eq:attn}. Standard Transformers incorporate positional encodings concatenated with the input values, unlike positional Transformers, where positional information is exclusively supplied through the matrix $P$. In \Cref{appdx:experiments}, we examine other configurations for standard Transformers, including relative positional encodings such as Rotary Positional Embedding (RoPE) \cite{su2024roformer} and find no significant differences in performance. Both variants share the same number of layers and dimensional configurations, with any specific differences explicitly noted. For details on the particular layer configurations and training details, refer to \Cref{appdx:experiments}. The training data is sampled in the range $[-2, 2]$. To ensure diversity in the data, for each input sample, we first select lower and upper bounds $\gamma_l$ and $\gamma_u$ uniformly in $[-2, 2]$, and then for each of the $n$ elements of the input sample, we select its value uniformly in  $[\gamma_l, \gamma_u]$. 

\begin{figure}[h!]
    \centering
    \includegraphics[width=\linewidth]{plots/induction_heads/indheads_id.pdf}
    \caption{In-distribution loss for standard Transformers (red) and positional Transformers (blue) in the $k$-hop induction heads task of \citet{pmlr-v235-sanford24a}, plotted as a function of the number of training samples (left, with eight layers) and the number of layers (right, with $50,\!000$ training samples).}
    \label{fig:indheads_id}
\end{figure}

\begin{figure*}[t!]
    \centering
    \includegraphics[width=\linewidth]{plots/fixed_length/all/scale.pdf}
    \caption{In-distribution and out-of-distribution loss for standard Transformers (red) and positional Transformers (blue) across all five tasks. The x-axis shows the OOD scale factor, with a scale factor of one indicating in-distribution performance.}
    \label{fig:fix_scale}
\end{figure*}

\subsection{In-distribution regime}
\label{sec:id-regime}
We evaluate our architecture in a fixed input length regime, with a detailed analysis of in-distribution generalization as a function of the number of samples and expressive power. The plots presented show the median over ten runs, with shaded areas representing the 10\textsuperscript{th} and 90\textsuperscript{th} percentile.


%
%
%

\textbf{Sample Size vs. Generalization:} In this setting, we demonstrate the learnability of positional Transformers, established in \Cref{thm:generalization}, as a function of the number of samples. To this end, we fix the input length $n = 8$ and examine in-distribution loss as a function of the number of training samples, ranging from $5,\!000$ to $50,\!000$.
\Cref{fig:fix_sample} shows that for all tasks, the in-distribution loss of both architectures consistently decreases with increasing number of samples.
%
%


\textbf{Expressive power vs. Generalization:} 
In this experiment, we empirically validate the theoretical results of \Cref{rem:expressivity}, analyzing the in-distribution capabilities of standard and positional Transformers as a function of their expressive power, measured by the number of layers. We evaluate their performance on a generalized version of the induction heads task \cite{olsson2022context}, which reflects the relational capabilities of self-attention.
This task, named \emph{$k$-hop induction heads}, applies a successive $k$ number of bigram completions to determine the next bigram in the sequence. We adapt the setting of \citet{pmlr-v235-sanford24a}, restricting the number of hops to at most three and the sequence length to 30. For a complete description of the task and additional experimental details, we refer readers to \citet{pmlr-v235-sanford24a} and \Cref{appdx:experiments}, respectively. We chose this task because it is designed to highlight the efficiency of self-attention in terms of the number of required layers. The in-distribution results in \Cref{fig:indheads_id} show that standard Transformers exhibit near-zero in-distribution loss with as few as three layers. In contrast, we find that positional Transformers generally perform worse. As established in \Cref{sec:experiments}, this results from the expressivity of the positional Transformers since achieving a level of performance comparable to that of the standard Transformers requires an increasing number of layers, aligning with our theoretical results.

\begin{figure}[h!]
    \centering
    \includegraphics[width=0.6\linewidth]{plots/induction_heads/indheads_ood.pdf}
    \caption{Out-of-distribution loss for standard Transformers (red) and positional Transformers (blue) in the $k$-hop induction heads task of \citet{pmlr-v235-sanford24a}, plotted as a function of the number of layers (with $50,\!000$ training samples).}
    \label{fig:indheads_ood}
\end{figure}


%
%


\begin{figure*}[t!]
    \centering
    \includegraphics[width=0.9\linewidth]{plots/nlp/alltasks_mse_noise_p_and_s.pdf}
    \caption{In-distribution and out-of-distribution losses for standard (red) and positional Transformers (blue) across mixed-type input task variations. The x-axis shows the OOD scale factor, where one indicates in-distribution performance.}
    \label{fig:nlp_noise_p_and_t}
\end{figure*}

\subsection{Out-of-distribution regime}
\label{sec:ood-regime}

We further evaluate the capabilities of both architectures in generalizing out-of-distribution (OOD) for the tasks previously outlined. To generate OOD data, we employ a similar sampling strategy as in training but extend the value range to $[-2c,2c]$, where $c>1$ is the OOD scale factor. Additionally, during the test sampling process, we apply a rejection step to ensure that either $\gamma_l < -2$ or $\gamma_u > 2$, while maintaining $-2c \le \gamma_l \le \gamma_u \le 2c$. This ensures that a test sample will not lie in the training domain with high probability.\footnote{This sampling strategy does not guarantee that every test sample falls outside the training distribution. However, as shown in \Cref{app:ood_bound}, the probability of a test sample being within the training distribution is at most $O(1/nc^2)$. In practice, most test instances lie outside this domain. For further details, see \Cref{app:ood_bound}.} This implies that our test results reflect the ``true" out-of-distribution performance of both architectures.

The results in \Cref{fig:fix_scale} show that positional Transformers achieve significantly more stable performance across increasing out-of-distribution ranges than standard Transformers across computational tasks. 
We \textit{hypothesize} that such differences in performance can be attributed to algorithmic alignment \cite{Xu2020What}, a concept that suggests a particular architecture demonstrates better generalization capabilities if its underlying computational structure aligns more closely with the target task. Specifically, the structure of these tasks aligns well with the core motivation for positional attention, i.e., enabling communication through positional information alone. For example, the minimum task is known to have a solution where the underlying computational graph is a binary tree and independent of the input values, see \Cref{app:illustration_algo}. In \Cref{app:failure}, we discuss potential reasons why standard Transformers fail to match the performance of positional Transformers and present additional experiments that illustrate the differing behaviors of self-attention and positional attention.

%

For the $k$-hop induction heads task, OOD data is generated by creating sequences sampled exclusively from a portion of the alphabet not used during training. \Cref{fig:indheads_ood} illustrates the OOD performance of both standard and positional Transformers. Note that the positional Transformer  performs poorly. We hypothesize that this is because the $k$-hop induction heads task deviates from our core motivation, as its inherent task requirement suggests a need for dynamic communication that changes with every input, see also our comments in the description of the tasks above.
%


\subsection{Additional experiments with mixed-type inputs}
\label{sec:nlp}


\begin{comment}
\begin{figure*}[ht]
    \centering
    \begin{subfigure}{0.9\linewidth}
        \includegraphics[width=\linewidth]{plots/nlp/alltasks_mse_noise_p_and_s.pdf}
    \end{subfigure}
    \begin{subfigure}{0.9\linewidth}
        \includegraphics[width=\linewidth]{plots/nlp/alltasks_mape_noise_p_and_s.pdf}
    \end{subfigure}
    \caption{This experiment shows ID/OOD losses (measured as mean squared error, MSE, and mean absolute percentage error, MAPE) for standard Transformers (red) and positional Transformers (blue) across all three variations of the mixed-type input task. The x-axis represents the OOD scale factor. The solid line and shaded area denote the median and the region between the 10\textsuperscript{th} and 90\textsuperscript{th} percentiles over ten trials, respectively.}
    \label{fig:nlp_noise_p_and_t}
\end{figure*}
\end{comment}

In this section, we further investigate the OOD performance when the input data is a mixture of tokens and numbers. This task involves conditional reasoning and simple pattern matching on top of computing an aggregation function. We begin by illustrating a training and test sample, followed by high-level overview of the task. For the definition of the task and the details about the setting, see \Cref{app:nlp}. 

Training sample: {\small
\texttt{Input = ['Cat2', $3.45$, \textcolor{blue}{'Cat+7'}, 'Cat5', $1.23$, 'Cat7', $0.65$, 'Cat8', $2.23$, 'Cat11', \textcolor{blue}{'Cat-8'}, $4.10$, 'Cat13', $1.10$, 'Cat14', $0.10$, 'Cat20', $2.75$, 'Find min of Cat5, Cat8, Cat11 and Cat20']},\quad \texttt{Output = $1.23$}
}

Test sample: {\small
\texttt{Input = ['Cat23', $7.28$, 'Cat24', $33.5$, 'Cat28', $9.17$, 'Cat30', $55.90$, 'Cat31', \textcolor{blue}{'Cat*24'}, $23.70$, 'Cat33', $12.47$, 'Cat34', $8.45$, \textcolor{blue}{'Cat\_40'}, 'Cat40', $1.50$, 'Find min of Cat28, Cat31, Cat33 and Cat40']},\quad \texttt{Output = $1.50$}
}

The input prompt presents the model with a number of textual categories and associated values and asks for the result of an aggregation operation on a subset of the categories. Noise is injected by adding irrelevant categories (with no associated value) in the prompt (colored blue in the examples). Note that the category identifiers, the range of values, and the type of noise all change between training and test data. Specifically, the train category identifiers are integers between $1$ and $20$, whereas the test category identifiers are integers between $21$ and $40$. The models are trained on category values in $[0,5]$ and tested on category values in $[0, 5c]$, for $c=1,2,\dots,10$. We experiment with the following aggregation operations: minimum (min), sum, and a multi-tasking setting that combines minimum (min) and maximum (max). %
This task involves distinguishing relevant from irrelevant categories and pattern matching across differing train and test distributions. The task also demands basic conditional reasoning since each query involves only a subset of all given categories, and algorithmic computation to derive the correct outputs. The textual input is tokenized and processed through an embedding layer, while the numeric input passes through a linear layer. The results, shown in \Cref{fig:nlp_noise_p_and_t}, indicate that the positional Transformer outperforms the Transformer.\footnote{In \Cref{app:nlp} we compare the performance of positional Transformer against a fine-tuned version of GPT2 \cite{radford2019language}. Positional Transformer has better OOD performance. Note that this task does not demonstrate/suggest that positional Transformer will outperform other Transformer-based architectures (like LLMs) on general-purpose reasoning on text data.}




==================================================
Filename: 08_limitations.tex
==================================================
\section{Limitations and future work}

Our work provides a theoretical and experimental investigation of learning algorithms with positional attention. However, further research is needed in some directions that we believe are important. Theoretically, more precise OOD analyses are necessary to capture finer differences between architectures, as existing OOD generalization bounds are not tight enough (see \Cref{app:ood} for an extended discussion). Empirically, our approach relies on fixed positional encodings, limiting its applicability to longer inputs. Developing positional encodings that support arbitrary input lengths would enable us to explore the length generalization capabilities of positional attention.


\begin{comment}
\begin{enumerate}[leftmargin=1cm]
    \item \textbf{OOD generalization theory:} It is very often the case that existing OOD generalization bounds are not tight (see \Cref{app:ood} for an extended discussion). For specific tasks, there is often a gap between what theory says about the worst-case performance and what one observes empirically. This highlights the need for a more fine-grained analysis that will be able to capture the difference in OOD generalization capabilities among different architectures.
    \item \textbf{Length generalization capability:} Our current proposal uses fixed positional encodings, making it difficult to test it on bigger length inputs. Designing positional encodings that can work with arbitrary input lengths will allow us to explore the length generalization capabilities of positional attention.
    \item \textbf{Complementary tasks:} Testing the positional attention mechanism on complementary tasks, such as graph algorithms, requires special treatment. In particular, graph algorithms require that the model effectively process graph connectivity rather than merely treating it as input for the data matrix. Adopting positional attention in various architectures that use graph attention is an exciting future work.
\end{enumerate}

\textcolor{red}{A final potential research direction is understanding why the standard Transformer fails to OOD-generalize even when trained on (seemingly) simple tasks. In \Cref{app:failure} we discuss some potential sources of issues that might be the underlying causes of poor OOD generalization in standard self-attention. For example, a potential problem seems to be related to the stability of self-attention weights against OOD data. In particular, the weights can be very sensitive to the scale of input values. We observe a dramatic change in attention weights as soon as we give the model the same input but scaled so that the values lie outside the domain of the training data. However, more research is needed before we can conclusively answer this question.}
\end{comment}


==================================================
Filename: 10_acknowledgements.tex
==================================================
\section*{Acknowledgements}
The authors would like to thank Federico Barbero (Google DeepMind) and Csaba Szepesvári (Google DeepMind) for their valuable comments and suggestions on this work.

K.~Fountoulakis would like to acknowledge the support of the Natural Sciences and Engineering Research Council of Canada (NSERC). Cette recherche a \'et\'e financ\'ee par le Conseil de recherches en sciences naturelles et en g\'enie du Canada (CRSNG), [RGPIN-2019-04067, DGECR-2019-00147].

S.~Yang would like to acknowledge the support of the Natural Sciences and Engineering Research Council of
Canada (NSERC).

G.~Giapitzakis would like to acknowledge the support of the Onassis Foundation - Scholarship ID: F ZU 020-1/2024-2025.


==================================================
Filename: 991_mpc.tex
==================================================
\section{Background and definitions
}
\label{app:background}
%
\subsection{Massively Parallel Computation (MPC)}
\label{app:mpc}
The MPC model is a computational model of the MapReduce framework \citep{10.1145/1327452.1327492}, widely used for computations over massive datasets.
It defines a parallel computing model that jointly executes a function across multiple machines, each constrained by limited memory capacity. The MPC model is capable of representing various parallel algorithms \citep{im2023massively} and is more powerful than other established parallel models, such as parallel random-access machine (PRAM) \citep{8555148}.

For completeness, we provide a simplified version of the definition of the MPC protocol by \cite{8555148}, which makes the connection to our proxy computational model more apparent\footnote{The original definition includes assumptions about the relation between the input size, the memory size, and the number of machines which 
we omit to better highlight the contrast between MPC and PCOC.}.

\begin{definition}[MPC protocol, Def. I.1 \citep{8555148}, simplified] Let $s$ be a parameter. There are $p \geq 1$ machines (processors), each with local memory of size $s$. The input is distributed on the local memory of some of the machines. The computation proceeds in rounds. In each round, each machine computes the data in its local memory and sends messages to other machines at the end of the round. The total size of messages sent or received by a machine in a round is bounded by $s$. In the next round, each machine only holds the received messages in its local memory. At the end of the computation, the output is in the memory of some machines. 
\end{definition}

%
%

%

%

%
%
%

\subsection{Parallel Computation with Oracle Communication (PCOC)}
\label{app:pcoc}
In this section, we present the definition of PCOC used in \Cref{sec:expressivity} and \Cref{app:expressivity} to establish the expressivity of positional Transformers.
We emphasize that PCOC is merely a proxy model to establish our expressivity results. It is not intended to serve as a competing parallel computational model, such as MPC.

The PCOC model consists of two steps at each round. The first step is communication, where machines send and receive data from other machines. The communication pattern can change at every round. The second step is computation, where all machines perform some local computation on data stored in their local memory.

\textbf{Oracle communication.} For each length $n$ and input data, we assume the existence of an oracle that provides the destination for each machine and message at each round. The oracle executes a communication pattern agnostic to the data being processed. This contrasts with other parallel models, such as the Massively Parallel Computation (MPC) model \citep{8555148}, where it is assumed that the processed data can influence communication patterns.
At first glance, introducing such an oracle might not seem particularly useful, mainly because it is fixed for each input data, where the data can be real-valued, which implies potentially unaccountably many communication oracles for a particular task. However, its importance lies in the fact that for a variety of parallel algorithms, the communication pattern at each round depends only on the length of the input and is independent of other input data. For example, in algorithms that compute basic statistics such as the sum or minimum or tasks such as sorting a list of numbers, the communication between machines is determined not by their values but by their positions. This means that if the values at each machine change, the communication pattern between machines for a given task and input length remains the same. 
In \Cref{app:illustration_algo}, we illustrate communication patterns for some of these tasks, which are also addressed in the experiments in \Cref{sec:experiments}.

%


%

%

%

\begin{definition}[PCOC model] \label{def:pcoc}
The PCOC model is described by a set of $n$ machines labeled from $1$ to $n$, the number of rounds $R$, an integer $s$ and an oracle $\mathsf{RCV}: [R] \times [n] \to [n] \times (\mathcal{P}([s])\setminus \{\emptyset\})$ (which is fixed for a given $n$ and input data) satisfying the following:
\begin{enumerate}
\item Each machine $i$ has a local memory $\mathtt{MEM}_i \in \mathbb{T}^s$ of size $s$, where $\mathbb{T}$ is some abstract data-type. The contents of the memory are indexed from $1$ to $s$ and we use the notation $\mathtt{MEM}_i[j]$ to refer to the element at position $j$ on the memory of the $i$-th machine. 
\item Each machine performs some local computation on the data it receives and overrides the results to its memory. A single machine can perform different local computations on different rounds and different machines (generally) perform different local computations. 
\item When $(r, i)$, where $r\in [R]$ and $i \in [n]$, is passed to the oracle it returns a subset $M$ of the set $[n] \times (\mathcal{P}([s])\setminus \{\emptyset\})$.
The oracle essentially returns a (possibly empty) set of machines that machine $i$ has to receive some data from in round $r$ along with the exact positions on the memories of those machines to retrieve.
\item The total size of data sent and received by each machine is at most $s$. Size here is measured in terms of the number of ``variables" of data-type $\mathbb{T}$.
\end{enumerate}
The protocol is executed in $R$ rounds. At the start, the input is distributed across the $n$ machines. At the beginning of round $r$, each machine $i$ simultaneously queries the oracle with input $(r, i)$ and receives data from the machines it returns. The machines then simultaneously perform their local computations on the received data.
\end{definition}

\begin{framed}
    %
    
    \textbf{Input:} $\mathtt{Data} = (\mathtt{Data}_1,\dots, \mathtt{Data}_n)$ distributed across the memories of $n$ machines, labeled in $[n]$. An oracle $\mathsf{RCV}_{n,\mathtt{Data}}: [R] \times [n] \to [n] \times (\mathcal{P}([s])\setminus \{\emptyset\})$.

\begin{enumerate}[leftmargin=0.33cm]
    \item[1] \textbf{For} each round $r=1,\dots, R$ \textbf{then} 
    \item[2] \hspace{0.5cm} Each machine $i$ simultaneously queries $\mathsf{RCV}_{n,\mathtt{Data}}$ with $(r, i)$ as input and receives \item[] \hspace{0.5cm} data from the machines and memory positions returned by the oracle.
    \item[3] \hspace{0.5cm} The machines simultaneously perform local computations on the received data \item[] \hspace{0.5cm} and write the results in their local memories.
    %
    %
\end{enumerate}
\end{framed}


\subsection{Illustration of parallel algorithms}
\label{app:illustration_algo}
In this section, we further expand on the discussion of \Cref{sec:expressivity}, stating that communication in several parallel algorithms depends only on the identification of the machines rather than their current values. To illustrate this, we provide some concrete examples.

\begin{figure}[h]
    \centering
    \includegraphics[width=0.8\linewidth]{plots/graphs.pdf}
    \caption{Illustration of the computational graph for algorithms such as (cumulative) sum/minimum (on the left) and sorting (on the right) for $n=4$. Circles indicate the machines, each indexed by the subscript. The superscript indicates each round of the algorithm. At round $0$, machine $i$ holds the $i$-th element of the input. 
    The cumulative version of the sum/minimum algorithm includes all arrows (black and orange), while the non-cumulative version is represented only by the orange arrows.}
    \label{fig:alg_examples}
\end{figure}


These tasks are examples of those presented in \Cref{sec:experiments}. Note that these illustrations do not indicate the computational graphs derived by our architecture, as there are multiple ways to achieve the same goal, and they do not necessarily involve the neat graph structures shown in \Cref{fig:alg_examples}. For a more in-depth analysis of the results obtained by our architecture, we refer the reader to \Cref{appdx:experiments}.

In these computational graphs, we represent each machine by a circle, distinguished by a subscript (from 1 to 4, since $n=4$). Furthermore, we use superscripts to denote the rounds of the algorithm, with superscript 0 representing the initial stage where no computation or communication is performed. Note that no specific values are provided in these examples. This indicates that the correct results can be achieved by following the computational graph for any set of values. In the subsequent rounds, each machine receives information from other machines (including itself) and performs some computation.
For each algorithm in \Cref{fig:alg_examples}, we will briefly describe the computations involved, noting that this is not the main focus of the paper and serves only as motivating examples.

For the computation of the minimum and the summing function, each machine applies the minimum (or sum) operator to all incoming values from the previous round. By iteratively applying this operator over multiple rounds, machine 4 ultimately obtains the global minimum value (or total sum) across all $n$ entries, while the other machines hold cumulative minima (or sums) up to their positions. For the non-cumulative versions of these algorithms, the local computations are the same as the cumulative versions, and the communication paths are denoted in orange and form a binary tree pattern.

For sorting, the graph on the right of \Cref{fig:alg_examples} represents the odd-even transposition sort algorithm \citep{habermann1972parallel}. This algorithm works by alternating communication between adjacent machines, starting with the odd-indexed pairs (machines with indices $1$ and $2$, $3$ and $4$, etc.), then switching to even-indexed pairs (machines with indices $2$ and $3$, in this example). In stages where two machines communicate, the machine with the lower index picks the minimum value among the two, while the machine with the higher index picks the maximum. The procedure runs for a total of $n-1$ rounds, after which the values are sorted in ascending order.


\subsection{Additional definitions}
\label{app:back_definition}
In this section, we provide important definitions that are used throughout our theoretical results.
For $n\in \mathbb{N} := \{1,\dots\}$, we use 
$[n]$ to refer to the set $\{1,\dots,n\}$. For two nonnegative integers $m, n \in \mathbb{N} \cup \{0\}$ we use 
$m \veebar n$ to refer to the exclusive or of $m$ and $n$. Next, we define the notion of Hamming neighbors:
\begin{definition}[j\textsuperscript{th} Hamming neighbors]
    Two nonnegative integers $x, y\in [n] \cup \{0\}$ are \textit{j\textsuperscript{th}-Hamming neighbors}, denoted $x \sim_j y$, if:
    \begin{equation*}
    x \sim_j y \iff (\text{bin($x$)}_i = \text{bin($y$)}_i \; \forall i \neq j) \text{ and } (\text{bin($x$)}_j \neq \text{bin($y$)}_j),
    \end{equation*}
    where \text{bin($x$)} is the binary representation of $x$ with $\lceil \log_2 n \rceil$ bits ordered from most to least significant bit.
\end{definition}

Finally, we give the definition of a Bene\v{s} network \cite{benevs1965mathematical}:
\begin{definition}[Bene\v{s} network]
    Let $n=2^k$, $k \in \mathbb{N}$. A Bene\v{s} network is a layered directed graph consisting of $L=2k$ layers with each layer having $n$ nodes and $2n$ directed edges between two sequential layers. Nodes in each layer are labeled with integers from $0$ to $n-1$. We identify each node in the network by a tuple $(\ell, i)$ where $\ell\in [2k]$ is the layer, and $i\in \{0,1,\dots,n-1\}$ is the node label within layer $\ell$. Two nodes $(\ell, i)$ and $(\ell', j)$ are connected by an edge \textit{iff}
    
    $$\ell'=\ell+1 \text{ and } \left(i=j \text{ or }i\sim_{k-\lvert k-\ell\rvert} j\right).$$
    \end{definition}

An example of a Bene\v{s} network with $n=2^3$ nodes in each layer is shown in \Cref{fig:benes}.

\begin{figure}[htb!]
    \centering
    \includegraphics[width=0.6\textwidth]{plots/benes.pdf}
    \caption{A Bene\v{s} network with $n=2^3$ nodes in each layer.}
    \label{fig:benes}
\end{figure}
















==================================================
Filename: 992_expressivity.tex
==================================================
\section{Expressivity results}
\label{app:expressivity}

\subsection{Proof of \Cref{thm:expressivity} (Positional Transformers can simulate MPC)}
\label{app:sim_pt_mpc}


In this section, we prove our main result that establishes the simulation of Massively Parallel Computation (MPC) model by Positional Transformers.
Towards this goal, we present two lemmas that bridge positional Transformers and MPC. 
More specifically, we utilize a proxy model that has a static communication network. We call this proxy model Parallel Computation with Oracle Communication (PCOC), established in \Cref{def:pcoc}. \textit{We note that PCOC is only used to establish our expressivity result. It is not meant as a model to compete with established parallel models, such as MPC}. 
We begin by outlining its main features, followed by the two lemmas that lead to our final result.

The PCOC model operates on $n$ machines labeled from $0$ to $n-1$ over $R$ rounds, with each machine having a local memory of size $s$. The model consists of two fundamental steps per round: computation and communication. During computation, machines perform local computations on data stored in their memory.
In the communication step, machines exchange data through an oracle that provides destination information for each machine and message, executing a communication pattern that is agnostic to the processed data and depends only on the input length. Critically, the oracle's communication pattern remains fixed for a given input length and task, meaning that if the values at each machine change, the communication pattern between machines remains the same, a property that is observed in many parallel algorithms such as computing basic statistics or sorting, where communication is determined by positions rather than values. 

This contrasts with other parallel models, such as MPC, where the communication pattern is more flexible and directly influenced by the data being processed. In the MPC model, each machine has a limited memory size of $s$, and the input data is distributed among the machines. The computation proceeds in rounds: during each round, a machine processes the data in its memory, exchanges messages with other machines (up to a total size of $s$), and retains only the messages it received for the next round. After all rounds, the final output is stored in the memory of one or more machines. Unlike PCOC, the communication in MPC is not predetermined by the input length or task alone; instead, it can dynamically depend on the data being processed, provided the total communication per machine per round does not exceed the memory limit $s$.

Even with limited communication capabilities, PCOC can perform the same functions as MPC, incurring only a logarithmic cost in rounds relative to the number of machines, $N$.
The proof of this lemma is presented in \Cref{app:proof_expr}.

\begin{restatable}{lemma}{MPCPCOC}
\label{lem:pcoc_mpc}
Consider an instance of MPC $\mathsf{T}$ with $R$ rounds, $N$ machines with local memory $s$. 
For any instance $\mathsf{T}$, there exists an instance of PCOC $\mathsf{P}$ with $2R\lceil\log N\rceil$ rounds, $N$ machines with local memory $2s$ that simulates $\mathsf{T}$.
\end{restatable}

An analogy to illustrate the distinction between the two approaches can be drawn from a mailing system. In the MPC model, each processor acts as an individual delivering letters directly to their recipients. Communication is flexible and determined by the specific needs of the task. In contrast, PCOC resembles a structured mailing system with predefined routes and sorting stations. Each individual sends their letters to the nearest sorting station, where the letters are organized into bins and passed through successive stages of sorting and delivery. After several rounds, all letters reach their destinations, following a fixed communication pattern regardless of the specific contents.

Having established the PCOC model, we now show that the positional Transformer model in \Cref{eq:arch} can simulate it.
More specifically, our results show that a $R$-round PCOC protocol can be simulated by a positional Transformer with $R$ layers. 

\begin{restatable}{lemma}{PCOC}
\label{lem:pcoc_sim}
Consider a PCOC instance $\mathsf{P}$ with $R$ rounds, $N$ machines with local memory $s$, and data type $\mathbb{T}=\mathbb{R}$. Let $\mathcal{M}$ be a model following the architecture in \Cref{eq:arch} with $n=N+1$ nodes, $R$ layers and $s$ attention heads. Then, for any instance $\mathsf{P}$ with Borel measurable local functions, there exists a configuration of $\mathcal{M}$ that approximates $\mathsf{P}$ to any desired degree of accuracy.
\end{restatable}

The proof of the above lemma can be found in \Cref{app:simulation}, where we show that each round of computation and communication performed by PCOC can be approximated by one layer of positional Transformers. Having established the two necessary lemmas \Cref{lem:pcoc_mpc} and \Cref{lem:pcoc_sim}, we are ready to re-state our main result.

\expressivity*

The proof of this result follows directly from the combination of the two preceding results: MPC can be simulated by PCOC, which can be simulated by positional Transformers. Notably, the structure of MPC enforces that local memory can only grow sublinearly relative to the input size. This constraint prevents trivial simulation configurations in which all data is sent to a single node, which would undermine the purpose of parallelism.

\subsubsection{Discussion of \Cref{rem:expressivity}}

In \citet{pmlr-v235-sanford24a}, the authors present not only a simulation of MPC using standard Transformers but also the reverse, i.e., a simulation of Transformers by MPC. From the latter, we derive the following remark that establishes an equivalence between standard and positional Transformers.

\equivalence*

An important observation from \citet{pmlr-v235-sanford24a} is that the embedding dimension, denoted as $d_V$, also governs other intermediate dimensions, such as the embedding dimension $d_m$ of the matrices $W_K$ and $W_Q$. However, in positional Transformers, these quantities are decoupled. Specifically, the embedding dimension $d_m$ is directly influenced by the positional encodings and scales linearly with the number of nodes.

It is worth noting that the observed quadratic dependency on the number of nodes in positional Transformers is primarily a consequence of the specific simulation strategy employed in \citet{pmlr-v235-sanford24a}. These dependencies should not be taken too literally, as the results do not provide a sharp characterization of the two computational models, an aspect acknowledged by \citet{pmlr-v235-sanford24a}. Moreover, alternative and more efficient simulation strategies may exist.
Additionally, this quadratic cost arises only when simulating the same sequence of operations across architectures. In practice, the model may converge to parameter configurations that do not correspond to interpretable algorithms.

\subsection{Proof of \Cref{lem:pcoc_mpc} (PCOC can simulate MPC)}
\label{app:proof_expr}

The goal of this section is to show that MPC can be efficiently simulated via PCOC (that is, in a data-agnostic way). For completeness, we re-state the lemma below:

\MPCPCOC*
%
%
%
%
%

We break the proof in two steps: first, we show that MPC routing can be simulated through a fixed communication pattern given by the connectivity of a Bene\v{s} network. Subsequently, we show how this routing can be accomplished in a data-agnostic way, consistent with the PCOC definition. 

Consider an MPC instance with $p$ machines, each with memory size $s$. Assume messages are encoded in the following way: machine $i$ holds a list $\texttt{SEND}_i$ containing tuples of the form $(i, \mathsf{b}, \mathrm{pos}, j)$ where $\mathsf{b}\in \mathbb{A}$, $\mathrm{pos} \in \mathbb{N}$ and $j \in [p]$, where $\mathsf{b}$ is the $\mathrm{pos}$-th character in the message machine $i$ sends to machine $j$. Here, 
$\mathbb{A}$ denotes the alphabet used to encode the messages (for example, $\mathbb{A}$ can be $\{0,1\}$ if we are representing messages by their bits). As a concrete example, if machine $\textcolor{blue}{0}$ sends $010110$ to machine $\textcolor{purple}{1}$ then 
$$\texttt{SEND}_{\textcolor{blue}{0}} = \{(\textcolor{blue}{0}, 0, 1, \textcolor{purple}{1}), (\textcolor{blue}{0}, 1, 2, \textcolor{purple}{1}), (\textcolor{blue}{0}, 0, 3, \textcolor{purple}{1}), (\textcolor{blue}{0}, 1, 4, \textcolor{purple}{1}), (\textcolor{blue}{0}, 1, 5, \textcolor{purple}{1}), (\textcolor{blue}{0}, 0, 6, \textcolor{purple}{1})\}$$
By the assumptions of the MPC model, $|\texttt{SEND}_i| \leq s$ for all $i \in \{0,1,\dots,p-1\}$. Similarly, denote by $\texttt{RCVD}_i$ the list of source-character-position-destination tuples that machine $i$ receives. That is:
$$\texttt{RCVD}_i = \left\{(k, \mathsf{b}, \mathrm{pos}, j) \in \bigcup\nolimits_{i'} \texttt{SEND}_{i'} \,\big|\, j = i \right\} \quad \text{for all } i \in \{0,1,\dots,p-1\}$$
Again, by the assumptions of the MPC model, we have $|\texttt{RCVD}_i|\leq s$ for all $i\in [p]$.
We further impose the following simplifying assumptions:
\begin{enumerate}
    \item The number of machines $p$ is a power of two, i.e. $p=2^k$ for some $k\in \mathbb{N}$, and
    \item For each pair of machines $i$, $j$, the number of tuples sent from machine $i$ to machine $j$ is either $0$ or at least $2^{k-1}$, where $k$ is as above.
\end{enumerate}
Note that the above assumptions do not limit the MPC model since we can always add a couple dummy machines to reach a power of two and pad shorter messages with a special token (potentially having to scale $s$). 
The second assumption can be lifted altogether at the expense of a more complicated proof. We chose to work with both assumptions in order to highlight the essence of the proposed routing algorithm.
As we mentioned previously, to highlight the routing strategy we split the analysis into two parts: first we show how we can simulate one round of MPC using a Bene\v{s} network routing without the need for extra memory, assuming access to destination information. 
Building on that, we then show how to extend the Bene\v{s} routing protocol to the case where message routing is independent of destinations, at the expense of doubling the memory of each machine. The latter shows that PCOC can simulate one round of MPC with the same number of machines and double the memory.
The result is then trivially extended to $R$ rounds of MPC by iterating the routing algorithm $R$ times.

\paragraph{Data-dependent Bene\v{s} routing}
We begin by providing a data-dependent routing algorithm that simulates one round of MPC using a Bene\v{s} network.
Under the two assumptions made above, the routing algorithm executes in $2\log_2(p) = 2k$ steps and is as follows:

\begin{itemize}
    \item Initially, the memory $M_i$ of each machine holds the list $\texttt{SEND}_i$ (i.e. the messages it needs to send in the source-character-position-destination format described above).
    \item At step $L \in \{1, \dots, k-1\}$, machine $i$ partitions the contents of its memory as 
    $$\bigcup_{j=1}^p M_{i,j}$$ where
    $M_{i,j} = \left\{(k,\mathsf{b}, \mathrm{pos}, j') \in M_i \mid j'=j \right\}$. Let $(L+1, m)$ denote the neighbor of $(L, i)$ at level $L+1$ with a label different from $i$ (so $m\neq i$). Machine $i$ partitions each $M_{i,j}$ into two equal parts (this is possible since every message is a power of two) and sends one part to machine $i$ and the other part to machine $m$. Steps $0$ through $k-1$ are collectively referred to as the ``fragmentation phase" of the algorithm. 
    \item At steps $L \in \{k, \dots, 2k-1\}$, denote by $\mathtt{RCVD}_i^{(L)}$ the source-character-position-destination tuples received by machine $i$ after $L-1$ steps. Also, let $(L+1, m)$ denote the neighbor of $(L, i)$ at level $L+1$ with a label different from $i$. At step $L$, machine $i$ sends all tuples $(k, \mathsf{b}, \mathrm{pos}, j) \in \mathtt{RCVD}_i^{(L)}$ such that there is a path between nodes $(L+1, i)$ and $(2k, j)$ in the Bene\v{s} network to machine $i$, and the rest to machine $m$. Notice that by the construction of the Bene\v{s} network, exactly one of following holds:
    \begin{itemize}
        \item There is a path between nodes $(L+1, i)$ and $(2k, j)$, or
        \item There is a path between nodes $(L+1, m)$ and $(2k, j)$
    \end{itemize}
    Steps $k$ through $2k-1$ are collectively referred to as the ``collection phase" of the algorithm. 
    \item At step $2k$, the nodes in the last layer recombine the received messages using the source and position information and the algorithm terminates.
\end{itemize}

We show correctness of our algorithm using a two-step argument. First, we show that during the fragmentation phase (first $k-1$ steps of the algorithm) the tuples are evenly distributed ensuring that the at every point no machine receives more than $s$ tuples. Then we proceed by showing that messages are correctly routed and reach their destination.

\begin{fact}
\label{fact:1}
    After $L \in \{0,1,\dots,k-1\}$ steps, machine $i$ receives exactly $1/2^L$ of each $\texttt{SEND}_{i'}$ where $i'$ is such that $i$ and $i'$ agree on the first $k-L$ least significant bits. Furthermore, no two machines receive the same part of the same list.
\end{fact}
\begin{proof}
    The proof is by induction. For $L=0$ the result trivially holds by the initialization step of the algorithm. Assume it holds for some $L < k-1$. Let $M_i$ be the memory of machine $i$ after $L$ steps. By the connectivity of the Bene\v{s} network, after $L+1$ steps, each machine $i$ will have received messages from itself and machine $m := i \veebar 2^{k-L-1}$. From the inductive hypothesis, $M_i$ contains  $1/2^L$ of each $\texttt{SEND}_{i'}$ where $i'$ is such that $i$ and $i'$ agree on the first $k-L$ least significant bits, and similarly, $M_m$ contains  $1/2^L$ of each $\texttt{SEND}_{i'}$ where $i'$ is such that $m$ and $i'$ agree on the first $k-L$ least significant bits. Now $i$ and $i'$ agree on the first $k-L-1$ least significant bits \text{if and only if} they agree on the first $k-L$ least significant bits or $m$ and $i'$ agree on the first $k-L$ least significant bits. By construction, machine $i$ will receive half the contents of $M_i$ and $M_m$ and it will hence hold exactly $1/2^{L+1}$ of each $\texttt{SEND}_{i'}$ where $i'$ is such that $i$ and $i'$ agree on the first $k-(L+1)$ least significant bits. Since there is at most one path between any two nodes in the Bene\v{s} network, the second part of the statement is trivially satisfied. This concludes the proof.
\end{proof}

\Cref{fact:1} immediately shows that for the first $k-1$ steps of the algorithm, the contents on the memory of each machine never exceed $s$ (since the length of each $\texttt{SEND}_i$ is bounded by $s$). Similarly, we can show the following fact:

\begin{fact}
\label{fact:2}
    After $L\in \{k,k+1,\dots, 2k-1\}$ steps of the algorithm, machine $i$ receives exactly $1/2^{2k-L-1}$ of all $\texttt{RCVD}_{j}$ where $j$ is such that $i'$ and $i$ agree on the first $L-k+1$ least significant bits. Furthermore, no two machines receive the same part of the same list.
\end{fact}
\begin{proof}
We shall only prove the base case $L=k$ as the inductive step is completely symmetric to the proof of \Cref{fact:1}. After the $k$-th step of the algorithm, machine $i$ will receive all tuples for which the destination $j$ is
such that $(2k, j)$ is reachable from $(L + 1, i)$ in the Bene\v{s} network from itself and from machine $m:=i\veebar 2^0$. Recall from \Cref{fact:1} that before the $k$-th step is executed, the memory of machine $i$, $M_i$, contains $1/2^{k-1}$ of every  $\texttt{SEND}_{i'}$ where $i'$ is such that $i$ and $i'$ agree on the least significant bit. Similarly the memory of machine $m$, $M_m$, contains $1/2^{k-1}$ of every $\texttt{SEND}_{i'}$ where $i'$ is such that $m$ and $i'$ agree on the least significant bit. But since the least significant bit of $m$ and $i$ differ, the memories of machines $m$ and $i$ combined, $M_i \cup M_m$, contain exactly $1/2^{k-1}$ of $\bigcup_{i'} \texttt{SEND}_{i'}$ (i.e. a $1/2^{k-1}$ fraction of all messages). Combining the above we see that after the $k$-th step is executed, machine $i$ will receive $1/2^{k-1} = 1/2^{2k-L-1}$ of every  $\texttt{RCVD}_j$ where $j$ is such that $i$ and $j$ agree on the first $L-k+1=1$ least significant bits. As in \Cref{fact:1}, the uniqueness of paths in the Bene\v{s} network guarantees that no two machines ever receive the same part of the same list.
\end{proof}

A direct application of \Cref{fact:2} shows that after $2k-1$ steps of the algorithm, machine each machine $i$ will have in memory exactly $\texttt{RCVD}_i$, showing correctness. Finally, the fact that $|\texttt{RCVD}_i| \leq s$ for all $i \in [p]$ shows that after every step $L \in \{k,k+1,\dots, 2k-1\}$ of the algorithm each machine receives no more than $s$ tuples.
This construction shows that even if we restrict the communication of MPC to some fixed routing pattern, we can still achieve any and all possible routing patterns by using the Bene\v{s} network efficiently, paying the cost of logarithmically more rounds. 
We now proceed to show how we can further extend this algorithm to be consistent with our PCOC model.

\paragraph{From data-dependent to data-agnostic Bene\v{s} routing}
The algorithm above brings us one step closer to our ultimate goal of data-agnostic simulation since at every step, each machine will send data to at most two other machines, respecting the connectivity of the Bene\v{s} network. However, both the fragmentation and collection phases of the algorithm (steps $0$ to $k-1$ and $k$ to $2k-1$) implicitly require access to destination information for each message. 
In terms of PCOC, assuming the oracle mimics the algorithm execution, it would require access to destination information, which is not allowed in the PCOC model. This is because of the following reasons:
\begin{enumerate}
    \item During the fragmentation phase it needs to know exactly which messages are destined for each machine $j$ to evenly divide them into two parts and route them to the two receiving machines.
    \item During the collection phase, destination information is directly used to route messages.
\end{enumerate}
While, at first, these obstacles might seem hard to circumvent, we can do so by using a few simple techniques that enforce a specific structure on the output of the local computation:
\begin{enumerate}
    \item Observe that we can eliminate the need for destination information during the fragmentation phase if we assume that the tuples in the memory of each machine are always sorted based on their destination. If this is the case each machine can simply send the tuples in the even and odd indices on its memory to the two receiving machines, respectively. This achieves the desired outcome: for each $j$, exactly half of the messages with destination $j$ will get routed to one machine and the rest to the other.
    \item As previously discussed, during the collection phase, the oracle would need access to the destination to decide where to route each tuple. However, if we assume that the memory size of each machine is $2s$ (instead of $s$) we can still overcome this issue. Suppose the  oracle needs to simulate step $L \in \{k,k+1,2k-1\}$ of the algorithm. For each machine $i$, let $i_m \neq i$ be such that $(L, i)$ and $(L+1,i_m)$ are connected in the Bene\v{s} network. Recall that at step $k$, machine $i$ will send data to itself and machine $i_m$, depending on the final destination. Given the fixed structure of the Bene\v{s} network, we can require the local computation to calculate where each tuple needs to be routed ($i$ or $i_m$) and place it either in the first half of the memory slots (if it needs to be routed to $i$) or to the second half (if it needs to be routed to $i_m$). The oracle will then route the first half of the memory to $i$ and the second half to $i_m$. Empty spaces on each half can be padded with special tokens.
\end{enumerate}

\subsection{Hardmax patterns using positional attention}
\label{appdx:hardmax}
In this section, we show that the positional attention architecture in \Cref{eq:attn} can approximate any unique hardmax pattern, a concept we define later in this section. This result is used to support the simulation results of \Cref{lem:pcoc_sim}. We begin by stating the definition of the row-wise hardmax transformation for a $p\times q$ matrix $X$ from \Cref{sec:prelim}:
\begin{equation}
    \hardmax(X)_{i,j} = \begin{cases}
    1 & \text{if } X_{i,j} = \max_{k\in[q]} X_{i, k}\\
    0 & \text{otherwise}
\end{cases}
\quad \text{for } i\in [p], j\in [q],
\end{equation}

where we implicitly extend the definition for vectors in $\RR^n$ by viewing them as $1\times n$ matrices.

We use the term \textit{hardmax pattern} to refer to any matrix in the image of hardmax (i.e. a binary matrix with at least one non-zero element in every row). Furthermore, we use the term \textit{unique hardmax pattern} to refer to hardmax patterns with exactly one non-zero element in every row. Unique hardmax patterns occur when the input matrix has a unique maximum value in every row.

We further define key concepts that will be used for the more formal re-statement of \Cref{lem:hardmax}. Let the input have $n$ rows, and the binary positional encodings be defined by the positional encoding matrix $P=I_n$, therefore having $d_P=n$.
Finally, let $T$ be a positive scalar that represents a temperature parameter that controls the approximation of softmax.

\begin{lemma}
\label{lem:hardmax}
For any given $n\times n$  unique hardmax pattern $\bar{A}$, there exists a configuration of node positional attention parameters in \Cref{eq:attn} and a temperature parameter $T$ such that the resulting softmax pattern $A$ approximates $\bar{A}$ to any desired level of accuracy. Formally, for any unique hardmax pattern $\bar{A}$ and any $\varepsilon > 0$, there exists some $T=T(\varepsilon) > 0$ such that the inequality $\lvert\bar{A}_{i,j} - A_{i,j}\rvert \le \varepsilon$ holds for all $i, j \in [n]$.
\end{lemma}

\begin{proof}
Without loss of generality, we may assume $\varepsilon < 1$. We start by setting the node positional attention parameters to be $W_K = I_n$ and $W_Q = T(2\bar{A}-1)$, where, $T > 0$ is our temperature scalar parameter and $\bar{A}$ is the target pattern.
Since, in this construction, node positional encodings are set to be the identity, the inner-product $PW_QW_K^\top P^\top$ reduces to $W_Q$, where each entry $(i,j)$ is $T$ if $\bar{A}_{i,j}=1$, and $-T$ otherwise.

This inner product is passed to the softmax operator, resulting in the attention matrix $A$. For each $i,j \in [n]$ we separately analyze the following two cases:
\begin{enumerate}[leftmargin=0.8cm]
    \item Case $\bar{A}_{i,j}=1$: In that case, the only non-zero element on the $i$-th row of $\bar{A}$ is $\bar{A}_{i,j}$, so we can express the difference as
    \begin{align*}
    \bar{A}_{i,j} - A_{i,j} &= 1 - \frac{\exp((W_Q)_{i,i})}{\sum_{k=1}^n\exp((W_Q)_{i, k})}
    = 1 - \frac{\exp(T)}{\exp(T)+\sum_{k\neq j}\exp(-T)}\\
    &\le 1- \frac{\exp(T)}{\exp(T)+n\exp(-T)}
    = 1-\frac{1}{1+\exp(\ln n-2T)}\\
    &= \frac{\exp(\ln n-2T)}{1+\exp(\ln n-2T)}
    \le \exp(\ln n-2T)\\
\end{align*}

\item Case $\bar{A}_{i,j} = 0$: Let $j_0 \neq j$ be the unique index for which $\bar{A}_{i,j_0}=1$, and we can express the difference as:
\begin{align*}
    A_{i,j} - \bar{A}_{i,j} &= \frac{\exp((W_Q)_{i,j})}{\sum_{k=1}^n\exp((W_Q)_{i, k})}
    = \frac{\exp(-T)}{\exp(T)+\sum_{k\neq j_0}\exp(-T)}\\
    &= \frac{1}{\exp(2T)+n-1}
    \le \frac{1}{\exp(2T-\ln n)+1}
    \le \exp(\ln n -2T)%
\end{align*}
\end{enumerate}

In any case, we have that $\lvert\bar{A}_{i,j}-A_{i,j}\lvert \le \exp(\ln n - 2T)$. Therefore, by taking $T \ge \nicefrac{1}{2}\ln(\nicefrac{n}{\varepsilon})$, we have $\lvert\bar{A}_{i,j}-A_{i,j}\rvert \le \varepsilon$.
\end{proof}


\subsection{Proof of \Cref{lem:pcoc_sim} 
 (Positional Transformers simulate PCOC)}
\label{app:simulation}
We begin this section by outlining the key concepts utilized in the routing protocol employed in our constructions. First, we describe the general structure of the input matrix $\mathbf{X}$.

\subsubsection{Encoding}
\textbf{Input matrix:} In alignment with the PCOC model, the input matrix $\mathbf{X}$ represents $N$ machines, where each machine is denoted by a row in the matrix, and its local memory, $\mathtt{MEM}_i\in\mathbb{T}^s$, is represented by the corresponding columns. The maximum size of data that any machine can send or receive is $s$ bits, with each bit corresponding to a column in $\mathbf{X}$.

However, the actual number of rows and columns in $\mathbf{X}$ differs from the number of machines and the local memory size for two reasons:
\begin{enumerate}
    \item \textbf{Sink node}: A dummy node is introduced to facilitate all possible communication patterns in PCOC using positional attention. This is necessary because PCOC allows for the possibility of information not being sent to any receiving machine. This scenario is incompatible with the softmax formulation, which requires at least one non-zero entry. The dummy node serves as a sink, collecting all messages that do not have a destination. Consequently, the number of rows in $\mathbf{X}$ is $n = N + 1$.
    \item \textbf{Unique node identifier:} Each machine also requires a unique identifier to enable element-wise local computations. To achieve this, we encode a unique scalar for each node in the last column of $\mathbf{X}$, resulting in a feature dimension of $d_X = s + 1$ for the input matrix.
\end{enumerate}

As discussed in \Cref{sec:expressivity}, in PCOC, routing is set by an oracle that decides how packets of data should be routed at each round. Under this framework, routing must be performed to prevent multiple data from being sent to the same destination. 
Since our construction relies on matrix operations, this leads to the following assumption:

\begin{assumption}
 (No-collision). For any layer $\ell \in [L]$, no two different machines $i_1, i_2 \in [N]$ should route data to the same destination $i_3 \in [N]$ for the same column $j \in [s]$, where each column represents a bit of local memory across the nodes.
\end{assumption}

Note that this assumption does not limit the generality of our PCOC model. It only defines how data should be stored in the memory of each receiving machine, and any valid PCOC routing has a corresponding no-collision configuration of bits that realizes it due to the restriction on the total size of received data. As demonstrated in the constructive proof, this directly influences the sparsity pattern generated by each attention head.

\textbf{Positional encodings:} As previously mentioned, although the connectivity at each layer may vary, the positional encodings remain consistent across all layers. Our architecture simulates MPC using node positional encodings with dimension $d_P = n$ by setting $P = I_n$, with each positional encoding functioning as a standard basis vector.

\subsubsection{Simulation results}

We now demonstrate that, with the established encoding, the architecture provided in \Cref{sec:arch} can simulate the execution of any PCOC instance.
Each round of such a PCOC instance can be decomposed into two stages: communication and computation. Our objective is to provide existential results for both stages.

In the communication stage, routing assigns destination machines for each message. In our architecture, this assignment is analogously captured by the attention weights, which determine whether a message should be received by a node using binary values. 

The no-collision assumption ensures that all routing patterns can be represented by unique hardmax patterns.
As expressed in \Cref{lem:hardmax}, since any unique hardmax pattern can be approximated by our attention layer using softmax, for simplicity, the subsequent proofs use hardmax instead of softmax.
With all details now established, we re-state our main simulation result:

\PCOC*
%
%
%
%

\begin{proof} 
Despite the desired degree of accuracy being influenced by the number of rounds performed, it suffices to show that one layer of our architecture can simulate one round of PCOC. The same constructive arguments can be extended to more rounds, ensuring the overall degree of approximation is respected\footnote{
To completely align the simulation results with PCOC execution, positional Transformers require an extra initial layer where the attention layer acts as the identity. This is to account for the fact that PCOC first starts with a computation stage rather than communication as in Transformers.}. To this end, we begin the proof with the communication stage.

\textbf{Communication:} In PCOC, communication is encoded as routing patterns determined by the oracle. At round $\ell \in [R]$, we denote by $H^{(\ell)} = \{((i, j), K) \mid i, j \in [N],\, K \in \mathcal{P}([s])\}$ the set of valid routing patterns provided by the oracle. This set specifies that the data at positions $K$ in the local memory of machine $i$ must be sent to machine $j$ at the same position. A valid routing pattern requires that no collisions occur (i.e., no two triplets in $H^{(\ell)}$ should have the same destination $j$ and memory position $k\in K$). We further denote by $H^{(\ell)}_z = \{((i, j), z) \mid ((i, j), K) \in H^{(\ell)},\, z\in K\}$ the subset of routing patterns corresponding to position $z$ in local memory.

The first part of our result constructively demonstrates that positional attention can reproduce any valid routing set by the oracle that adheres to the PCOC model.
We construct $s$ attention heads indexed by $h \in [s]$, which handle routing for the corresponding subset $H_h$.

For clarity in the construction phase, we introduce an augmented set to simplify notation.
We begin by extracting the set of source nodes for each set  $H^{(\ell)}_z$, denoted as $I^{(\ell)}_z = \{i \mid ((i, j), z) \in H^{(\ell)}_z,\, j\in [N],\, z\in [s]\}$.
Next, we create a complement set $\cancel{H}^{(\ell)}_z$, which routes all unused sources (i.e., those not in $I^{(\ell)}_z$) to the sink node labeled $n=N+1$. We denote this complement set by $\cancel{H}^{(\ell)}_z = \{(i, n, z) \mid i \in [n]\setminus I^{(\ell)}_z\}$.
Finally, we define the union of these sets as $\hat{H}^{(\ell)}_z = H^{(\ell)}_z \cup \cancel{H}^{(\ell)}_z$.

The attention parameters are then set as follows:
\begin{equation*}
    \left(W_K^{(\ell, h)}\right)_{i, j} = \begin{cases}
        1 &\text{if } i=j,\\
        0 &\text{otherwise,}
    \end{cases} \quad \text{ and }\quad
    \left(W_Q^{(\ell,h)}\right)_{i, j} = \begin{cases}
        1 &\text{if } (j, i, h)\in \hat{H}^{(\ell)}_h\\
        0 &\text{otherwise,}
    \end{cases}
\end{equation*}

In this construction, we first observe that both the node positional encodings and the key matrix are identity matrices, reducing the inner product in attention to be solely defined by the query matrix. The query matrix is then designed to encode the source node $i$ as a standard basis vector in the row corresponding to the destination node $j$.
This effectively represents the routing set $\hat{H}^{(\ell)}_h$ as a binary matrix, which is also preserved after applying hardmax.
Additionally, the no-collision assumption, combined with the sink node strategy, ensures exactly one non-zero entry in the first $N$ rows of the attention weights matrix for each attention head $h$.

For the value and output transformation, we set all value matrices $W_V^{(\ell, h)}$ to be the identity $I_{d_X}$ and define the output matrix $W_O^{(\ell)}\in \RR^{(H\cdot d_X) \times (d_X-1)}$ as follows:
\begin{equation*}
    \left(W_O^{(\ell)}\right)_{i, j} = \begin{cases}
        1 &\text{if } i=k+(h-1)s,\,j=h\\
        0 &\text{otherwise.}
    \end{cases}
\end{equation*}

Here, the output matrix $W_O^{(\ell)}$ ensures that only the correct memory position receives the information and places it in the corresponding column. Note that since the outputs of the attention heads are concatenated before being processed by $W_O^{(\ell)}$, 
the values along the rows of the output matrix also depend on the attention head.

We now focus on the computation stage for the second part of the proof.

\textbf{Computation:} At round $\ell \in [R]$ of a PCOC model, let $[\phi^{(\ell, z)}_i]^{n}_{i=1}$ be the local functions applied by each machine $i\in [N]$ and let $\phi^{(\ell, z)}_{n}$ correspond to the function of the augmented sink node, which effectively erases all data received.
Each function $\phi^{(\ell, z)}_i$ operates on received data in each machine's local memory, which corresponds to the output of the attention layer, denoted by $z$ and outputs a vector of the same dimension $s$, that is, $\phi^{(\ell, z)}_i: \RR^s \to \RR^s$.

Furthermore, let $\phi^{(\ell, x)}: \RR^{d_X} \to \RR^{d_X}$ be a function common to all nodes, which operates solely on the residual connection $x$.
This function outputs a vector where all entries are zero except the last entry. The value in this last entry corresponds to the unique node identifier extracted from the residual input $x$.

We aim to approximate both $\phi_i^{(\ell, z)}$ and $\phi^{(\ell, x)}$ using neural networks. To this end, we define the combined function $\phi_i^{(\ell)}: \RR^{d_X+s}\to \RR^{d_X}$ by:

\begin{equation}
    \phi^{(\ell)}_i(z_i \oplus x_i) := \phi^{(\ell, x)}(x_i) + \phi_i^{(\ell, z)}(z_i)\oplus0,
\end{equation}

where $z_i \oplus x_i$ denotes the concatenation of output of $z_i \in \mathbb{R}^s$ and $x_i \in \mathbb{R}^{d_X}$. We further augment the output of $\phi_i^{(\ell, z)}$ with a zero scalar to match the dimension $d_X=s+1$.

Let $[\hat{\phi}^{(\ell)}_i]^{n}_{i=1}$ correspond to multilayer perceptrons (MLPs), each applied to each input $z_i \oplus x_i$.
By invoking universal approximation results such as those by \cite{cybenko1989approximation, hornik1989multilayer}, we assert that as long as the local functions $\phi^{(\ell)}_i$ are Borel measurable, there exist neural networks $\hat{\phi}^{(\ell)}_i$ that can approximate the functions $\phi^{(\ell)}_i$ to any desired degree of accuracy.
Additionally, note that the function $\phi^{(\ell, z)}_n$ of the sink node, as well as the function $\phi^{(\ell, x)}$ that operates on the residual connection, are both linear and therefore Borel measurable.

The final step in this argument is to relate these approximations to the proposed architecture in \Cref{sec:arch}. 
Specifically, we use the MLP $\Phi^{(\ell)}$ in \Cref{eq:arch}
 and leverage the aforementioned universality results to approximate all the element-wise functions $\phi^{(\ell)}_i$.

A crucial aspect of this step is the need for the input of each machine to be uniquely identifiable. This ensures that a single model can injectively encode multiple functions. Intuitively, it guarantees that each approximation of the local function can identify that it is processing the right row.
The unique identification of each machine is guaranteed by the scalar encodings of every node, which, regardless of the contents in local memory, ensure that the input rows are unique. Therefore, the function that $\Phi^{(\ell)}$ has to approximate is a piecewise Borel function with each branch being one of the $\phi_i^{(\ell)}$, based on the unique machine identifier. Such function is Borel measurable, and so the universal approximation results of \cite{hornik1989multilayer} hold, guaranteeing the existence of the desired MLP $\Phi^{(\ell)}$.

This demonstrates that our neural network architecture can emulate the computations performed by the local functions $\phi_i^{(\ell, z)}$ acting on the output of the attention layer (with their outputs zero-padded to match the required dimension) and the function $\phi^{(\ell, x)}$ acting on the residual connection, even though they act on distinct parts of the input.

Therefore, we establish that our proposed architecture can approximate the computations in each round of the PCOC model.
%
\end{proof}

An important observation is that the computational model and expressive results for the proposed architecture are specific to a fixed input length $n$.
Furthermore, one could extend such results to a model with communication and local computations that also consider the input length as an input.
For local computations, proof in \Cref{lem:pcoc_sim} can also cover such cases, provided that the information about the length is also encoded in the input.
For communication, we present the following remark.

\begin{remark}
For any collection of unique hardmax patterns $\{\bar{A}^{(k)}\}_{k=1}^n$, where $\bar{A}^{(k)}$ is $k \times k$, there exists a configuration of node positional attention parameters in \Cref{eq:attn} and a temperature parameter $T$ such that the resulting softmax patterns $\{A^{(k)}\}_{k=1}^n$ approximate each $\bar{A}^{(k)}$ to any desired level of accuracy. Formally, for any collection of unique hardmax patterns $\{\bar{A}^{(k)}\}_{k=1}^n$ and any $\varepsilon > 0$, there exists a temperature parameter $T = T(\varepsilon) > 0$ and corresponding attention parameters such that for all $i, j \in [n]$ and for all $k \in [n]$, the following inequality holds: $\left|\bar{A}^{(k)}_{i,j} - A^{(k)}_{i,j} \right| \leq \varepsilon$.
\end{remark}

The proof of this remark relies on a slight modification of the proof of \Cref{lem:hardmax}.
However, to cover all possible patterns, the embedding dimension of positional encodings should also encode the input length and be of the order of $O(n^3)$.
Although this embedding dimension is theoretically large, in practice, one does not need as many dimensions for positional encodings, as demonstrated in the variable length experiments in \Cref{sec:experiments}.

\subsection{Softmax patterns using positional attention}
\label{app:softmax}
We conclude the discussion on expressivity by showing a final, standalone, result, namely that the positional attention architecture in \Cref{eq:attn} can represent any softmax pattern. This result is further discussed in \Cref{app:failure} regarding reasons for poor out-of-distribution performance by standard Transformers.

We begin by stating the definition of the row-wise softmax transformation for a matrix $X \in \RR^{p\times q}$:

\begin{equation}
    \softmax(X)_{i,j} = \frac{\exp(X_{i,j})}{\sum_{k=1}^q \exp(X_{i,k})} \quad \text{for } i\in [p], j\in [q]
\end{equation}

As with hardmax, the definition is implicitly extended to vectors in $\RR^n$ by viewing them as $1\times n$ matrices. The image of the softmax function is the set of row-stochastic matrices with entries in $(0, 1)$. Indeed, it is easy to see that when softmax is applied to a matrix, the resulting matrix satisfies the above property. On the other hand, for a matrix $\mathbf{B} = (b_{ij})_{i\in[p], j\in [q]}$ with $b_{ij}\in (0,1)$ and $\sum_{j\in [q]} b_{ij} = 1$ for all $i \in [p]$ we have $\softmax(\mathbf{\Tilde{X}}) = B$ where $\mathbf{\Tilde{X}}_{i,j} = \ln(b_{ij})$. We use the term \textit{softmax pattern} to refer to any matrix in the image of softmax.

Consider attention weights $A^{(\ell, h)}$ that are defined by positional encodings in \Cref{eq:attn}. Let $B\in (0, 1)^{n\times n}$ be a softmax pattern. We would like to find parameters $W_Q^{(\ell, h)}$ and $W_K^{(\ell, h)}$ that induce $B$, that is $A^{(\ell, h)} = B$. From the properties of softmax described above, it suffices to solve the matrix equation $(PW^{(\ell, h)}_Q)\cdot(PW^{(\ell, h)}_K)^\top = \Tilde{B}$ where $\Tilde{B}_{ij} = \ln(B_{ij})$. This equation always has a solution when $d_P = n$ and $P$ is invertible. We summarize the above observation in the following expressivity remark:
\begin{remark}[Positional attention is expressive]
    \label{rem:softmax}
    Positional attention can realize all softmax patterns at every layer provided that $d_P = n$ and $P$ is invertible. This is not necessarily true in the case of standard attention where, in subsequent layers, positional encodings are modified and, therefore, not guaranteed to be linearly independent.
\end{remark}



==================================================
Filename: 993_complexity.tex
==================================================
\section{Proof of Generalization Bound (\Cref{thm:generalization})}
\label{app:gen_bound}

This section provides a complete proof of \Cref{thm:generalization}. The high level strategy which inductively builds a cover for the class of multi-layer architectures is similar to that of \citet{edelman2022inductive}. We made a couple of changes to account for positional attention. For the reader's convenience we re-state all relevant definitions below.

\subsection{Preliminaries}
\label{app:gen_prelim}

\subsubsection{Notations and the network architecture}

For a matrix $Z$ we denote by $Z_{i,:}$ and $Z_{:,i}$ the $i$-th row and column of $Z$, respectively. $\|\cdot\|_2$ denotes the spectral norm for matrices, and $\|\cdot\|_{p,q}$ denotes the $(p,q)$ matrix norm where the $p$-norm is over columns and $q$-norm over rows. For vectors, $\|\cdot\|_p$ denotes the $\ell_p$ norm. For example, $\|Z\|_{2,\infty} := \max_{j} \|Z_{:,j}\|_2$ and $\|Z\|_{2,1} = \sum_{j}\|Z_{:,j}\|_2$.


We are interested in the positional transformer architecture of $L$ layers, which is obtained by composing \Cref{eq:arch} $L$ times. In addition, we parameterize $\Phi^{(\ell)}$ to be a 2-layer MLP for each $\ell$. Given input $X \in \mathbb{R}^{n\times d_X}$ and positional encodings $P\in\mathbb{R}^{n\times d_P}$, the resulting architecture has $L$ layers, each layer has $H$ attention heads, and it is written as
\begin{allowdisplaybreaks}
\begin{align}
  &F^{(0)}(X;W^{1:0}) := X,\nonumber\\
  &\forall \ell = 1,2,\ldots,L:\nonumber\\
  &\hspace{5mm}F^{(\ell)}(X;W^{1:\ell}) := \sigma\left(\left(\left(\bigoplus_{h=1}^H A^{(\ell,h)}F^{(\ell-1)}(X;W^{1:\ell-1})W_V^{(\ell,h)}\right)\oplus F^{(\ell-1)}(X;W^{1:\ell-1})\right)W_1^{(\ell)}\right)W_2^{(\ell)},\label{eq:F}\\
  &\hspace{20mm} \mbox{with} ~~ A^{(\ell,h)} := \softmax(PW_Q^{(\ell,h)}{W_K^{(\ell,h)}}^\top P^\top ), \quad \forall h = 1,2,\ldots,H.\nonumber
\end{align}
\end{allowdisplaybreaks}
In the above, $\sigma$ is an $L_\sigma$-Lipschitz activation function with $\sigma(0) = 0$. We write 
\[
  W^{(\ell)} := \{W_Q^{(\ell,1)},W_Q^{(\ell,2)},\ldots,W_Q^{(\ell,H)}, W_K^{(\ell,1)}, W_K^{(\ell,2)},\ldots, W_K^{(\ell,H)}, W_V^{(\ell,1)}, W_V^{(\ell,2)}, \ldots, W_V^{(\ell,H)}, W_1^{(\ell)}, W_2^{(\ell)}\}
\]
for the set of weight matrices at layer $\ell$, and
\[
  W^{1:\ell} := (W^{(1)}, W^{(2)},\ldots,W^{(\ell)})
\]
for the set of weight matrices up to layer $\ell$. Note that we may assume without loss of generality that the output projection matrix $W_O^{(\ell)}$ is absorbed into the weight matrix $W_1^{(\ell)}$ of the MLP, and therefore we omit $W_O^{(\ell)}$ in \Cref{eq:F}. Denote $d=\max\{d_X,d_P,d_V,d_{\mathrm{out}}\}$. Since we are to derive an upper bound which increases with $d_X,d_P,d_V,d_{\mathrm{out}}$, we sometimes replace $d_X,d_P,d_V,d_{\mathrm{out}}$ with their maximum $d$ in the derivations to simplify the notation.

Given the final output $X_{\mathsf{out}} = F^{(L)}(X;W^{1:L}) \in \mathbb{R}^{n \times n}$, we extract a scalar output from the architecture as $w^\top X_{\mathsf{out}}[n]$ with trainable weights $w\in\mathbb{R}^d$. This defines a class of scalar-output architectures as in \citet{edelman2022inductive}. The index $i\in[n]$ on which we apply the linear function $x \mapsto w^\top x$ can be arbitrary, and here we simply set to the last index, $n$. Our generalization bound applies to the following class of scalar-output architectures with bounded-norm parameters:
\begin{equation}\label{eq:bounded_norm_arch_class}
\begin{split}
  \mathcal{F} := \Bigg\{&X \mapsto w^\top \left(F^{(L)}(X;W^{1:L})[n]\right): \|w\|_2 \le B_w,\\
  &\left\|W_V^{(\ell,h)}\right\|_2, \left\|W_1^{(\ell)}\right\|_2, \left\|W_2^{(\ell)}\right\|_2 \le B_2, \forall \ell \in [L], ~ \forall h \in [H],\\
  &\left\|W_K{^{(\ell,h)}}^\top W_Q{^{(\ell,h)}}\right\|_{2,1}, \left\|W_V^{(\ell,h)}\right\|_{2,1}, \left\|W_1^{(\ell)}\right\|_{2,1}, \left\|W_2^{(\ell)}\right\|_{2,1} \le B_{2,1}, ~ \forall \ell \in [L], \forall h \in [H]\Bigg\}.
\end{split}
\end{equation}

\subsubsection{Generalization bound, Rademacher complexity, and the covering number}

\begin{definition*}[Risk]
Let $\mathcal{D}$ be a distribution over $\mathcal{X}\times \mathbb{R}$, and let $\ell : \mathbb{R}\times\mathbb{R} \rightarrow \mathbb{R}$ be the loss function. For a given class of functions $\mathcal{F} = \{ f : \mathcal{X} \rightarrow \mathbb{R} \}$ and $f\in\mathcal{F}$, the generalization error, or the risk, and the empirical risk, are defined as
\[
  \mathsf{risk}(f;\mathcal{D}) := \mathbb{E}_{(x,y)\sim\mathcal{D}} \left[\ell(f(x),y)\right], \quad \widehat{\mathsf{risk}}(f;(x{(i)},y^{(i)})_{i=1}^m) := \frac{1}{m}\sum_{i=1}^m \ell(f(x^{(i)}), y^{(i)}).
\]
\end{definition*}

\begin{definition}[Empirical Rademacher complexity]
For a given class of functions $\mathcal{F} = \{ f : \mathcal{X} \rightarrow \mathbb{R} \}$ and $\{x^{(i)}\}_{i=1}^m \subset \mathcal{X}$, the empirical Rademacher complexity $\widehat{\mathcal R}(\mathcal{F};\{x^{(i)}\}_{i=1}^m)$ is defined as
\[
  \widehat{\mathcal R}(\mathcal{F};\{x^{(i)}\}_{i=1}^m) = \frac{1}{m} \mathbb{E}_\varepsilon \left[\sup_{f \in \mathcal F} \sum_{i=1}^m \varepsilon_i f(x^{(i)})\right],
\]
where $\varepsilon$ is a vector of $m$ i.i.d. Rademacher random variables, that is, $\mathbb{P}(\varepsilon_i = 1) = \mathbb{P}(\varepsilon_i = -1) = 1/2$.
\end{definition}


The following theorem bounds the generalization gap $|\mathsf{risk}-\widehat{\mathsf{risk}}|$ using the Rademacher complexity of a function class.

\begin{theorem}[\citet{bartlett2003rademacher}]\label{thm:generalization_from_rademacher}
Let $\mathcal{D}$ be a distribution over $\mathcal{X} \times \mathbb{R}$ and let $\ell : \mathbb{R} \times \mathbb{R} \rightarrow \mathbb{R}$ be a
$b$-bounded loss function that is $L$-Lipschitz in its first argument. Let $\mathcal{F} = \{f : \mathcal{X} \rightarrow \mathbb{R}\}$ be a given class of functions. Then for any $\delta > 0$, with probability at least $1-\delta$, simutaneously for all $f \in \mathcal{F}$, it holds that
\[
  \left| \mathsf{risk}(f;\mathcal{D}) - \widehat{\mathsf{risk}}(f;(x{(i)},y^{(i)})_{i=1}^m) \right|
  \le 
  4L \widehat{\mathcal R}(\mathcal{F};\{x^{(i)}\}_{i=1}^m) + 2b \sqrt{\frac{\log(1/\delta)}{2m}}.
\]
\end{theorem}

We may bound the Rademacher complexity of a function class $\mathcal{F}$ using its covering number.

\begin{definition}[Covering number]
For a given class of (scalar-valued or vector-valued) functions $\mathcal F$, the covering number $\mathcal{N}_{\infty}(\mathcal{F};\varepsilon;\{x^{(i)}\}_{i=1}^m;\|\cdot\|)$ is the smallest size of a collection (a cover) $\mathcal{C} \subset \mathcal{F}$ such that for all $f\in\mathcal{F}$ there is $\hat{f}\in\mathcal{C}$ satisfying
\[
  \max_{i\in[m]} \left\|f(x^{(i)} - \hat{f}(x^{(i)})\right\| \le \varepsilon.
\]
\end{definition}

\begin{lemma}[Rademacher complexity via covering number; Corollary A.3 in \citet{edelman2022inductive}]\label{lem:rademacher_from_cover}
For a class of functions $\mathcal{F} = \{ f : \mathcal{X} \rightarrow \mathbb{R} \}$
such that $|f| \le A$ for all $f \in \mathcal{F}$, suppose that $\log \mathcal{N}_{\infty}(\mathcal{F};\varepsilon;\{x^{(i)}\}_{i=1}^m) \le C_{\mathcal{F}}/\varepsilon^2$, then for some consstant $c>0$:
\[
  \widehat{\mathcal R}(\mathcal{F};\{x^{(i)}\}_{i=1}^m) \le c \cdot \sqrt{\frac{C_{\mathcal{F}}}{m}} \cdot \left( 1 + \log\left(A\sqrt{m/C_{\mathcal{F}}}\right)\right).
\]
\end{lemma}

Combining \Cref{thm:generalization_from_rademacher} and \Cref{lem:rademacher_from_cover} we may bound the generalization gap $|\mathsf{risk}-\widehat{\mathsf{risk}}|$ using the covering number of a function class, as stated in the following lemma.

\begin{lemma}\label{lem:generalization_from_cover}
For a class of functions $\mathcal{F} = \{ f : \mathcal{X} \rightarrow \mathbb{R} \}$
such that $f$ is bounded for all $f \in \mathcal{F}$ and that $\log \mathcal{N}_{\infty}(\mathcal{F};\varepsilon;\{x^{(i)}\}_{i=1}^m) \le C_{\mathcal{F}}/\varepsilon^2$ for all $\{x^{(i)}\}_{i=1}^m \subset \mathcal{X}$. Then for any $\delta>0$, with probability at least $1-\delta$, simultaneously for all $f\in\mathcal{F}$, it holds that
\[
  \left| \mathsf{risk}(f;\mathcal{D}) - \widehat{\mathsf{risk}}(f;(x{(i)},y^{(i)})_{i=1}^m) \right|
  \le
  \tilde O \left(\sqrt{\frac{C_{\mathcal{F}}}{m}} + \sqrt{\frac{\log(1/\delta)}{m}}\right).
\].
\end{lemma}

\subsubsection{Some useful lemmas}

\begin{lemma}[Corollary A.7 in \citet{edelman2022inductive}]\label{lem:lip_softmax}
For vectors $x,y\in\mathbb{R}^n$ we have that 
$$
\normOne{\softmax(x) - \softmax(y)} \le 2 \norminf{x - y}.
$$
\end{lemma}

\begin{lemma}[Lemma 4.6 in \citet{edelman2022inductive}]\label{lem:basic_cover}
Let $\mathcal{W}:= \{W\in\mathbb{R}^{d \times d'} \ : \ \|W^\top \|_{2,1} \le B_W\}$, and consider the function class $\mathcal{F}_W := \{x \mapsto Wx  :  W\in\mathcal{W}\}$. For any $\varepsilon > 0$ and $x^{(1)},x^{(2)},\ldots,x^{(m)}\in\mathbb{R}^{d}$ satisfying $\|x^{(i)}\|_2 \le B_X$ for all $i\in[m]$, we have
\begin{equation*}
  \log \mathcal{N}_\infty(\mathcal{F}_W;\varepsilon; x^{(1)},\ldots,x^{(m)};\|\cdot\|_2) \lesssim \frac{B_W^2 B_X^2}{\varepsilon^2} \log (d m).
\end{equation*}
\end{lemma}

\begin{lemma}[Lemma A.8 in \citet{edelman2022inductive}]\label{lem:setting_epsilons}
For $a_1, b_i \ge 0$, the solution to the following optimization problem
\begin{align*}
  \min_{\varepsilon_1,\ldots,\varepsilon_t} & \quad \sum_{i=1}^t \frac{a_i}{\varepsilon_i^2} \\
  \mathrm{s.t} & \quad \sum_{i=1}^t b_i \varepsilon_i = \varepsilon
\end{align*}
is $\nu^3/\varepsilon^2$ and is achieved at $\varepsilon_i = \varepsilon( a_1/b_i)^{1/3}/\nu$ where $\nu:=\sum_{i=1}^t a_i^{1/3} b_i^{2/3}$.
\end{lemma}

\subsection{Lipschitz Property of the Architecture}
\label{app:gen_lip}

Denote 
\begin{align*}
  f(Z;\{W_Q^{(\cdot,h)}\}_{h=1}^H,\{W_K^{(\cdot,h)}\}_{h=1}^H,\{W_V^{(\cdot,h)}\}_{h=1}^H) 
  &:= \left(\bigoplus_{h=1}^H A^{(\cdot,h)}ZW_V^{(\cdot,h)}\right)\oplus Z \\
  &= \left[A^{(\cdot,1)}ZW_V^{(\cdot,1)}, A^{(\cdot,2)}ZW_V^{(\cdot,2)}, \cdots, A^{(\cdot,H)}ZW_V^{(\cdot,H)}, Z \right],
\end{align*}
where $A^{(\cdot,h)} = \softmax(PW_Q^{(\cdot,h)}{W_K^{(\cdot,h)}}^\top P^\top )$.

\begin{lemma}\label{lem:lip_pos_attn}
For any $Z$, $\widehat{Z}$, $\{W_Q^{(\cdot,h)}\}_{h=1}^H$, $\{\widehat{W}_Q^{(\cdot,h)}\}_{h=1}^H$, $\{W_K^{(\cdot,h)}\}_{h=1}^H$, $\{\widehat{W}_K^{(\cdot,h)}\}_{h=1}^H$, $\{W_V^{(\cdot,h)}\}_{h=1}^H$, $\{\widehat{W}_V^{(\cdot,h)}\}_{h=1}^H$,
\begin{align*}
  &~\left\|
  \left(f\left((Z;\{W_Q^{(\cdot,h)}\}_{h=1}^H,\{W_K^{(\cdot,h)}\}_{h=1}^H,\{W_V^{(\cdot,h)}\}_{h=1}^H\right) - f\left(\widehat{Z};\{\widehat{W}_Q^{(\cdot,h)}\}_{h=1}^H,\{\widehat{W}_K^{(\cdot,h)}\}_{h=1}^H,\{\widehat{W}_V^{(\cdot,h)}\}_{h=1}^H\right)\right)^\top \right\|_{2,\infty}\\
  &~\le \max\left(\max_{h\in[H]}\|W_V^{(\cdot,h)}\|_2,1\right) \\
  &\hspace{10mm} \cdot \Bigg(\|Z^\top -\widehat{Z}^\top \|_{2,\infty} + 2\|Z^\top \|_{2,\infty}\|P^\top \|_{2,\infty}\max_{h\in[H]}\left\|\left(W_Q^{(\cdot,h)}{W_K^{(\cdot,h)}}^\top -\widehat{W}_Q^{(\cdot,h)}\widehat{W}_K{^{(\cdot,h)}}^\top \right)P^\top \right\|_{2,\infty}\Bigg)\\
  &\qquad + \sqrt{H}\max_{h\in[H]} \left\|\left(W_V^{(\cdot,h)}-\widehat W_V^{(\cdot,h)}\right)^\top \widehat Z^\top \right\|_{2,\infty}
\end{align*}
\end{lemma}

\begin{proof}
Denote $A^{(\cdot,h)} = \softmax(PW_Q^{(\cdot,h)}{W_K^{(\cdot,h)}}^\top P^\top )$ and $\widehat A^{(\cdot,h)} = \softmax(P\widehat W_Q^{(\cdot,h)}\widehat W_K{^{(\cdot,h)}}^\top P^\top )$:
\begin{allowdisplaybreaks}
\begin{align}
  &\left\|
  \left(f\left((Z;\{W_Q^{(\cdot,h)}\}_{h=1}^H,\{W_K^{(\cdot,h)}\}_{h=1}^H,\{W_V^{(\cdot,h)}\}_{h=1}^H\right) - f\left(\widehat{Z};\{\widehat{W}_Q^{(\cdot,h)}\}_{h=1}^H,\{\widehat{W}_K^{(\cdot,h)}\}_{h=1}^H,\{\widehat{W}_V^{(\cdot,h)}\}_{h=1}^H\right)\right)^\top \right\|_{2,\infty}\nonumber\\
  =~&\left\|\begin{bmatrix}W_V{^{(\cdot,1)}}^\top (A^{(\cdot,1)}Z)^\top \\ \vdots \\ W_V{^{(\cdot,H)}}^\top (A^{(\cdot,H)}Z)^\top \\ Z^\top \end{bmatrix} - \begin{bmatrix}\widehat W_V{^{(\cdot,1)}}^\top (\widehat A^{(\cdot,1)}\widehat Z)^\top \\ \vdots \\ \widehat W_V{^{(\cdot,H)}}^\top (\widehat A^{(\cdot,H)}\widehat Z)^\top \\ \widehat Z^\top \end{bmatrix}\right\|_{2,\infty} \nonumber\\
  \le~&\left\|\begin{bmatrix}W_V{^{(\cdot,1)}}^\top \left(A^{(\cdot,1)}Z-\widehat A^{(\cdot,1)}\widehat Z\right)^\top \\ \vdots \\ W_V{^{(\cdot,H)}}^\top \left(A^{(\cdot,H)}Z-\widehat A^{(\cdot,H)}\widehat Z\right)^\top \\ Z^\top -\widehat Z^\top \end{bmatrix}\right\|_{2,\infty} + \quad\left\|\begin{bmatrix}\left(W_V^{(\cdot,1)}-\widehat W_V^{(\cdot,1)}\right)^\top (\widehat A^{(\cdot,1)}\widehat Z)^\top \\ \vdots \\ \left(W_V^{(\cdot,H)}-\widehat W_V^{(\cdot,H)}\right)^\top (\widehat A^{(\cdot,H)}\widehat Z)^\top \\ 0\end{bmatrix}\right\|_{2,\infty}\nonumber\\
  =~& \left\|\begin{bmatrix}W_V{^{(\cdot,1)}}^\top  & & & \\ & \ddots & & \\ & & W_V{^{(\cdot,H)}}^\top  & \\ & & & I \end{bmatrix} \begin{bmatrix}(A^{(\cdot,1)}Z)^\top -\widehat A^{(\cdot,1)}\widehat Z)^\top \\ \vdots \\ (A^{(\cdot,H)}Z)^\top -\widehat A^{(\cdot,H)}\widehat Z)^\top \\ Z^\top -\widehat Z^\top \end{bmatrix}\right\|_{2,\infty} +\quad\left\|\begin{bmatrix}\left(W_V^{(\cdot,1)}-\widehat W_V^{(\cdot,1)}\right)^\top (\widehat A^{(\cdot,1)}\widehat Z)^\top \\ \vdots \\ \left(W_V^{(\cdot,H)}-\widehat W_V^{(\cdot,H)}\right)^\top (\widehat A^{(\cdot,H)}\widehat Z)^\top \end{bmatrix}\right\|_{2,\infty}\nonumber\\
  \le~&\left\|\begin{bmatrix}W_V{^{(\cdot,1)}}^\top  & & & \\ & \ddots & & \\ & & W_V{^{(\cdot,H)}}^\top  & \\ & & & I \end{bmatrix}\right\|_2\left\|\begin{bmatrix}(A^{(\cdot,1)}Z)^\top -\widehat A^{(\cdot,1)}\widehat Z)^\top \\ \vdots \\ (A^{(\cdot,H)}Z)^\top -\widehat A^{(\cdot,H)}\widehat Z)^\top \\ Z^\top -\widehat Z^\top \end{bmatrix}\right\|_{2,\infty} +\quad\left\|\begin{bmatrix}\left(W_V^{(\cdot,1)}-\widehat W_V^{(\cdot,1)}\right)^\top (\widehat A^{(\cdot,1)}\widehat Z)^\top \\ \vdots \\ \left(W_V^{(\cdot,H)}-\widehat W_V^{(\cdot,H)}\right)^\top (\widehat A^{(\cdot,H)}\widehat Z)^\top \end{bmatrix}\right\|_{2,\infty}\nonumber\\
  \le~&\max\left(\max_{h\in[H]}\|W_V^{(\cdot,h)}\|_2,1\right)\cdot \left\|\begin{bmatrix}(A^{(\cdot,1)}Z)^\top -\widehat A^{(\cdot,1)}\widehat Z)^\top \\ \vdots \\ (A^{(\cdot,H)}Z)^\top -\widehat A^{(\cdot,H)}\widehat Z)^\top \\ Z^\top -\widehat Z^\top \end{bmatrix}\right\|_{2,\infty} + \quad \sqrt{H}\max_{h\in[H]} \left\|\left(W_V^{(\cdot,h)}-\widehat W_V^{(\cdot,h)}\right)^\top \left(\widehat A^{(\cdot,h)}\widehat Z\right)^\top \right\|_{2,\infty}\nonumber\\
  \le~&\max\left(\max_{h\in[H]}\|W_V^{(\cdot,h)}\|_2,1\right)\cdot\left(\left\|Z^\top \left[A{^{(\cdot,1)}}^\top -\widehat A{^{(\cdot,1)}}^\top , \ldots, A{^{(\cdot,H)}}^\top  - \widehat A{^{(\cdot,H)}}^\top \right]\right\|_{2,\infty}\right.\nonumber\\
  &\hspace{10mm}+\left.
  \left\|(Z-\widehat Z)^\top \left[\widehat A{^{(\cdot,1)}}^\top ,\ldots,\widehat A{^{(\cdot,H)}}^\top , I \right]\right\|_{2,\infty}\right) + \sqrt{H}\max_{h\in[H]} \left\|\left(W_V^{(\cdot,h)}-\widehat W_V^{(\cdot,h)}\right)^\top \left(\widehat A^{(\cdot,h)}\widehat Z\right)^\top \right\|_{2,\infty}.\label{eq:AZ-A_hatZ_hat}
\end{align}
\end{allowdisplaybreaks}
Using $\|Pv\|_2 \le \|P\|_{2,\infty}\|v\|_1$ (and therefore $\|PQ\|_{2,\infty}\le \|P\|_{2,\infty}\|Q\|_{1,\infty}$):
\begin{equation}\label{eq:(Z-Z_hat)A_hat}
  \left\|(Z-\widehat Z)^\top \left[\widehat A{^{(\cdot,1)}}^\top ,\ldots,\widehat A{^{(\cdot,H)}}^\top , I \right]\right\|_{2,\infty} 
  \le \left\|(Z-\widehat Z)^\top \right\|_{2,\infty} \left\|\left[\widehat A{^{(\cdot,1)}}^\top ,\ldots,\widehat A{^{(\cdot,H)}}^\top , I \right]\right\|_{1,\infty} 
  = \left\|Z^\top -\widehat Z^\top \right\|_{2,\infty},
\end{equation}
and
\begin{allowdisplaybreaks}
\begin{align}
  \left\|\left(W_V^{(\cdot,h)}-\widehat W_V^{(\cdot,h)}\right)^\top \left(\widehat A^{(\cdot,h)}\widehat Z\right)^\top \right\|_{2,\infty}
  &\le \left\|\left(W_V^{(\cdot,h)}-\widehat W_V^{(\cdot,h)}\right)^\top  \widehat Z^\top \right\|_{2,\infty} \left\|\widehat A{^{(\cdot,h)}}^\top \right\|_{1,\infty}\nonumber\\
  &= \left\|\left(W_V^{(\cdot,h)}-\widehat W_V^{(\cdot,h)}\right)^\top  \widehat Z^\top \right\|_{2,\infty}. \label{eq:(W-W_hat)AZ}
\end{align}
\end{allowdisplaybreaks}
Using $\|PQ\|_{2,\infty}\le \|P\|_{2,\infty}\|Q\|_{1,\infty}$,  $\|PQ\|_{\infty,\infty} = \max_{i,j}|P_{i,:}Q_{:,j}| \le \max_i \|P_{i,:}^\top \|_2 \max_j\|Q_{:,j}\|_2 = \|P^\top \|_{2,\infty}\|Q\|_{2,\infty}$, and  \Cref{lem:lip_softmax},
\begin{allowdisplaybreaks}
\begin{align}
  &\left\|Z^\top \left[A{^{(\cdot,1)}}^\top -\widehat A{^{(\cdot,1)}}^\top , \ldots, A{^{(\cdot,H)}}^\top  - \widehat A{^{(\cdot,H)}}^\top \right]\right\|_{2,\infty} \nonumber\\
  =~& \max_{h \in [H]} \left\|Z^\top \left(A{^{(\cdot,h)}}^\top -\widehat A{^{(\cdot,h)}}^\top \right)\right\|_{2,\infty}\nonumber\\
  \le~& \max_{h\in[H]}\|Z^\top \|_{2,\infty} \left\|A{^{(\cdot,h)}}^\top -\widehat A{^{(\cdot,h)}}^\top \right\|_{1,\infty}\nonumber\\
  \le~& \max_{h\in[H]}2\|Z^\top \|_{2,\infty} \left\|A{^{(\cdot,h)}}^\top -\widehat A{^{(\cdot,h)}}^\top \right\|_{\infty,\infty}\nonumber\\
  =~& \max_{h\in[H]}2\|Z^\top \|_{2,\infty} \left\|P\left(W_K{^{(\cdot,h)}}W_Q{^{(\cdot,h)}}^\top  - \widehat W_K{^{(\cdot,h)}}\widehat W_Q{^{(\cdot,h)}}^\top \right)P^\top \right\|_{\infty,\infty}\nonumber\\
  \le~&\max_{h\in[H]}2\|Z^\top \|_{2,\infty}\|P^\top \|_{2,\infty} \left\|\left(W_Q{^{(\cdot,h)}}W_K{^{(\cdot,h)}}^\top  - \widehat W_Q{^{(\cdot,h)}}\widehat W_K{^{(\cdot,h)}}^\top \right)P^\top \right\|_{2,\infty}\label{eq:Z(A-A_hat)}
\end{align}
\end{allowdisplaybreaks}
Combining \Cref{eq:AZ-A_hatZ_hat,eq:(Z-Z_hat)A_hat,eq:(W-W_hat)AZ,eq:Z(A-A_hat)} gives the result.
\end{proof}

We introduce one more term to simplify the notation. For $\ell = 1,2,\ldots, L$ denote 
\begin{equation}\label{eq:g}
  g^{(\ell)}(X;W^{1:\ell}) = \left(\bigoplus_{h=1}^H A^{(\ell,h)}F^{(\ell-1)}(X;W^{1:\ell-1})W_V^{(\ell,h)}\right) \oplus F^{(\ell-1)}(X;W^{1:\ell-1})
\end{equation}
where $F^{(\ell-1)}(X;W^{1:\ell-1})$ and $A^{(\ell,h)}$ are the same as in \Cref{eq:F}. That is, $g^{(\ell)}(X;W^{1:\ell})$ is the input to the MLP in the $\ell$-th layer of the positional transformer architecture.

\begin{lemma}\label{lem:lip_pos_arch}
For any $W^{1:\ell}$ and $\widehat W^{1:\ell}$,
\begin{allowdisplaybreaks}
\begin{align*}
  &\left\|\left(F^{(\ell)}(X;W^{1:\ell})-F^{(\ell)}(X;\widehat W^{1:\ell})\right)^\top \right\|_{2,\infty}\\
  &\le~ \left\|\left(W_2^{(\ell)}-\widehat W_2^{(\ell)}\right)^\top \sigma\left(g^{(\ell)}(X;\widehat W^{1:\ell})\widehat W_1^{(\ell)}\right)^\top \right\|_{2,\infty}\\
  &\hspace{10mm} + L_\sigma \|W_2^{(\ell)}\|_2 \left\|\left(W_1^{(\ell)}-\widehat W_1^{(\ell)}\right)^\top g^{(\ell)}(X;\widehat W^{1:\ell})^\top \right\|_{2,\infty}\\
  &\hspace{10mm} +  L_\sigma \|W_1^{(\ell)}\|_2 \|W_2^{(\ell)}\|_2 \max\big(\max_{h\in[H]}\|W_V^{(\ell,h)}\|_2, 1\big)\left\|\left(F^{(\ell-1)}(X;W^{1:\ell-1})-F^{(\ell-1)}(X;\widehat W^{1:\ell-1})\right)^\top \right\|_{2,\infty}\\
  &\hspace{10mm} +  2 L_\sigma \|W_1^{(\ell)}\|_2 \|W_2^{(\ell)}\|_2 \|F^{(\ell-1)}(X;W^{1:\ell-1})^\top \|_{2,\infty}\|P^\top \|_{2,\infty}\max\big(\max_{h\in[H]}\|W_V^{(\ell,h)}\|_2, 1\big) \\
  &\hspace{20mm} \cdot \max_{h\in[H]}\left\|\left(W_Q^{(\ell,h)}{W_K^{(\ell,h)}}^\top -\widehat{W}_Q^{(\ell,h)}\widehat{W}_K{^{(\ell,h)}}^\top \right)P^\top \right\|_{2,\infty}\\
  &\hspace{10mm} +  L_\sigma \|W_1^{(\ell)}\|_2 \|W_2^{(\ell)}\|_2 \sqrt{H}\max_{h\in[H]} \left\|\left(W_V^{(\ell,h)}-\widehat W_V^{(\ell,h)}\right)^\top  F^{(\ell-1)}(X;W^{1:\ell-1})^\top \right\|_{2,\infty}
\end{align*}
\end{allowdisplaybreaks}
\end{lemma}

\begin{proof}
Using the definition of $F^{(\ell)}(X;W^{1:\ell})$ in \Cref{eq:F} and the definition of $g^{(\ell)}(X;W^{1:\ell})$ in \Cref{eq:g}:
\begin{allowdisplaybreaks}
\begin{align}
  &\hspace{4.6mm}\left\|\left(F^{(\ell)}(X;W^{1:\ell})-F^{(\ell)}(X;\widehat W^{1:\ell})\right)^\top \right\|_{2,\infty}&&\nonumber\\
  &=~\left\|W_2{^{(\ell)}}^\top \sigma\left(g^{(\ell)}(X;W^{1:\ell})W_1^{(\ell)}\right)^\top  - \widehat W_2{^{(\ell)}}^\top \sigma\left(g^{(\ell)}(X;\widehat W^{1:\ell})\widehat W_1^{(\ell)}\right)^\top \right\|_{2,\infty}\nonumber\\
  &\mbox{(Using triangle inequality):}\nonumber\\
  &\le~\left\|\left(W_2{^{(\ell)}}-\widehat W_2{^{(\ell)}}\right)^\top \sigma\left(g^{(\ell)}(X;\widehat W^{1:\ell})\widehat W_1^{(\ell)}\right)^\top \right\|_{2,\infty}\label{eq:(W-W_hat)g_hatW_hat}\\
  &\hspace{10mm} + ~ \left\|W_2{^{(\ell)}}^\top \left(\sigma\left(g^{(\ell)}(X;W^{1:\ell})W_1^{(\ell)}\right)-\sigma\left(g^{(\ell)}(X;\widehat W^{1:\ell})\widehat W_1^{(\ell)}\right)\right)^\top \right\|_{2,\infty}\label{eq:W(gW-g_hatW_hat)}.
\end{align}
\end{allowdisplaybreaks}
The term in (\ref{eq:(W-W_hat)g_hatW_hat}) gives the first part of the required bound. We focus on the term in (\ref{eq:W(gW-g_hatW_hat)}) below and show that it is bounded by the second to the last parts of the required bound.
\begin{allowdisplaybreaks}
\begin{align*}
  &\hspace{4.6mm}\left\|W_2{^{(\ell)}}^\top \left(\sigma\left(g^{(\ell)}(X;W^{1:\ell})W_1^{(\ell)}\right)-\sigma\left(g^{(\ell)}(X;\widehat W^{1:\ell})\widehat W_1^{(\ell)}\right)\right)^\top \right\|_{2,\infty}\nonumber\\
  &\mbox{(Using $\|PQ\|_{2,\infty}\le\|P\|_2\|Q\|_{2,\infty}$ and $\|P\|_2 = \|P^\top \|_2$):}\nonumber\\
  &\le~\|W_2^{(\ell)}\|_2\left\|\left(\sigma\left(g^{(\ell)}(X;W^{1:\ell})W_1^{(\ell)}\right)-\sigma\left(g^{(\ell)}(X;\widehat W^{1:\ell})\widehat W_1^{(\ell)}\right)\right)^\top \right\|_{2,\infty}\nonumber\\
  &\mbox{(Using Lipschitzness of $\sigma$ and $\sigma(0)=0$):}\nonumber\\
  &\le~L_\sigma \|W_2^{(\ell)}\|_2 \left\|\left(g^{(\ell)}(X;W^{1:\ell})W_1^{(\ell)}-g^{(\ell)}(X;\widehat W^{1:\ell})\widehat W_1^{(\ell)}\right)^\top \right\|_{2,\infty}\nonumber\\
  &\mbox{(Using triangle inequality):}\nonumber\\
  &\le~L_\sigma \|W_2^{(\ell)}\|_2 \left\|\left(W_1^{(\ell)}-\widehat W_1^{(\ell)}\right)^\top  g^{(\ell)}(X;\widehat W^{1:\ell})^\top \right\|_{2,\infty}\nonumber\\
  &\hspace{10mm} + L_\sigma \|W_2^{(\ell)}\|_2\left\|W_1{^{(\ell)}}^\top \left(g^{(\ell)}(X;W^{1:\ell}) - g^{(\ell)}(X;\widehat W^{1:\ell})\right)^\top \right\|_{2,\infty}\nonumber\\
  &\mbox{(Using $\|PQ\|_{2,\infty}\le\|P\|_2\|Q\|_{2,\infty}$ and $\|P\|_2 = \|P^\top \|_2$):}\nonumber\\
  &\le~L_\sigma \|W_2^{(\ell)}\|_2 \left\|\left(W_1^{(\ell)}-\widehat W_1^{(\ell)}\right)^\top  g^{(\ell)}(X;\widehat W^{1:\ell})^\top \right\|_{2,\infty}\nonumber\\
  &\hspace{10mm} + L_\sigma \|W_1^{(\ell)}\|_2\|W_2^{(\ell)}\|_2\left\|\left(g^{(\ell)}(X;W^{1:\ell}) - g^{(\ell)}(X;\widehat W^{1:\ell})\right)^\top \right\|_{2,\infty}\nonumber\\
  &\mbox{(Using \Cref{lem:lip_pos_attn} and definition of $g^{(\ell)}(X;W^{1:\ell})$):}\nonumber\\
  &\le~L_\sigma \|W_2^{(\ell)}\|_2 \left\|\left(W_1^{(\ell)}-\widehat W_1^{(\ell)}\right)^\top  g^{(\ell)}(X;\widehat W^{1:\ell})^\top \right\|_{2,\infty}\nonumber\\
  &\hspace{10mm} + ~ L_\sigma \|W_1^{(\ell)}\|_2 \|W_2^{(\ell)}\|_2 \max\big(\max_{h\in[H]}\|W_V^{(\ell,h)}\|_2, 1\big) \cdot \Bigg(\left\|\left(F^{(\ell-1)}(X;W^{1:\ell-1})-F^{(\ell-1)}(X;\widehat W^{1:\ell-1})\right)^\top \right\|_{2,\infty}\nonumber\\
  &\hspace{25mm}  + ~  2 \|F^{(\ell-1)}(X;W^{1:\ell-1})^\top \|_{2,\infty}\|P^\top \|_{2,\infty} \max_{h\in[H]}\left\|\left(W_Q^{(\ell,h)}{W_K^{(\ell,h)}}^\top -\widehat{W}_Q^{(\ell,h)}\widehat{W}_K{^{(\ell,h)}}^\top \right)P^\top \right\|_{2,\infty}\Bigg)\\
  &\hspace{10mm} + ~ L_\sigma \|W_1^{(\ell)}\|_2 \|W_2^{(\ell)}\|_2\sqrt{H}\max_{h\in[H]} \left\|\left(W_V^{(\ell,h)}-\widehat W_V^{(\ell,h)}\right)^\top  F^{(\ell-1)}(X;W^{1:\ell-1})^\top \right\|_{2,\infty}.
\end{align*}
\end{allowdisplaybreaks}
This gives the required result.
\end{proof}

\subsection{Covering Number Upper Bound}
\label{app:gen_cov}

Let $X^{(1)}, X^{(2)},\ldots, X^{(m)} \in \mathbb{R}^{n \times d}$ be a collection of inputs. Let $\mathcal{F}$ be the class of functions $\mathbb{R}^{n \times d} \rightarrow \mathbb{R}$ as defined in \Cref{eq:bounded_norm_arch_class}. Given $\epsilon > 0$, let $\mathcal{C} \subset \mathcal{F}$ be such that, for all $f\in\mathcal{F}$ there is $\hat{f} \in \mathcal{C}$ such that
\[
  \max_{i \in [m]} |f(X^{(i)}) - \hat{f}(X^{(i)})| \le \epsilon.
\]
Our goal is to bound the size of the cover set $\mathcal{C}$. To do this we will inductively construct $\mathcal{C}$. Let us start with a simple bound on the norm of the output. 

For $W^{1:1}$ satisfying the norm bounds in \Cref{eq:bounded_norm_arch_class}, using $\|Pv\|_2 \le \|P\|_2\|v\|_2$ repeatedly, and the fact that $\sigma$ is $L_\sigma$-Lipschitz with $\sigma(0)=0$,
\[
  \left\|F^{(1)}(X;W^{1:1})^\top \right\|_{2,\infty} \le L_\sigma\|W_2^{(1)}\|_2\|W_1^{(1)}\|_2\sqrt{1+\sum_{h\in H}\|W_V^{(1,h)}\|_2^2}\|X^\top \|_{2,\infty} \le L_\sigma B_2^2\sqrt{1+HB_2^2}\|X^\top \|_{2,\infty},
\]
and inductively for each $\ell \in [L]$ and $W^{1:\ell}$ satisfying the norm bounds in \Cref{eq:bounded_norm_arch_class} we get
\[
  \left\|F^{(\ell)}(X;W^{1:\ell})^\top \right\|_{2,\infty} \le B^{(\ell)}\|X^\top \|_{2,\infty}, \; \mbox{where} \; B^{(\ell)} := L_{\sigma}^\ell B_2^{2\ell}(1+HB_2^2)^{\ell/2}.
\]

Let us now build a cover $\mathcal{C}^{(1)}$ for the first layer of the architecture. For any $X^{(1)}, X^{(2)},\ldots, X^{(m)} \in \mathbb{R}^{n \times d}$ satisfying $\|X{^{(i)}}^\top \|_{2,\infty}\le B_X$ for all $i$, and $W,\widehat W \in \mathcal{W} = \{W \in \mathbb{R}^{d \times d'}: \|W^\top \|_{2,1} \le B_W\}$, because
\[
  \max_{i\in[m]}\left\|(W - \widehat W)^\top  X{^{(i)}}^\top \right\|_{2,\infty} = \max_{i\in[m],j\in[n]}\left\|(W - \widehat W)^\top X{^{(i)}_{j,:}}^\top \right\|_2,
\]
we may apply \Cref{lem:basic_cover} and get that there is an $\varepsilon_W$-cover $\widehat{\mathcal{W}}\subset \mathbb{R}^{d\times d'}$ such that
\[
  \log |\widehat{\mathcal{W}}| \lesssim \frac{B_W^2B_X^2}{\varepsilon_W^2}\log(dmn),
  \quad\mbox{and}\quad
  \max_{W \in \mathcal{W}}\min_{\widehat W \in \widehat{\mathcal{W}}}\max_{i\in[m]}\left\|(W - \widehat W)^\top  X{^{(i)}}^\top \right\|_{2,\infty} \le \varepsilon_W.
\]
Therefore, given inputs $X^{(1)}, X^{(2)},\ldots, X^{(m)} \in \mathbb{R}^{n \times d}$ such that $\|X{^{(i)}}^\top \|_{2,\infty}\le B_X$ for all $i$, using \Cref{lem:basic_cover} we get for each $h\in[H]$ there is an $\varepsilon_V^{(1)}$-cover $\mathcal{C}_V^{(1,h)} \subset \mathbb{R}^{d \times d}$ with
\[
   \log \left|\mathcal{C}^{(1,h)}_V\right| \lesssim \frac{B_{2,1}^2B_X^2}{\varepsilon{^{(1)}_V}^2}\log(dmn)
\]
such that for all $h\in[H]$ and $W_V^{(1,h)}$ satisfying $\|W_V^{(1,h)}\|_{2,1} \le B_{2,1}$ there is $\widehat W_V^{(1,h)}\in \mathcal{C}_V^{(1,h)}$ such that
\[
  \max_{i\in[m]}\left\|\left(W_V^{(1,h)}-\widehat W_V^{(1,h)}\right)^\top X{^{(i)}}^\top \right\|_{2,\infty} \le \varepsilon{^{(1)}_V}.
\]
It follows that there is an $\varepsilon_V^{(1)}$-cover $\mathcal{C}_V^{(1)} \subset \bigotimes_{h=1}^H\mathbb{R}^{d \times d}$ with 
\[
  |\mathcal{C}_V^{(1)}| \le \prod_{h=1}^H |\mathcal{C}_V^{(1,h)}| \quad \implies \quad \log \left|\mathcal{C}^{(1)}_V\right| \lesssim \frac{HB_{2,1}^2B_X^2}{\varepsilon{^{(1)}_V}^2}\log(dmn)
\]
such that for any $W_V^{(1,1)},W_V^{(1,2)},\ldots,W_V^{(1,H)}$ satisfying $\|W_V^{(1,h)}\|_{2,1} \le B_{2,1}$ for all $h$, there is $(\widehat W_V^{(1,1)},\ldots,\widehat W_V^{(1,H)})\in \mathcal{C}_V^{(1)}$ such that
\[
  \max_{h\in[H]}\max_{i\in[m]}\left\|\left(W_V^{(1,h)}-\widehat W_V^{(1,h)}\right)^\top X{^{(i)}}^\top \right\|_{2,\infty} \le \varepsilon{^{(1)}_V}.
\]
Similarly, there is an $\varepsilon_{QK}^{(1)}$-cover $\mathcal{C}_{QK}^{(1)} \subset \prod_{h=1}^H\mathbb{R}^{d_P \times d_P}$ with
\[
   \log \left|\mathcal{C}_{QK}^{(1)}\right| \lesssim \frac{HB_{2,1}^2B_P^2}{\varepsilon{^{(1)}_{QK}}^2}\log(d_Pn)
\]
such that for any $W_Q^{(1,1)}W_K{^{(1,1)}}^\top ,W_Q^{(1,2)}W_K{^{(1,2)}}^\top ,\ldots,W_Q^{(1,H)}W_K{^{(1,H)}}^\top $ satisfying $\|W_Q^{(1,h)}W_K{^{(1,h)}}\|_{2,1} \le B_{2,1}$ for all $h$, there is $(\widehat W_Q^{(1,1)}\widehat W_K{^{(1,1)}}^\top ,\ldots,\widehat W_Q^{(1,H)}\widehat W_K{^{(1,H)}}^\top )\in \mathcal{C}_{QK}^{(1)}$ such that
\[
  \max_{h\in[H]}\left\|\left(W_Q^{(1,h)}W_K{^{(1,h)}}^\top -\widehat W_Q^{(1,h)}\widehat W_K{^{(1,h)}}^\top \right)^\top  P^\top \right\|_{2,\infty} \le \varepsilon{^{(1)}_{QK}}.
\]
Given $(\widehat W_Q^{(1,1)}\widehat W_K{^{(1,1)}}^\top ,\ldots,\widehat W_Q^{(1,H)}\widehat W_K{^{(1,H)}}^\top )\in \mathcal{C}_{QK}^{(1)}$ and $(\widehat W_V^{(1,1)},\ldots,\widehat W_V^{(1,H)})\in \mathcal{C}_V^{(1)}$, we have that $\|g^{(1)}(X;\widehat{W}^{1:1})\|_{2,\infty}\le B^{(1)}B_X$. Therefore, using \Cref{lem:basic_cover} there is an $\varepsilon^{(1)}_{W_1}$-cover $\mathcal{C}^{(1)}_{W_1} \subset \mathbb{R}^{d(H+1) \times d}$ with
\[
  \log \left|\mathcal{C}^{(1)}_{W_1}\right| \lesssim \frac{B{^{(1)}}^2 B_{2,1}^2 B_X^2}{\varepsilon{^{(1)}_{W_1}}^2}\log(dHmn)
\]
such that for any $W_1^{(1)}$ satisfying $\|W_1^{(1)}\|_{2,1}\le B_{2,1}$ there is $\widehat W_1^{(1)} \in \mathcal{C}^{(1)}_{W_1}$ such that
\[
  \left\|\left(W_1^{(1)} - \widehat W_1^{(1)}\right)^\top  g^{(1)}(X;\widehat W^{1:1})^\top \right\|_{2,\infty} \le \varepsilon{^{(1)}_{W_1}}.
\]
Similarly, given $(\widehat W_Q^{(1,1)}\widehat W_K{^{(1,1)}}^\top ,\ldots,\widehat W_Q^{(1,H)}\widehat W_K{^{(1,H)}}^\top )\in \mathcal{C}_{QK}^{(1)}$, $(\widehat W_V^{(1,1)},\ldots,\widehat W_V^{(1,H)})\in \mathcal{C}_V^{(1)}$, and $\widehat W_1^{(1)}\in\mathcal{C}_{W_1}^{(1)}$, there is an $\varepsilon^{(1)}_{W_2}$-cover $\mathcal{C}^{(1)}_{W_2} \in \mathbb{R}^{d\times d}$ with
\[
  \log \left|\mathcal{C}^{(1)}_{W_2}\right| \lesssim \frac{B{^{(1)}}^2 B_{2,1}^2 B_X^2}{\varepsilon{^{(1)}_{W_2}}^2}\log(dmn)
\]
such that for any $W_2^{(1)}$ satisfying $\|W_2^{(1)}\|_{2,1}\le B_{2,1}$ there is $\widehat W_2^{(1)} \in \mathcal{C}^{(1)}_{W_2}$ such that
\[
  \left\|\left(W_2^{(1)}-\widehat W_2^{(1)}\right)^\top \sigma\left(g^{(1)}(X;\widehat W^{1:1})\widehat W_1^{(1)}\right)^\top \right\|_{2,\infty} \le \varepsilon^{(1)}_{W_2}.
\]
Therefore, using \Cref{lem:lip_pos_arch}, an $\varepsilon^{(1)}$-cover $\mathcal{C}^{(1)}$ for the first layer of the architecture, where
\[
  \varepsilon^{(1)} = \varepsilon^{(1)}_{W_2} + L_\sigma B_2\varepsilon^{(1)}_{W_1} + 2 L_\sigma B_2^3  B_P \varepsilon^{(1)}_{QK} + \sqrt{H}L_\sigma B_2^2 \varepsilon^{(1)}_{V},
\]
has size
\begin{align*}
  &|\mathcal{C}^{(1)}| 
  \le 
  |\mathcal{C}^{(1)}_{V}| 
  \cdot 
  |\mathcal{C}^{(1)}_{QK}| 
  \cdot 
  \max\left\{\left|\mathcal{C}_{W_1}^{(1)}\left(\widehat W_V^{(1,h)}, \widehat W_Q^{(1,h)}\widehat W_K{^{(1,h)}}^\top ,\forall h\right)\right|:\left\{\widehat W_V^{(1,h)}\right\}_{h=1}^H \in \mathcal{C}^{(1)}_{V}, \left\{\widehat W_Q^{(1,h)} \widehat W_K{^{(1,h)}}^\top \right\}_{h=1}^H \in \mathcal{C}^{(1)}_{QK}\right\}\\
  &\cdot \max\left\{\left|\mathcal{C}_{W_2}^{(1)}\left(\widehat W_V^{(1,h)}, \widehat W_Q^{(1,h)}\widehat W_K{^{(1,h)}}^\top ,\forall h; \widehat W_1^{(1)}\right)\right|:\left\{\widehat W_V^{(1,h)}\right\}_{h=1}^H \in \mathcal{C}^{(1)}_{V}, \left\{\widehat W_Q^{(1,h)} \widehat W_K{^{(1,h)}}^\top \right\}_{h=1}^H \in \mathcal{C}^{(1)}_{QK}, \widehat W_1^{(1)}\in\mathcal{C}_{W_1}^{(1)}\right\}.
\end{align*}
This means that
\[
  \log |\mathcal{C}^{(1)}|\lesssim \left(\frac{HB_{2,1}^2B_X^2}{\varepsilon{^{(1)}_{V}}^2} + \frac{HB_{2,1}^2B_P^2}{\varepsilon{^{(1)}_{QK}}^2} + \frac{B{^{(1)}}^2B_{2,1}^2B_X^2}{\varepsilon{^{(1)}_{W_1}}^2} + \frac{B{^{(1)}}^2B_{2,1}^2B_X^2}{\varepsilon{^{(1)}_{W_2}}^2}\right)\log(H(d+d_P)mn).
\]
We now inductively construct a cover for deeper layers as follows. For each element $\widehat W^{1:\ell} \in \mathcal{C}^{1:\ell}$ where $\mathcal{C}^{1:\ell}$ denotes a cover for the architecture up to (and including) layer $\ell$ on inputs $X^{(1)},\ldots,X^{(m)}$, we construct an $\varepsilon'_{\ell+1}$-cover $\mathcal{C}^{(\ell+1)}(\widehat W^{1:\ell})$ on inputs $\left\{F^{(\ell)}(X^{(i)};\widehat W^{1:\ell})\right\}_{i\in[m]}$ for the $(\ell+1)$-th layer by following the same steps as above, where
\[
  \varepsilon'_{\ell+1} = \varepsilon^{(\ell+1)}_{W_2} + L_\sigma B_2\varepsilon^{(\ell+1)}_{W_1} + 2 L_\sigma B^{(\ell)} B_2^3  B_P \varepsilon^{(\ell+1)}_{QK} + \sqrt{H}L_\sigma B_2^2 \varepsilon^{(\ell+1)}_{V}.
\]
It follows that from \Cref{lem:lip_pos_arch} that
\[
  \mathcal{C}^{1:\ell+1} := \left\{(\widehat W^{1:\ell}, \widehat W^{(\ell)}) : \widehat W^{1:\ell} \in \mathcal{C}^{1:\ell}, \widehat W^{(\ell)}\in\mathcal{C}^{(\ell+1)}(\widehat W^{1:\ell})\right\}
\]
is an $\varepsilon^{(\ell+1)}$-cover of the architecture up to (and including) layer $\ell+1$, where
\[
  \varepsilon^{(\ell+1)} = L_\sigma B_2^3 \varepsilon^{(\ell)} + \varepsilon^{(\ell+1)}_{W_2} + L_\sigma B_2\varepsilon^{(\ell+1)}_{W_1} + 2 L_\sigma B^{(\ell)} B_2^3  B_P \varepsilon^{(\ell+1)}_{QK} + \sqrt{H}L_\sigma B_2^2 \varepsilon^{(\ell+1)}_{V},
\]
and
\[
  |\mathcal{C}^{1:\ell+1}| \le |\mathcal{C}^{1:\ell}| \cdot \max_{\widehat W^{1:\ell} \in \mathcal{C}^{1:\ell}} |\mathcal{C}^{(\ell+1)}(\widehat W^{1:\ell})|.
\]
By induction, we obtain an $\varepsilon^{(L)}$-cover $\mathcal{C}^{1:L}$ of the function $F^{(L)}(X;W^{1:L})$ on inputs $X^{(1)},\ldots,X^{(m)}$
with
\[
  \varepsilon^{(L)} = \sum_{\ell=1}^L \alpha^{(\ell)} \left( \varepsilon_{W_2}^{(\ell)} + \beta^{(\ell)}\varepsilon_{W_1}^{(\ell)} + \gamma^{(\ell)}\varepsilon_{QK}^{(\ell)} + \eta^{(\ell)}\varepsilon_V^{(\ell)}\right)
\]
where
\[
 \alpha^{(\ell)} = L_\sigma^{\ell-1} B_2^{3\ell-3}, \quad
 \beta^{(\ell)} = L_\sigma B_2, \quad
 \gamma^{(\ell)} = 2 L_\sigma B^{(\ell-1)} B_2^3 B_P, \quad
 \eta^{(\ell)} = \sqrt{H}L_\sigma B_2^2.
\]
The size of the cover satisfies
\[
  \log \big|\mathcal{C}^{1:L}\big|\lesssim \sum_{\ell=1}^L \left(\frac{HB_{2,1}^2B_X^2}{\varepsilon{^{(\ell)}_{V}}^2} + \frac{HB{^{(\ell-1)}}^2B_{2,1}^2B_P^2}{\varepsilon{^{(\ell)}_{QK}}^2} + \frac{B{^{(\ell)}}^2B_{2,1}^2B_X^2}{\varepsilon{^{(\ell)}_{W_1}}^2} + \frac{B{^{(\ell)}}^2B_{2,1}^2B_X^2}{\varepsilon{^{(\ell)}_{W_2}}^2}\right)\log(H(d+d_P)mn).
\]
To build the final cover $\mathcal{C}$ for the class of scalar-output architectures $\mathcal{F}$, note that for any $W^{1:L}, \widehat W^{1:L}, w, \widehat w$,
\begin{align*}
  &\left|w^\top  \left(F^{(L)}(X; W^{1:L})[n]\right) - \widehat w^\top  \left(F^{(L)}(X; \widehat W^{1:L})[n]\right)\right|\\
  \le~
  &\left|w^\top  \left(F^{(L)}(X; W^{1:L})[n] - F^{(L)}(X; \widehat W^{1:L})[n]\right)\right| + \left|(w - \widehat w)^\top  \left(F^{(L)}(X; \widehat W^{1:L})[n]\right)\right|.
\end{align*}
Therefore, we may apply \Cref{lem:basic_cover} and build $\mathcal{C}$ inductively for every element in $\mathcal{C}^{1:L}$, and get that $\mathcal{C}$ is an $\varepsilon$-cover for $\mathcal{F}$ with
\begin{equation}\label{eq:cover_epsilon}
  \varepsilon = \varepsilon_w + B_w \sum_{\ell=1}^L \alpha^{(\ell)} \left(\varepsilon_{W_2}^{(\ell)} + \beta^{(\ell)}\varepsilon_{W_1}^{(\ell)} + \gamma^{(\ell)}\varepsilon_{QK}^{(\ell)} + \eta^{(\ell)}\varepsilon_V^{(\ell)}\right)
\end{equation}
and that the size of $\mathcal{C}$ is bounded as
\begin{equation}\label{eq:cover_log_size}
  \log |\mathcal{C}|
  \lesssim 
  \frac{B{^{(L)}}^2B_X^2B_w^2}{\varepsilon_w^2}\log m + \log \big|\mathcal C^{1:L}\big|.
\end{equation}
To determine a good upper bound on $\log |\mathcal C|$, we pick $\varepsilon^{(\ell)}_V, \varepsilon^{(\ell)}_{QK}, \varepsilon^{(\ell)}_{W_1}, \varepsilon^{(\ell)}_{W_2}, \forall \ell \in [L]$, and $\varepsilon_w$ that minimize the right-hand size of \Cref{eq:cover_log_size} while satisfying \Cref{eq:cover_epsilon}. Invoking \Cref{lem:setting_epsilons}, this gives
\begin{allowdisplaybreaks}
\begin{align}
  \log |\mathcal C|
  &\lesssim \frac{1}{\varepsilon^2} 
  \Bigg( 
  \left(B{^{(L)}}^2B_X^2B_w^2\log m\right)^{1/3} 
  + 
  \left(B_w^2B_X^2B_{2,1}^2\log(H(d+d_P)mn)\right)^{1/3} \cdot \nonumber\\
  &\hspace{20mm}
  \sum_{\ell=1}^L \left(H^{1/3}\alpha{^{(\ell)}}^{2/3}\eta{^{(\ell)}}^{2/3} + H^{1/3}B{^{(\ell-1)}}^{2/3}B_P^{2/3}\alpha{^{(\ell)}}^{2/3}\gamma{^{(\ell)}}^{2/3} \right.\nonumber\\
  &\hspace{30mm}\left. + ~ B{^{(\ell)}}^{2/3}\alpha{^{(\ell)}}^{2/3}\beta{^{(\ell)}}^{2/3} + B{^{(\ell)}}^{2/3}\alpha{^{(\ell)}}^{2/3}\right)\Bigg)^3 \nonumber\\
  &\lesssim \frac{H B_w^2 B_X^2 B_P^2 B_{2,1}^2 B{^{(L)}}^2 \alpha{^{(L)}}^2 \left(\beta{^{(L)}}^2 + \gamma{^{(L)}}^2 + \eta{^{(L)}}^2\right)}{\varepsilon^2}\log(H(d+d_P)mn) \nonumber\\
  &\lesssim \frac{L_\sigma^{2L} B_2^{3L} B{^{(L)}}^2 H^2 B_w^2 B_{2,1}^2 B_X^2 B_P^4}{\varepsilon^2} \log(H(d+d_P)mn) \nonumber\\
  &=(H L_\sigma B_2)^{O(L)} \frac{B_w^2 B_{2,1}^2 B_X^2 B_P^4}{\varepsilon^2}\log(H(d+d_P)mn). \label{eq:cover_number_bound}
\end{align}
\end{allowdisplaybreaks}
Letting $\|w\|_2\le B_2$, $\|X^\top\|_{2,\infty}=\|P^\top\|_{2,\infty}=1$ and using $d\ge d_P$ gives
\begin{equation}
\label{eq:cover_number_bound_simp}
  \log |\mathcal{C}| \lesssim \frac{(HL_\sigma B_2)^{O(L)}B_{2,1}^2}{\varepsilon^2} \log(Hdmn).
\end{equation}

Combining \Cref{eq:cover_number_bound_simp} and \Cref{lem:generalization_from_cover} proves \Cref{thm:generalization}.

{\bf Remark.} Note that if we additionally apply layer normalization by projecting the output of each layer to the $\ell_2$ unit ball as in \citet{edelman2022inductive}, such that
\[
  \|g^{(\ell)}(X;W^{1:\ell})^\top \|_{2,\infty}, \|F^{(\ell)}(X;W^{1:\ell})^\top \|_{2,\infty} \le B^{(\ell)} = 1,
\]
then the bound on $\log |\mathcal C|$ in \Cref{eq:cover_number_bound_simp} becomes
\begin{equation}\label{eq:cover_number_bound_with_layer_normalization}
  \log |\mathcal C| \lesssim \frac{(L_\sigma B_2)^{O(L)} H^2 B_{2,1}^2}{\varepsilon^2} \log(Hdmn).
\end{equation}

%

%
%
%
%
%
%
%
%
%
%
%
%
%
%
%
%
%
%
%
%
%
%
%
%
%
%
%
%
%
%
%
%
%
%
%
%
%
%
%
%
%
%
%
%
%
%
%
%
%
%
%
%
%
%
%
%
%
%
%
%
%
%
%
%
%
%
%
%
%
%
%
%
%
%
%
%
%
%
%
%
%
%
%
%
%
%
%
%
%
%
%
%
%
%
%
%
%
%
%
%
%
%
%
%
%
%
%
%
%
%
%
%
%
%
%
%
%
%
%
%
%
%
%
%
%
%
%
%
%
%
%
%
%
%
%
%
%
%
%
%
%
%
%
%
%
%
%
%
%
%
%
%
%
%
%
%
%
%
%
%
%
%
%
%
%
%
%
%
%
%
%
%
%
%
%
%
%
%
%
%
%
%
%
%
%
%
%
%




==================================================
Filename: 994_experiments.tex
==================================================
%
%
\section{Experiments}
\label{appdx:experiments}
This section presents detailed results for the experiments reported in \Cref{sec:experiments} as well as additional results not present in the main paper. 
Most reported results combine in-distribution (ID) and out-of-distribution (OOD) in one plot. This is shown by the scale factor on the x-axis, which indicates the level of out-of-distribution. A scale of one indicates in-distribution, whereas a scale greater than one shows an increasing level of out-of-distribution.

\begin{itemize}[label={}, leftmargin=*]
  \item \ref{app:num} - \textbf{Numeric tasks} 
  \begin{itemize}[label={}]
    \item \ref{app:num_setting} - Experimental setting
    \item \ref{app:num_length} - Input length vs. generalization experiments
    \item \ref{app:num_sample} - Sample complexity experiments
    \item \ref{app:num_var} - Performance over variable input lengths 
    \item \ref{app:num_comp} - Performance of compact Transformers 
    \item \ref{app:num_posenc} - Performance with alternative positional encodings 
    \item \ref{app:num_ablation} - Ablation study on architectural design choices
    \item \ref{app:num_fine_tuning} - Performance on OOD data with fine-tuned V projection matrices
    \item \ref{app:num_gnns} - Empirical results using Graph Neural Networks (GNNs)
    \item \ref{app:num_acc} - Accuracy measures
  \end{itemize}
  \item \ref{app:indheads} - \textbf{$k$-hop induction heads}
  \begin{itemize}[label={}]
    \item \ref{app:indheads_setting} - Experimental setting
    \item \ref{app:indheads_sample} - Sample complexity experiment
    \item \ref{app:indheads_exp} - Expressive power experiment
  \end{itemize}
  \item \ref{app:nlp} - \textbf{Mixed-type input task}
  \begin{itemize}[label={}]
    \item \ref{app:nlp_task} - Task description
    \item \ref{app:nlp_setting} - Experimental setting
    \item \ref{app:nlp_sample} - Sample complexity experiments
    \item \ref{app:nlp_gpt2} - Experiments with GPT2
  \end{itemize}
\end{itemize}

\begin{comment}
Briefly, we examine the following:
\begin{itemize}
    \item The performance of standard and positional Transformers on algorithmic computational tasks over variable-length inputs.
    \item The performance of more compact standard Transformers on algorithmic computational tasks.
    \item The performance of standard and positional Transformers using different types of positional encodings.
    \item The performance of standard and positional Transformers as a function of training samples.
    \item The performance of standard and positional Transformers as a function of input lengths.
    \item The performance of standard and positional Transformers when evaluated on a relational task with mixed-type inputs.
    \item The performance of standard and positional Transformers when evaluated on the previous relational task with additional noise.
\end{itemize}

Finally, we present two new transformations that label each prediction as ``correct" or ``incorrect" and use the labels to calculate the OOD accuracies of both standard and positional Transformer for all tasks.

Unless otherwise specified, all configurations are consistent with those described in  \Cref{sec:experiments}. Both training and testing lists are generated using the sampling strategy discussed in \Cref{sec:experiments}.
\end{comment}

\subsection{Numeric tasks}
\label{app:num}
\subsubsection{Experimental setting}
\label{app:num_setting}
All numeric tasks employ the same model configuration. The model uses the structure of \Cref{eq:arch}, augmented with linear encoding and decoding layers.

We compare the standard Transformer, which utilizes the attention mechanism in \Cref{eq:selfattn}, and the positional Transformer, which employs the attention defined in \Cref{eq:attn}. Both variants share the same number of layers and dimensional configurations, with any specific differences explicitly noted.

In all configurations, the total number of layers is set to $\lceil\log_2 n\rceil +1$, where $n$ denotes the maximum input length, and each layer uses $2$ attention heads.
Along with each input sequence, we also append an empty scratchpad entry. This extra entry does not count toward the total number of layers and is not used to compute the loss. It is included solely to aid in the computation of the tasks. For the function $\Phi^{(\ell)}$, we employ a 2-layer MLP with ReLU activation functions. The embedding dimension of the encoder and the hidden dimensions of the MLP are both set to $64$. 

We use one-hot encoded vectors of dimension $n$ for positional encodings, where the non-zero entry corresponds to the node position. Consequently, the embedding dimensions of $W_Q$ and $W_K$ are set to $n$. A key difference between the models is that standard Transformers concatenate positional encodings to the input, whereas positional Transformers supply positional information exclusively through the matrix $P$. Therefore, in positional Transformers, input values are solely encoded in the input, and positional information is exclusively encoded in the positional encoding matrix.

Both models are trained end-to-end 
using the squared loss between the predicted and target vectors of size $n$, with no intermediate supervision.
We train models with Adam, starting with a learning rate of $5\cdot 10^{-4}$ and a learning rate scheduler for a total of 2000 epochs.

Our training data consists of samples from the range $[-2, 2]$. To ensure diversity in the data, for each input sample, we first select lower and upper bounds $\gamma_l$ and $\gamma_u$ uniformly in $[-2, 2]$, and then for each of the $n$ elements of the input sample, we select its value uniformly from  $[\gamma_l, \gamma_u]$.

\subsubsection{Input length vs. Generalization experiments}
\label{app:num_length}
In this section, we validate that our results hold for multiple values of $n$. We train models for each fixed length $n \in \{2, 4, 8, 16, 32\}$ using $30,\!000$ samples across all settings. The model depth varies with the input length $n$, with the number of layers set to $\lceil \log_2 n \rceil + 1$. We report the performance in-distribution and out-of-distribution using $1,\!000$ test samples. As illustrated across \Cref{fig:app_sum_length,fig:app_min_length,fig:app_median_length,fig:app_sort_length,fig:app_maxsub_length}, positional Transformers exhibit comparable in-distribution loss to standard Transformers across various sequence lengths. Naturally, for a fixed number of samples, the in-distribution loss slightly increases as the sequence length grows, indicating the need for more samples.
For out-of-distribution performance, we observe that the positional Transformer maintains a more stable performance even with inputs of length $n=32$. In contrast, the standard Transformer's performance rapidly degrades with the input length.

\begin{figure}[htb]
    \centering
    \includegraphics[width=0.9\linewidth]{plots/fixed_length/sum/size.pdf}
    \vspace{-4mm}
    \caption{Training, in-distribution, and out-of-distribution test performance for the summing task is shown for standard Transformers (red) and positional Transformers (blue) across different input lengths. The x-axis indicates the fixed input length on which the model was trained. Models are trained on the range $[-2, 2]$ with 30,000 samples, validated on the same domain, and tested on an extended domain, $[-6, 6]$, each with 1,000 samples.}
    \label{fig:app_sum_length}
\end{figure}

\vspace{-5mm}

\begin{figure}[hbt]
    \centering
    \includegraphics[width=0.9\linewidth]{plots/fixed_length/min/size.pdf}
    \vspace{-4mm}
    \caption{Training, in-distribution, and out-of-distribution test performance for the minimum task is shown for standard Transformers (red) and positional Transformers (blue) across different input lengths. The x-axis indicates the fixed input length on which the model was trained. Models are trained on the range $[-2, 2]$ with 30,000 samples, validated on the same domain, and tested on an extended domain, $[-6, 6]$, each with 1,000 samples.}
    \label{fig:app_min_length}
\end{figure}

\begin{figure}[H]
    \centering
\includegraphics[width=0.9\linewidth]{plots/fixed_length/median/size.pdf}
    \vspace{-4mm}
    \caption{Training, in-distribution, and out-of-distribution test performance for the median task is shown for standard Transformers (red) and positional Transformers (blue) across different input lengths. The x-axis indicates the fixed input length on which the model was trained. Models are trained on the range $[-2, 2]$ with 30,000 samples, validated on the same domain, and tested on an extended domain, $[-6, 6]$, each with 1,000 samples.}
    \label{fig:app_median_length}
\end{figure}

\begin{figure}[H]
    \centering
    \includegraphics[width=0.9\linewidth]{plots/fixed_length/sort/size.pdf}
    \caption{Training, in-distribution, and out-of-distribution test performance for the sorting task is shown for standard Transformers (red) and positional Transformers (blue) across different input lengths. The x-axis indicates the fixed input length on which the model was trained. Models are trained on the range $[-2, 2]$ with 30,000 samples, validated on the same domain, and tested on an extended domain, $[-6, 6]$, each with 1,000 samples.}
    \label{fig:app_sort_length}
\end{figure}

\begin{figure}[H]
    \centering
    \includegraphics[width=0.9\linewidth]{plots/fixed_length/max_subarray/size.pdf}
    \caption{Training, in-distribution, and out-of-distribution test performance for the maximum sum subarray task is shown for standard Transformers (red) and positional Transformers (blue) across different input lengths. The x-axis indicates the fixed input length on which the model was trained. Models are trained on the range $[-2, 2]$ with 30,000 samples, validated on the same domain, and tested on an extended domain, $[-6, 6]$, each with 1,000 samples.}
    \label{fig:app_maxsub_length}
\end{figure}


\subsubsection{Sample size vs. Generalization experiments}
\label{app:num_sample}
In this section, we provide detailed results showcasing the training, in-distribution, and out-of-distribution test performance for each of the five numerical tasks as a function of the number of training samples used. From the results, we can draw two conclusions about the behavior of the models as the number of samples increases. First, both models achieve good in-distribution performance. Second, only the positional Transformer achieves better OOD performance. The results for this experiment are presented in \Cref{fig:app_sum_sample,fig:app_min_sample,fig:app_median_sample,fig:app_sort_sample,fig:app_maxsub_sample}.

\begin{figure}[htb]
    \centering
    \includegraphics[width=0.9\linewidth]
    {plots/fixed_length/sum/sample.pdf}
    \caption{Training, in-distribution, and out-of-distribution test performance for the summing task is shown for standard Transformers (red) and positional Transformers (blue) as a function of the number of training samples (indicated on the x-axis). Models are trained on the range $[-2, 2]$ with varying training set sizes. In-distribution testing is performed on the same domain, and testing is conducted on an extended domain, $[-6, 6]$, each using 1,000 samples.}
    \label{fig:app_sum_sample}
\end{figure}

\begin{figure}[hbt]
    \centering
    \includegraphics[width=0.9\linewidth]{plots/fixed_length/min/sample.pdf}
    \caption{Training, in-distribution, and out-of-distribution test performance for the minimum task is shown for standard Transformers (red) and positional Transformers (blue) as a function of the number of training samples (indicated on the x-axis). Models are trained on the range $[-2, 2]$ with varying training set sizes. In-distribution testing is performed on the same domain, and testing is conducted on an extended domain, $[-6, 6]$, each using 1,000 samples.}
    \label{fig:app_min_sample}
\end{figure}

\begin{figure}[H]
    \centering
\includegraphics[width=0.9\linewidth]{plots/fixed_length/median/sample.pdf}
    \caption{Training, in-distribution, and out-of-distribution test performance for the median task is shown for standard Transformers (red) and positional Transformers (blue) as a function of the number of training samples (indicated on the x-axis). Models are trained on the range $[-2, 2]$ with varying training set sizes. In-distribution testing is performed on the same domain, and testing is conducted on an extended domain, $[-6, 6]$, each using 1,000 samples.}
    \label{fig:app_median_sample}
\end{figure}

\begin{figure}[H]
    \centering
    \includegraphics[width=0.9\linewidth]{plots/fixed_length/sort/sample.pdf}
    \caption{Training, in-distribution, and out-of-distribution test performance for the sorting task is shown for standard Transformers (red) and positional Transformers (blue) as a function of the number of training samples (indicated on the x-axis). Models are trained on the range $[-2, 2]$ with varying training set sizes. In-distribution testing is performed on the same domain, and out-of-distribution testing is conducted on an extended domain, $[-6, 6]$, each using 1,000 samples.}
    \label{fig:app_sort_sample}
\end{figure}

\begin{figure}[H]
    \centering
    \includegraphics[width=0.9\linewidth]{plots/fixed_length/max_subarray/sample.pdf}
    \caption{Training, in-distribution, and out-of-distribution test performance for the maximum sum subarray task is shown for standard Transformers (red) and positional Transformers (blue) as a function of the number of training samples (indicated on the x-axis). Models are trained on the range $[-2, 2]$ with varying training set sizes. In-distribution testing is performed on the same domain, and out-of-distribution testing is conducted on an extended domain, $[-6, 6]$, each using 1,000 samples.}
    \label{fig:app_maxsub_sample}
\end{figure}

\subsubsection{Performance over variable input lengths}
\label{app:num_var}
In this section, we present generalization results for models operating on variable-length inputs. This setting aims to verify the models' ability to generalize across different scales while maintaining the flexibility to handle inputs of varying lengths.

We evaluate the models' ability to process sequences of varying lengths up to a maximum size of $n = 8$ Specifically, the model is required to perform tasks on input sequences with lengths ranging from 1 to 8. 
Due to the more challenging setting, we train models with 500,000 samples and ensure that all input lengths are equally represented.

We then evaluate the in-distribution (ID) and out-of-distribution (OOD) loss across different scale factors $c \in \{1, 2, \dots, 10\}$. Note that when $c = 1$, the setting actually corresponds to ID generalization. The losses reported are calculated using 3,000 samples for each scale. As shown in \Cref{fig:var_scale}, positional Transformers maintain the same in-distribution performance while consistently outperforming standard Transformers across all out-of-distribution scales and tasks. Additionally, our architecture maintains robust OOD performance even in tasks where the output can exceed the input magnitude (e.g., sum and maximum sum subarray).

\begin{figure}[ht]
    \centering
\includegraphics[width=\linewidth]{plots/variable_length/all/scale.pdf}
    \caption{ID/OOD loss (measured as mean squared error, MSE) for standard Transformers (red) and positional Transformers (blue) across all five tasks for \emph{variable} lengths (up to $n = 8$). The x-axis represents the ID/OOD scale factor. The solid line and shaded area denote the median and the region between the 10\textsuperscript{th} and 90\textsuperscript{th} percentiles over ten trials, respectively.}
    \label{fig:var_scale}
\end{figure}


\subsubsection{Performance of compact Transformers}
\label{app:num_comp}
This section presents additional generalization results for simpler models, aiming to rule out potential overfitting caused by the excessive complexity of the standard Transformer. We examine two configuration variants: one with $\log n + 1$ layers ($4$ layers) and another with a single layer. The plots also illustrate the outcomes for different hidden dimensions in the MLP. We report the generalization results for the cumulative sum (\Cref{fig:combined_scale_sum}) and cumulative minimum (\Cref{fig:combined_scale_min}) tasks. As observed, in both cases, the OOD performance deteriorates as the network size decreases. Although single-layer networks exhibit slightly better performance, they remain inferior to the performance of positional attention reported in the main paper.

\begin{figure}[hbt]
    \centering
    \begin{minipage}{0.45\linewidth}
        \centering
        \includegraphics[width=\linewidth]{plots/small_transformers/sum/log_layers/scale.pdf}
    \end{minipage}
    \hspace{0.05\linewidth} %
    \begin{minipage}{0.45\linewidth}
        \centering
        \includegraphics[width=\linewidth]{plots/small_transformers/sum/one_layer/scale.pdf}
    \end{minipage}
    \caption{ID/OOD loss for various standard Transformer models on the cumulative sum task with fixed length ($n = 8$). The left plot displays results for models with 4 layers, while the right plot shows results for single-layer models, both featuring varying hidden dimensions in the MLPs. The x-axis represents the out-of-distribution scale factor, indicating the distance from the training distribution. The solid lines and shaded areas denote the median and the regions between the 10\textsuperscript{th} and 90\textsuperscript{th} percentiles across ten trials, respectively.}
    \label{fig:combined_scale_sum}
\end{figure}

\begin{figure}[htb]
    \centering
    \begin{minipage}{0.45\linewidth}
        \centering
        \includegraphics[width=\linewidth]{plots/small_transformers/min/log_layers/scale.pdf}
    \end{minipage}
    \hspace{0.05\linewidth} %
    \begin{minipage}{0.45\linewidth}
        \centering
        \includegraphics[width=\linewidth]{plots/small_transformers/min/one_layer/scale.pdf}
    \end{minipage}
    \caption{ID/OOD loss for various standard Transformer models on the cumulative minimum task with fixed length ($n = 8$). The left plot displays results for models with 4 layers, while the right plot shows results for single-layer models, both featuring varying hidden dimensions in the MLPs. The x-axis represents the out-of-distribution scale factor, indicating the distance from the training distribution. The solid lines and shaded areas denote the median and the regions between the 10\textsuperscript{th} and 90\textsuperscript{th} percentiles across ten trials, respectively.}
    \label{fig:combined_scale_min}
\end{figure}

\subsubsection{Performance with alternative positional encodings}
\label{app:num_posenc}

We compare the performance of positional Transformers and standard Transformers using alternative positional encodings, such as binary and sinusoidal encodings \cite{vaswani2017attention} and Rotary Positional Embedding (RoPE) \citep{10.1016/j.neucom.2023.127063}, a widely adopted technique in natural language processing contexts that has also been applied to algorithmic tasks \citep{bounsi2024transformers}.

For binary positional encodings, we use $\lceil \log_2 n \rceil$ dimensions, where each entry represents the binary encoding of its index in $\lceil \log_2 n \rceil$ bits, with zeros encoded as $-1$. The result for binary encodings is shown in \Cref{fig:fix_scale_binary}.
For sinusoidal positional encodings, we follow the strategy outlined in \cite{vaswani2017attention}, with the encoding dimension set to $\lceil \nicefrac{n}{2} \rceil$. The result for sinusoidal encodings is shown in \Cref{fig:fix_scale_sine}.
From an expressivity perspective, while these encodings are less expressive than one-hot positional encodings, they maintain consistent out-of-distribution (OOD) performance across all ranges. Furthermore, positional Transformers outperform standard Transformers in every task tested.

As for RoPE, while it manages to decrease the OOD test loss,
this improvement is far from what is observed in positional Transformers across every task. The results of this experiment are presented in \Cref{fig:fix_scale_rope}.

\begin{figure}[H]
    \centering
    \includegraphics[width=\linewidth]{plots/fixed_length/all/scale_binary.pdf}
    \caption{ID/OOD loss (measured as mean squared error, MSE) for standard Transformers (red) and positional Transformers (blue) using binary positional encodings across all five tasks for fixed length ($n = 8$). The x-axis represents the OOD scale factor, indicating the distance from the training distribution. The solid line and shaded area denote the median and the region between the 10\textsuperscript{th} and 90\textsuperscript{th} percentiles over ten trials, respectively.
    }
    \label{fig:fix_scale_binary}
\end{figure}


\begin{figure}[H]
    \centering
    \includegraphics[width=\linewidth]{plots/fixed_length/all/scale_sinusoidal.pdf}
    \caption{ID/OOD loss (measured as mean squared error, MSE) for standard Transformers (red) and positional Transformers (blue) using sinusoidal positional encodings across all five tasks for fixed length ($n = 8$). The x-axis represents the OOD scale factor, indicating the distance from the training distribution. The solid line and shaded area denote the median and the region between the 10\textsuperscript{th} and 90\textsuperscript{th} percentiles over ten trials, respectively.
    }
    \label{fig:fix_scale_sine}
\end{figure}


\begin{figure}[H]
    \centering
    \includegraphics[width=\linewidth]{plots/fixed_length/all/scale_rope.pdf}
    \caption{ID/OOD loss (measured as mean squared error, MSE) for standard Transformers with RoPE (red) and positional Transformers (blue) across all five tasks for fixed length ($n = 8$). The x-axis represents the OOD scale factor, indicating the distance from the training distribution. The solid line and shaded area denote the median and the region between the 10\textsuperscript{th} and 90\textsuperscript{th} percentiles over ten trials, respectively.
    }
    \label{fig:fix_scale_rope}
\end{figure}

\subsubsection{Ablation study on architectural design choices}
\label{app:num_ablation}

We carried out additional experiments to explore variations in the input data format and the placement of positional encodings. This leads to the following seven architectural choices (including two from the main paper and five additional variations):

\begin{enumerate}
  \item Standard Transformers: Input numbers and positional encodings are fed to MLPs, value, query, and key matrices.
  \item Standard Transformers with positional input (but no positional encodings): Input numbers are placed in one-hot positions, that is, the input is $\mbox{Diag}(X)$ where $X$ is the list of input numbers we give to other architectures. No additional positional encodings are used.
  \item Positional Transformers: Input numbers are fed to the MLPs and value matrix; positional encodings are fed to query and key matrices.
  \item Misaligned Positional Transformers: Input numbers are fed to the MLPs and value matrix; positional encodings are fed to the MLPs, value, query, and key matrices. That is, compared with Positional Transformers, we add positional encodings to the input.
  \item Input-regularized Standard Transformers: Input numbers are fed to the MLPs, value, query, and key matrices; positional encodings are only used in query and key matrices.
  \item No Positional Encodings: input numbers are fed to the MLPs, value, query, and key matrices. No positional encodings are used.
  \item Using RoPE Only: Input numbers are fed to MLPs, value, query, and key matrices, removing absolute positional encodings, and using only RoPE in standard transformers.
\end{enumerate}
We test the architectures for the cumulative sum task and the sorting task as representative tasks, for fixed input length $n=8$. We keep the train/valid/test setup the same as before. The results are shown in \Cref{fig:model_variations_sum} and \Cref{fig:model_variations_sort}. We note that our proposed positional Transformer architecture has all of the following three important factors that enabled OOD generalization: (1) use positional encodings, (2) do not use positional encodings in the value matrix of the attention, (3) use only fixed positional encodings in the key and query matrices. This design principle aligns with algorithms that are typically used to solve algorithmic tasks.

\begin{figure}[ht]
    \centering
    \includegraphics[width=\linewidth]{plots/model_variations/sum_model_variations.pdf}
    \caption{ID/OOD loss for the cumulative sum task under various architectural choices. We fix input length $n = 8$. The x-axis represents the OOD scale factor. The solid line and shaded area denote the median and the region between the 10\textsuperscript{th} and 90\textsuperscript{th} percentiles for $10$ trials, respectively.}
    \label{fig:model_variations_sum}
\end{figure}

\begin{figure}[ht]
    \centering
    \includegraphics[width=\linewidth]{plots/model_variations/sort_model_variations.pdf}
    \caption{ID/OOD loss for the sorting task under various architectural choices. We fix input length $n = 8$. The x-axis represents the OOD scale factor. The solid line and shaded area denote the median and the region between the 10\textsuperscript{th} and 90\textsuperscript{th} percentiles for $10$ trials, respectively.}
    \label{fig:model_variations_sort}
\end{figure}

\subsubsection{Experiments on fine-tuning V projection matrices only}\label{app:num_fine_tuning}

Fine-tuning pre-trained models is very common in practice. In this case, one takes a model that has been trained on a different (and usually larger) dataset, and re-trains a small subset of layers on a new dataset, while keeping all other model parameters frozen. To further study the potential advantage of positional attention, we consider the following fine-tuning experiment. We take both positional Transformer and standard Transformer models that have been trained over in-distribution data (cf. \Cref{sec:ood-regime} where $c=1$, and \Cref{app:num_setting}) for the numeric algorithmic tasks, fix all their model parameters except the weight parameters in the V projection matrices, and fine-tune the V projection matrices over OOD data with 10x scale factor (cf. \Cref{sec:ood-regime} where $c=10$). We consider three algorithmic tasks: cumulative sum, cumulative minimum, and sorting. For these tasks, we find that fine-tuning only the V projection matrices in the OOD case allows the positional Transformer to achieve the same level of generalization as in the in-distribution case, while, on the other hand, the same does not hold for the standard Transformer.

\Cref{tab:fine_tuning_results} reports the median test error (MSE) for both positional and standard transformer models when (i) they are trained and tested on in-distribution data (in-dist), (ii) they are trained on in-distribution data but tested on numbers 10x larger than the original in-distribution data (10x OOD), (iii) they are trained on in-distribution data, fine-tuned V projection matrices on numbers 10x larger than the original in-distribution data, and then tested on numbers 10x larger than the original in-distribution data (10x OOD after fine-tuning). We use the same empirical setting as before, where we train both models for 2000 epochs on in-distribution data. For fine-tuning, we adopt the common practice and train for only 20 epochs. We normalize the MSE by the scale factor so that the results indicate relative accuracy. The results in \Cref{tab:fine_tuning_results} indicate that fine-tuning only the V projection matrices brings the accuracy of the positional Transformer to the same level as the in-distribution case. However, for the standard Transformer, even after fine-tuning the V projection matrices, the accuracy on 10x test data is still very high compared to that of the positional Transformer. The comparison highlights another potential advantage of the positional Transformer for executing algorithmic tasks: one might efficiently fine-tune a pre-trained positional Transformer on OOD data and achieve good performance.

\begin{table}[ht!]
\caption{Performance (MSE test loss) of fine-tuned positional and standard transformer models}
\label{tab:fine_tuning_results}
\centering
\small
\begin{tabular}{ccrrr}
\toprule
Task & Architecture & in-dist & 10x OOD & 10x OOD after fine-tuning \\
\midrule
\multirow{2}{*}{Cumulative Sum} & Positional & 6.07e-06 & 1.31e-03 & 1.32e-05\\
& Standard & 5.20e-06 & 1.09e+00 & 1.23e-01\\
\midrule
\multirow{2}{*}{Cumulative Minimum} & Positional & 1.03e-05 & 3.92e-04 & 2.19e-05\\
& Standard & 1.74e-06 & 1.39e-01 & 4.00e-02\\
\midrule
\multirow{2}{*}{Sorting} & Positional & 1.20e-04 & 7.37e-04 & 1.23e-04\\
& Standard & 9.05e-04 & 5.18e-01 & 6.90e-02\\
\bottomrule
\end{tabular}
\end{table}

\subsubsection{Empirical results using Graph Neural Networks (GNNs)}
\label{app:num_gnns}

Graph Neural Networks are a popular choice for solving algorithmic tasks on graphs. We tested the performance of Graph Convolutional Network (GCN) and Graph Attention Network (GAT) in solving the cumulative sum and sorting tasks. Since the tasks tested have no underlying native graph, we tested these models on complete and star graphs. Notably, the original GAT architecture on a complete graph is similar to a standard transformer but differs in that the value, key, and query weights are shared in GAT. In \Cref{tab:gnn_results}  we present results reporting the median MSE loss. As the results indicate, neither GCN nor GAT works very well even for in-distribution data (OOD scaling factor = 1), let alone achieving OOD generalization. We believe this is because these tasks are inherently unsuitable for standard message-passing architectures. 

\begin{table}[H]
\caption{ID/OOD loss of Graph Convolutional Network (GCN) and Graph Attention Network (GAT)}
\label{tab:gnn_results}
\centering
\small
\begin{tabular}{cccrrrrrr}
\toprule
& & & \multicolumn{6}{c}{OOD scale factor}\\
\cmidrule(l{2pt}r{2pt}){4-9}
Task & Graph & Architecture & 1 & 2 & 3 & 4 & 5 & 10\\
\midrule
\multirow{4}{*}{Cumulative Sum}  
& \multirow{2}{*}{Complete} & GCN & 3.9e+0 & 2.1e+1 & 4.3e+1 & 7.4e+1 & 1.2e+2  & 4.8e+2 \\
& & GAT & 3.7e+0 & 2.0e+1 & 4.3e+1 & 7.5e+1 & 1.2e+2 & 5.4e+2 \\
\cmidrule(l{2pt}r{2pt}){2-9}
& \multirow{2}{*}{Star} & GCN & 4.0e+0 & 2.0e+1 & 4.2e+1 & 7.4e+1 & 1.2e+2 & 4.6e+2 \\
& & GAT & 3.8e+0 & 2.0e+1 & 4.1e+1 & 7.5e+1 & 1.2e+2 & 5.0e+2\\
\midrule
\multirow{4}{*}{Sorting}
& \multirow{2}{*}{Complete} & GCN & 1.9e-1 & 9.8e-1 & 2.2e+0 & 3.9e+0 & 6.1e+0 & 2.7e+1 \\ 
& & GAT & 1.9e-1 & 9.6e-1 & 2.1e+0 & 3.7e+0 & 5.6e+0 & 2.4e+1 \\
\cmidrule(l{2pt}r{2pt}){2-9}
& \multirow{2}{*}{Star} & GCN & 1.9e-1 & 1.0e+0 & 2.1e+0 & 3.7e+0 & 6.1e+0 & 2.6e+1\\
& & GAT & 1.9e-1 & 9.9e-1 & 2.0e+0 & 3.5e+0 & 5.6e+0 & 2.3e+1 \\
\bottomrule
\end{tabular}
\end{table}

\subsubsection{Accuracy Measures}
\label{app:num_acc}
Given the regressive nature of the tasks considered in this work, the notion of model accuracy is not properly defined. However, we propose the following transformation strategies that assign binary labels (``correct"/``incorrect") to the models' outputs allowing us to measure their accuracy (with respect to these transformations):
\begin{itemize}
    \item Rounding transformation: We evaluate the model on lists containing integers (while training is still done using real numbers) and round the model’s output to the nearest integer (or nearest $0.5$ for the median task). A prediction is considered ``correct" if the rounded and ground truth lists are the same. The generalization accuracies for all arithmetic tasks in the main paper using this transformation strategy are presented in \Cref{fig:acc_strict}.
    \item Closeness transformation: We evaluate the model on lists of real numbers, considering a prediction ``correct" if each entry in the predicted list is within an absolute precision of $0.05$ and a relative precision of $5\%$ compared to the corresponding entry of the ground truth list. The generalization accuracies for all arithmetic tasks in the main paper using this transformation strategy are presented in \Cref{fig:acc_rel}.
\end{itemize}

It is important to note that the above metrics are quite unforgiving, as even a single element in the predicted list failing to meet the corresponding criterion results in the entire list being classified as ``incorrect". It is therefore expected that increasing the OOD scale factor causes the model's accuracy to decrease rapidly. Nevertheless, the positional Transformer significantly outperforms the standard Transformer in terms of accuracy in all five tasks. 

\begin{figure}[H]
    \centering
    \includegraphics[width=\linewidth]{plots/accuracy/all_tasks_integers_scale.pdf}
    \caption{ OOD accuracy when using ``rounding transformation" across all five tasks for standard Transformers (red) and positional Transformers (blue) as a function of the scale factor. The solid line and shaded area denote the median and the region between the 10\textsuperscript{th} and 90\textsuperscript{th} percentiles over ten trials, respectively.}
    \label{fig:acc_strict}
\end{figure}

\begin{figure}[H]
    \centering
    \includegraphics[width=\linewidth]{plots/accuracy/all_tasks_relative_scale.pdf}
    \caption{ OOD accuracy when using ``closeness transformation" across all five tasks for standard Transformers (red) and positional Transformers (blue) as a function of the scale factor. The solid line and shaded area denote the median and the region between the 10\textsuperscript{th} and 90\textsuperscript{th} percentiles over ten trials, respectively.}
    \label{fig:acc_rel}
\end{figure}


%
\subsection{$k$-hop induction heads}
\label{app:indheads}
\subsubsection{Experimental setting}
\label{app:indheads_setting}
We adopt a slight modification of the configuration used in the numeric task for the $k$-hop induction heads task of \citet{pmlr-v235-sanford24a}. Specifically, since the input consists of tokens, we use an embedding layer as the encoding layer, with an embedding dimension of 60. The decoding step remains a linear layer to produce the logits, whose output dimension is set to the vocabulary size. Additionally, due to the sequential nature of the task, we apply a causal mask over the attention coefficients in both architectures.

The dataset creation is adopted from \citet{pmlr-v235-sanford24a} and consists of the following types of characters: numbers from $0$ to $k$, which indicate the number of hops required for each sample (as mentioned in \Cref{sec:experiments}, we set $k = 3$ for our experiments). The sequences themselves are constructed from specific character sets: \texttt{\{a,b,c,d,e\}} for in-distribution and \texttt{\{f,g,h,i,j\}} for out-of-distribution. An additional character ($\sim$) denotes cases where a $k$-hop character cannot be found. To ensure that most output tokens can be reached within the specified number of hops (thus avoiding $\sim$ as an output), we set the sequence length to 30.

\subsubsection{Sample complexity Experiment}
\label{app:indheads_sample}
In this section, we provide detailed results showcasing the training, in-distribution, and out-of-distribution test performance for the $k$-hop induction heads task as a function of the number of training samples used. From the results on \Cref{fig:indheads_sample}, we can observe a sharp contrast between the learnability of the standard and positional Transformers, indicating that, for this particular task, the sample complexity of positional Transformers is higher than that of standard Transformers. As discussed in \Cref{sec:experiments}, this is likely a result of the nature of the induction heads task, which is purposely designed to showcase the relational capabilities of self-attention. As a result, positional Transformers require more data to replicate this same behavior.

\begin{figure*}[hbt]
    \centering
    \includegraphics[width=0.9\linewidth]{plots/induction_heads/indheads_sample_complexity.pdf}
    \caption{Training, in-distribution, and out-of-distribution test performance for the $k$-hop induction heads task are shown for standard Transformers (red) and positional Transformers (blue) as a function of the number of training samples (indicated on the x-axis).}
    \label{fig:indheads_sample}
\end{figure*}

%

\subsubsection{Expressive power experiment}
\label{app:indheads_exp}
In this section, we provide detailed results showcasing the training, in-distribution, and out-of-distribution test performance for the $k$-hop induction heads task as a function of the expressive power of the architectures, measured by the number of layers. From the results presented in \Cref{fig:indheads_layers}, we can conclude that, for this particular task, positional Transformers require many more layers to achieve a similar performance as standard Transformers, which empirically validates our findings established in \Cref{rem:expressivity}.

\begin{figure*}[hbt]
    \centering
    \includegraphics[width=0.9\linewidth]{plots/induction_heads/indheads_layers.pdf}
    \caption{Training, in-distribution, and out-of-distribution test performance for the $k$-hop induction heads task are shown for standard Transformers (red) and positional Transformers (blue) as a function of the number of layers (indicated on the x-axis).}
    \label{fig:indheads_layers}
\end{figure*}

\subsection{Mixed-type input task}
\label{app:nlp}

In this section, we provide details on the task of \Cref{sec:nlp}. Specifically, we give a formal description of the task, provide details on our experimental setting, and present additional experiments using a fine-tuned version of GPT2-large from HuggingFace \cite{radford2019language}. Lastly, we present additional experiments with a varying training sample size for both standard and positional Transformer.

\subsubsection{Task description}
\label{app:nlp_task}
Let $n, k, m\in \mathbb{N}$ with $n\geq k$ and $n \geq m$. The input to the model is a list consisting of text and numerical values of the following form:\\
\begin{center}
\small{
\texttt{['Cat$i_1$', $v_{i_1}$, 'Cat$i_2$', $v_{i_2}$, \dots, 'Cat$i_n$', $v_{i_n}$, 'Find \{type\} of Cat$j_1$, Cat$j_2$, \dots, Cat$j_{k-1}$ and Cat$j_k$']}}\footnote{The ellipses are not part of the prompt}
\end{center}
augmented by adding $m$ more alphanumeric elements at random positions in the list which we call irrelevant structure.
Here, $i_1,\dots,i_n \in \mathbb{N}$, $v_{i_1},\dots,v_{i_n} \in \RR_{\geq 0}$, $j_1,\dots,j_k \in \{i_1,\dots,i_n\}$, and \texttt{\{type\}} is one of \texttt{'min'}, \texttt{'max'}, \texttt{'sum'}. The answer to the query is either the minimum, maximum, or sum of the set $\{v_{j_1},\dots, v_{j_k}\}$ depending on \texttt{\{type\}}. The irrelevant structure is generated by concatenating the string \texttt{'Cat'} with one of the characters in \texttt{\{'+', '-', '\_', '*'\}} and a category identifier. Notice how this setting allows for one model to potentially process multiple query types. In fact, we present one such experiment later.

\paragraph{OOD generalization} For this task, we measure OOD generalization in three ways (simultaneously):
\begin{enumerate}
    \item We test using category identifiers (i.e $i_1, i_2, \dots, i_n$) that the models haven't encountered during training, and 
    \item The range of values $v_1, v_2, \dots, v_n$ used for testing is larger than the range used for training.
    \item The type of irrelevant structure used for testing is different than the type used for training. We detail this difference below.
\end{enumerate}

\subsubsection{Experimental setting}
\label{app:nlp_setting}
We fix $n=8$, $k=4$ and $m=2$. Both standard and positional Transformers consist of 3 layers, two attention heads
per layer, and an embedding dimension of 32. All characters of the non-numeric parts of the input are tokenized and passed through an embedding layer, while the numeric ones are passed through a linear layer. The category identifiers $i_1, i_2, \dots, i_8$  are sampled randomly from $\{1,2,\dots, 20\}$ for training and $\{21, 22, \dots, 40\}$ for testing. The query category identifiers are then sampled randomly from $\{i_1, i_2, \dots, i_8\}$. The values $v_{i_1},\dots, v_{i_8}$ are sampled using the technique of \Cref{sec:experiments} from $[0, 5]$ for training and from $[0, 5c]$ where $c\in\{1,2,\dots, 10\}$ is the scaling factor for testing. We also apply the rejection step of \Cref{sec:ood-regime} when generating testing samples. Furthermore, when sampling a training list we form irrelevant structure by concatenating the string \texttt{'Cat'} with either one of \texttt{'+'} or \texttt{'-'} (chosen at random) and a category identifier that is present in the actual input. For testing, irrelevant structure is formed similarly by concatenating the string \texttt{'Cat'} with either one of \texttt{'\_'} or \texttt{'*'} (chosen at random) and a category identifier that is present in the actual (test) input. In both cases, irrelevant structure is injected into the ``true" list at random positions. We run experiments for the following three variants of the task: 

\begin{enumerate}
    \item One where training and testing prompts consist exclusively of prompts where \texttt{\{type\}} is \texttt{'min'}.
    \item One where training and testing prompts consist exclusively of prompts where \texttt{\{type\}} is \texttt{'max'}.
    \item One where training and testing prompts consist exclusively of prompts where \texttt{\{type\}} is either \texttt{'min'} or \texttt{'max'}. We refer to this experiment as ``multitasking" since it allows a single trained model to process different query types. In fact, the choice of minimum and maximum as the two query types is, in some sense, an ``extreme" case, given the ``opposite" nature of minimum and maximum operations.
\end{enumerate}

Both positional and standard Transformer
models are trained using Adam for 2000 epochs with a sample size of 50,000, a batch size of 1024, and an initial learning rate of $5\cdot 10^{-4}$ which is controlled by a learning rate scheduler.

\subsubsection{Sample complexity Experiments}
\label{app:nlp_sample}
In Figures \ref{fig:nlp_sample_min} to \ref{fig:nlp_sample_multitask} we present the in-distribution and out-of-distribution performance of standard and positional Transformer on the three variations of the mixed-type input task. The conclusions are similar to those of \Cref{app:num_sample}, namely, both models fit the training data but positional Transformer achieves significantly better out-of-distribution performance.

\begin{figure*}[hbt]
    \centering
    \includegraphics[width=0.9\linewidth]{plots/nlp/min/sample.pdf}
    \caption{Training, in-distribution, and out-of-distribution test performance for the minimum variation of the mixed-type input task is shown for standard Transformers (red) and positional Transformers (blue) as a function of the number of training samples (indicated on the x-axis). Models are trained on category values in the range $[0, 5]$ with varying training set sizes. In-distribution testing is performed on the same domain, and out-of-distribution testing is conducted on an extended domain of category values, $[0, 50]$, each using 1,000 samples.}
    \label{fig:nlp_sample_min}
\end{figure*}

\begin{figure*}[hbt]
    \centering
    \includegraphics[width=0.9\linewidth]{plots/nlp/sum/sample.pdf}
    \caption{Training, in-distribution, and out-of-distribution test performance for the sum variation of the mixed-type input task is shown for standard Transformers (red) and positional Transformers (blue) as a function of the number of training samples (indicated on the x-axis). Models are trained on category values in the range $[0, 5]$ with varying training set sizes. In-distribution testing is performed on the same domain, and out-of-distribution testing is conducted on an extended domain of category values, $[0, 50]$, each using 1,000 samples.}
    \label{fig:nlp_sample_sum}
\end{figure*}

\begin{figure*}[hbt]
    \centering
    \includegraphics[width=0.9\linewidth]{plots/nlp/multitask/sample.pdf}
    \caption{Training, in-distribution, and out-of-distribution test performance for the multitask variation of the mixed-type input task is shown for standard Transformers (red) and positional Transformers (blue) as a function of the number of training samples (indicated on the x-axis). Models are trained on category values in the range $[0, 5]$ with varying training set sizes. In-distribution is performed on the same domain, and out-of-distribution testing is conducted on an extended domain of category values, $[0, 50]$, each using 1,000 samples.}
    \label{fig:nlp_sample_multitask}
\end{figure*}

\subsubsection{Experiments with GPT2}
\label{app:nlp_gpt2}
For completeness, and given the nature of the task, we also perform experiments with a fined-tuned version of GPT2-large from HuggingFace \cite{radford2019language} on the above task. For all three variations (min, max, and multitasking), we fix the precision of the input numbers to 4 digits (2
digits for the whole part and 2 digits after the decimal point) and the precision of the output to 4 digits for min and max (2 digits for the whole part and 2 digits after the decimal point) and 5 digits for sum (3 digits for the whole part and 2 digits after the decimal
point). This covers all possible numbers that can be sampled. We then use byte pair encoding for tokenization and fine-tuned using 50,000 training samples for 3 epochs. Since this setting is autoregressive (next-token prediction), there were cases where
the model’s output did not correspond to a real number. As this did not occur frequently enough to be considered a problem, we report the median OOD losses (MSE and MAPE) ignoring samples for which no numeric value could be extracted from the model’s output. The results of this experiment are presented in \Cref{fig:nlp_noise_p_and_gpt}. For easier comparison, we also include the losses of the positional Transformer on the same tasks. 

\begin{figure}[hbt]
    \centering
    \begin{subfigure}{0.9\linewidth}
        \includegraphics[width=\linewidth]{plots/nlp/alltasks_mse_noise_p_and_gpt.pdf}
    \end{subfigure}
    \begin{subfigure}{0.9\linewidth}
        \includegraphics[width=\linewidth]{plots/nlp/alltasks_mape_noise_p_and_gpt.pdf}
    \end{subfigure}
    \caption{This experiment shows ID/OOD losses (measured as mean squared error, MSE, and mean absolute percentage error, MAPE) for fine-tuned GPT2-large (red) and positional Transformers (blue) across all three variations of the mixed-type input task. The x-axis represents the OOD scale factor. The solid line and shaded area denote the median and the region between the 10\textsuperscript{th} and 90\textsuperscript{th} percentiles over ten trials, respectively.}
    \label{fig:nlp_noise_p_and_gpt}
\end{figure}




==================================================
Filename: 995_ood_probability.tex
==================================================
\section{Probability of generating OOD test data in our empirical setting}
\label{app:ood_bound}

Recall that we sample the training and test data in the following way. The training data consists of i.i.d samples whose values are drawn from the range $[-2, 2]$. To ensure diversity, for each training sample, we first select lower and upper bounds $\gamma_l$ and $\gamma_u$ uniformly in $[-2, 2]$, and then for each of the $n$ elements of the training sample, we select its value uniformly from the interval $[\gamma_l, \gamma_u]$. We employ a similar sampling strategy for testing but extend the value range to $[-2c,2c]$, where $c>1$ is the OOD scale factor. Additionally, during the test sampling process, we apply a rejection step to ensure that either $\gamma_l < -2$ or $\gamma_u > 2$, while maintaining $-2c \le \gamma_l \le \gamma_u \le 2c$.

We will compute the probability that a randomly sampled test instance $x \in \mathbb{R}^n$ lies in the domain of the training distribution, i.e., we will compute $\mathbb{P}_{x \sim \mathcal{D}_\text{test}(\mathcal{X})}(x \in [-2,2]^n)$. In particular, we will show that this probability is proportional to $1/nc^2$. Consequently, in our experiments, the majority of the test data lie outside of the domain of the training distribution.

Without loss of generality, let us assume that the training data are sampled within the interval $[-1,1]$ and the test data are sampled within the interval $[-c,c]$, where $c$ is the OOD scale factor. Note that this does not affect the probability that we want to compute. In the test sampling process, when we sample two uniform numbers $\gamma_\ell$ and $\gamma_u$ from $[-c,c]$, exactly one of the following 3 disjoint events can happen.
\begin{itemize}[leftmargin=0.5cm]
  \item {\em Event A.} Exactly one of $\gamma_\ell$ and $\gamma_u$ lies in $[-1,1]$. This happens with probability $2(\frac{c-1}{c})(\frac{1}{c})$.
  \item {\em Event B.} Neither $\gamma_\ell$ nor $\gamma_u$ is inside the interval $[-1,1]$. This happens with probability $(\frac{c-1}{c})^2$.
  \item {\em Event C.} Both $\gamma_\ell$ and $\gamma_u$ are inside the interval $[-1,1]$. This happens with probability $\frac{1}{c^2}$.
\end{itemize}
Our rejection step rejects the samples generated under Event C. Therefore, in our setting when we sample a pair of $\gamma_\ell$ and $\gamma_u$ in order to generate a single instance of the test list, we have that
\begin{equation}
\mathbb{P}(\mbox{Event A} | \mbox{Rejecting Event C}) = \frac{2(c-1)}{c^2-1}, \quad \mathbb{P}(\mbox{Event B} | \mbox{Rejecting Event C}) =\frac{(c-1)^2}{c^2-1}. \label{eq:p_AB}
\end{equation}
We analyze the probability of generating OOD test data under each event. First, let us suppose that Event A happens when we sample $\gamma_\ell$ and $\gamma_u$. This means that either $\gamma_\ell \in [-c,-1)$ or $\gamma_u \in (1,c]$ (but not both). More precisely, the following two sub-events partition Event A:
\begin{itemize}[leftmargin=1cm]
  \item {\em Event A.1.} $\gamma_\ell \in [-1,1]$ and $\gamma_u \in (1,c]$. Given that Event A happens, this sub-event happens with probability 1/2.
  \item {\em Event A.2.} $\gamma_\ell \in [-c,-1)$ and $\gamma_u \in [-1,1]$. Given that Event A happens, this sub-event happens with probability 1/2.
\end{itemize}
By symmetry of the probability distributions, the probability that we wish to compute remains the same under both of the above sub-events. Therefore, let us focus on Event A.1. Suppose that Event A.1 happens, i.e., $\gamma_\ell \in [-1,1]$ and $\gamma_u \in (1,c]$. Conditioning on this event, we know that $\gamma_\ell$ is uniform on $[-1,1]$ and $\gamma_u$ is uniform on $(1,c)$. The probability density functions for $\gamma_\ell$ and $\gamma_u$ are
\[
  f_{\gamma_\ell}(s) = \frac{1}{2}\cdot\mathbb{I}(-1\le s \le 1), \quad f_{\gamma_u}(t) = \frac{1}{c-1}\cdot\mathbb{I}(1< t \le c).
\]
Let $X \in \mathbb{R}^n$ be a random vector whose $i$th coordinate $X_i$ is independently and uniformly sampled from the interval $[\gamma_\ell, \gamma_u]$, then the conditional probability density function for $X$ given $\gamma_\ell,\gamma_u$ is
\[
  f_X(x_1,x_2,\ldots,x_n | \gamma_\ell=s, \gamma_u=t) = \left(\frac{1}{t-s}\right)^n \cdot \mathbb{I}(s \le x_i \le t, \forall i).
\]
Therefore, the joint density function is
\[
  f_{X,\gamma_\ell,\gamma_u}(x_1,\ldots,x_n,s,t) = \frac{1}{2}\frac{1}{(c-1)}\left(\frac{1}{t-s}\right)^n \cdot \mathbb{I}(-1 \le s \le 1, 1 < t \le c, s \le x_i \le t, \forall i).
\]
It follows that
\begin{allowdisplaybreaks}
\begin{align}
  &p_{\mathsf{A},\mathsf{in}} := \mathbb{P}_{X,\gamma_\ell,\gamma_u}\Big(X_i \in [-1,1], \forall i \in [n] \Big| \mbox{Event A}\Big)\nonumber\\
  =~& \int_{s \in [-1,1]}\int_{t \in (1,c]}\int_{x \in [s,1]^n} \frac{1}{2}\frac{1}{(c-1)}\left(\frac{1}{t-s}\right)^n dx\,dt\, ds \nonumber\\
  =~& \frac{1}{2}\frac{1}{(c-1)}\int_{s \in [-1,1]}\int_{t \in (1,c]} \left(\frac{1-s}{t-s}\right)^n dt\, ds\nonumber\\
  =~&\frac{1}{2}\frac{1}{(c-1)}\frac{1}{(n-1)}\int_{s \in [-1,1]} (1-s)\left(1-\left(\frac{1-s}{c-s}\right)^{n-1}\right) ds\nonumber\\
  =~&\frac{1}{2}\frac{1}{(c-1)}\frac{1}{(n-1)}\left[\int_{s \in [-1,0]} (1-s)\left(1-\left(\frac{1-s}{c-s}\right)^{n-1}\right) ds \right.\nonumber\\
  &\hspace{3cm} + \left.\int_{s \in [0,1]} (1-s)\left(1-\left(\frac{1-s}{c-s}\right)^{n-1}\right) ds\right]\nonumber\\
  \le~& \frac{1}{2}\frac{1}{(c-1)}\frac{1}{(n-1)}\left[\int_{s \in [-1,0]} (1-s)\left(1-\left(\frac{1}{c}\right)^{n-1}\right) ds + \int_{s \in [0,1]} (1-s) ds\right]\nonumber\\
  =~&\frac{3\left(1-1/c^{n-1}\right) + 1}{4(n-1)(c-1)}. \label{eq:p_A}
\end{align}
\end{allowdisplaybreaks}

This is the probability that, under Event A, a randomly generated test sample lies within the domain of the training distribution. Again, recall that by scaling down the domain of the test distribution to $[-c,c]$ accordingly, we have assumed that the domain of the training distribution is $[-1,1]^n$ without loss of generalization. 

Now suppose that Event B happens. In this case both $\gamma_\ell$ and $\gamma_u$ are uniformly distributed over $[-c,-1)\cup(1,c]$. The following two sub-events partition Event B:
\begin{itemize}[leftmargin=1cm]
  \item {\em Event B.1.} Either both $\gamma_\ell,\gamma_u > 1$ or both $\gamma_\ell,\gamma_u < -1$. Given that Event B happens, this sub-event happens with probability 1/2.
  \item {\em Event B.2.} $\gamma_\ell < -1$ and $\gamma_u > 1$. Given that Event B happens, this sub-event happens with probability 1/2.
\end{itemize}
Let $X \in \mathbb{R}^n$ be a random vector whose $i$th coordinate $X_i$ is independently and uniformly sampled from the interval $[\gamma_\ell, \gamma_u]$. Note that under Event B.1, one always has that $X \not\in [-1,1]^n$, i.e.,
\begin{equation}\label{eq:p_B1}
  p_{\mathsf{B.1},\mathsf{in}} := \mathbb{P}_{X,\gamma_\ell,\gamma_u}\Big(X_i \in [-1,1], \forall i \in [n] \Big| \mbox{Event B.1}\Big) = 0.
\end{equation}
Therefore let us consider Event B.2. Conditioning on this event, we know that $\gamma_\ell$ is
uniform on $[-c, -1)$ and $\gamma_u$ is uniform on $(1, c]$. The joint density function (conditional on Event B.2) for $X,\gamma_\ell, \gamma_u$ is
\[
    f_{X,\gamma_\ell,\gamma_u}(x_1,\ldots,x_n,s,t) = \frac{1}{(c-1)^2}\left(\frac{1}{t-s}\right)^n \cdot \mathbb{I}(-c \le s \le 1, 1 < t \le c, s \le x_i \le t, \forall i).
\]
Therefore, we have that
\begin{align}
  &p_{\mathsf{B.2},\mathsf{in}} := \mathbb{P}_{X,\gamma_\ell,\gamma_u}\Big(X_i \in [-1,1], \forall i \in [n] \Big| \mbox{Event B.2}\Big)\nonumber\\ 
  =~& \int_{s \in [-c,1)}\int_{t \in (1,c]}\int_{x \in [s,1]^n} \frac{1}{(c-1)^2}\left(\frac{1}{t-s}\right)^n dx\,dt\,ds\nonumber\\
  =~& \frac{1}{(c-1)^2} \int_{s \in [-c,1)}\int_{t \in (1,c]} 2^n\left(\frac{1}{t-s}\right)^n dt\, ds\nonumber\\
  =~& \frac{2^n}{(c-1)^2(n-1)}\int_{s \in [-c,1)} \left(\frac{1}{(1-s)^{n-1}} - \frac{1}{(c-s)^{n-1}}\right)ds\nonumber\\
  =~&\left\{\begin{array}{ll} \dfrac{4 - 8(\frac{2}{1+c})^{n-2} + 4(\frac{1}{c})^{n-2}}{(c-1)^2(n-1)(n-2)}, & \mbox{if}~ n\ge 3,\\ \dfrac{2\log(c+1) - \log c - 2\log 2}{(c-1)^2}, & \mbox{if} ~ n = 2.
  \end{array}\right. \label{eq:p_B2}
\end{align}

Combining \Cref{eq:p_AB}, \Cref{eq:p_A}, \Cref{eq:p_B1}, \Cref{eq:p_B2}, we get that, if $n \ge 3$,
\begin{align}
  p_{\mathsf{in}} := \mathbb{P}_{x\sim\mathcal{D}_\text{test}(\mathcal{X})}(x\in[-1,1]^n) 
  & \le \frac{3(1-1/c^{n-1}) + 1}{2(c^2-1)(n-1)} + \frac{2 - 4(\frac{2}{1+c})^{n-2} + 2(\frac{1}{c})^{n-2}}{(n-1)(n-2)(c^2-1)} \label{eq:ood_prob1}\\
  & \le \frac{2}{(c^2-1)(n-1)} + \frac{2}{(n-1)(n-2)(c^2-1)}\nonumber\\
  & = O\left(\frac{1}{nc^2}\right)\nonumber
\end{align}
and if $n=2$,
\begin{equation}
\label{eq:ood_prob2}
  p_{\mathsf{in}} \le \frac{3(1-1/c) + 1 + 2\log(c+1) - \log c - 2\log 2}{2(c^2-1)} \le \frac{3(1-1/c) + 9/8}{2(c^2-1)}.
\end{equation}
In the above, $p_{\mathsf{in}}$ is the probability that a randomly sampled test list $x \in \mathbb{R}^n$ has all its elements lie within $[-1,1]$, that is, the probability that $x$ lies within the domain of the training distribution. This probability is at most $O(1/nc^2)$. Suppose that we generate $N$ test instances, then a straightforward application of the multiplicative Chernoff bound yields that with probability at least $N^{-C}$ for some constant $C>0$, at most $O(\frac{N}{nc^2})$ samples will lie in the domain of the training distribution.

\begin{figure}[ht!]
  \centering
  \includegraphics[width=0.6\linewidth]{plots/misc/sampling_prob_contour.pdf}
  \caption{Contour plot of $\mathbb{P}_{x\sim\mathcal{D}_\text{test}(\mathcal{X})}(x \in \text{supp}(\mathcal{D}_\text{train}(\mathcal{X})))$, i.e., the probability (upper bound in \Cref{eq:ood_prob1} and \Cref{eq:ood_prob2}) that a randomly sampled test instance $x \in \mathbb{R}^n$ lies in the domain of the training distribution.}
  \label{fig:ood_prod}
\end{figure}

For $n \in \{2,4,8,16,32\}$ and $c \in \{1,2,\ldots,10\}$ which we consider in our experiments, \Cref{fig:ood_prod} shows a contour plot of the probability upper bound \Cref{eq:ood_prob1} and \Cref{eq:ood_prob2}. This probability is sufficiently small such that the majority of test instances in our test data does not belong to the domain of the training distribution.

For small $n$ and $c$, to determine the fraction of sampled test instances that will be within the domain of the training distribution, it is more informative to directly invoke the additive Chernoff bound with $p_{\mathsf{in}}$. Let $N$ denote the total number of test instances that we sample, and further let $N_{\mathsf{in}}$ denote the number of sampled instances that lie in the domain of the training distribution. Then by the additive Chernoff bound we have that
\begin{equation}
\label{eq:ood_prod_chernoff}
  \mathbb{P}(N_{\mathsf{in}} \ge N(p_{\mathsf{in}}+\epsilon)) \le \exp(-2N\epsilon^2).
\end{equation}
For example, suppose that we sample $N=1000$ test instances from the test distribution. Suppose that we generate the test data using list length $n=2$ and OOD scale factor $c=2$. Then in this case $p_{\mathsf{in}} \le 0.4375$. Take $\epsilon = 0.0625$. Then (\ref{eq:ood_prod_chernoff}) says that with probability at least 0.9995, at least $N/2$ samples do not lie in the domain of the training distribution. For another example with slightly larger $n$ and $c$, suppose that we generate the test data using list length $n=8$ and OOD scale factor $c=10$. Then $p_{\mathsf{in}} \le 0.0034$. Take $\epsilon = 0.0466$. Then (\ref{eq:ood_prod_chernoff}) says that with probability at least 0.98, more than 95\% of test instances do not lie in the domain of the train distribution.



==================================================
Filename: 996_failure_attention.tex
==================================================
\section{Potential reasons for failure in self-attention for out-of-distribution}
\label{app:failure}
In this section, we discuss potential reasons for the shortcomings in out-of-distribution generalization of standard Transformers for algorithmic tasks. While it is more straightforward to elicit reasons for the success of positional attention -- motivated by the \emph{algorithmic alignment} \citep{Xu2020What} between positional Transformers and parallel computational models -- it is considerably more challenging to pinpoint the causes of failure in standard Transformers.

Firstly, Transformers can simulate parallel algorithms, as demonstrated by \cite{pmlr-v235-sanford24a}. Intuitively, a single layer of self-attention should be more powerful than positional attention, as it leverages attention beyond positional encodings and allows for a more flexible structure in response to input variations. However, as discussed in \Cref{sec:expressivity}, executing parallel algorithms does not require using anything beyond positional information in attention.

Assuming that standard Transformers should adopt positional information to effectively execute parallel algorithms, the operations required by standard Transformers become increasingly difficult than positional Transformers for two main reasons:
\begin{enumerate}
\item Self-attention layers must learn to ignore input values and exploit positional information.
\item Transformer layers must preserve positional encodings for subsequent layers.
\end{enumerate}

Namely, these desirable properties of positional Transformers present two significant challenges for standard Transformers. The first challenge arises naturally from the differences between standard and positional attention mechanisms. The second challenge highlights the compositional structure of attention layers, which can be detrimental during training. Specifically, the operations performed by each attention and MLP layer can degrade the inputs of subsequent layers.

This issue is further emphasized in \Cref{rem:softmax}, where we state that while positional Transformers can represent any softmax pattern at any layer, standard Transformers may fail to do so due to potential degradation of the attention inputs. Although residual connections can mitigate this issue by preserving input information, they must ensure that no overlaps hinder the use of positional encodings in subsequent layers. Moreover, this problem compounds across layers, making training more difficult as errors in earlier layers adversely affect subsequent computations.

 Nevertheless, these remain speculative reasons for the observed failure of standard Transformers. Determining the exact causes and the difficulty of achieving the two aforementioned goals through training requires a thorough analysis of the training dynamics, which is inherently challenging. Future in-depth work within the mechanistic interpretability framework \citep{nanda2023progress} can potentially shed light on these issues by inspecting network parameters at convergence, thereby uncovering the underlying reasons for the failure of standard Transformers.

 Along this direction, we present some empirical evidence, in Figure~\ref{fig:attn_weights_1} and Figure~\ref{fig:attn_weights_2}, that self-attention layers in a trained Transformer model can be highly sensitive to input values. In particular, the attention weights change dramatically when the input values of the test data do not necessarily lie in the domain of the training data. This suggests that self-attention potentially overfits the training data and, therefore, offers a plausible explanation for why the standard Transformers exhibit such poor value generalization in our experiments. 

\begin{figure}[ht]
  \centering
  \includegraphics[width=0.9\linewidth]{plots/misc/self_attention_weights_sort1.pdf}
  \caption{Visualization of learned attention weights in the standard Transformer model trained to solve the sorting task in our experiments. The input list to the model is $cX$ where $X = [1.75,1.25,0.75,0.25,-0.25,-0.75,-1.25,-1.75]$ and $c = 1,2,\ldots,8$ is a scaling factor. The model is trained on data whose input values range from -2 to 2. Therefore $c=1$ gives in-distribution data and larger $c$ yields OOD data. For each layer in the architecture, we plot 1 of the 2 attention heads for illustration purposes. The trend for the other head is similar. Observe that the attention weights change dramatically as we increase the scaling factor of input values, with deeper layers suffering from more radical changes in the attention pattern under even a small change in the scale (e.g. going from $c=1$ to $c=2$). This behavior potentially explains why the standard Transformer model performs poorly on OOD test data in our experiments.}
  \label{fig:attn_weights_1}
\end{figure}

\begin{figure}[h]
  \centering
  \includegraphics[width=0.9\linewidth]{plots/misc/self_attention_weights_sort2.pdf}
  \caption{Another visualization of learned attention weights in the standard Transformer model. We use the same setting as described in Figure~\ref{fig:attn_weights_1}, except that the input list is $cX$ where $X = [2,2,-2,-2,-2,-2,2,2]$. Again, we observe that the attention weights are highly sensitive to the scaling factor $c$, especially those at deeper layers.}
  \label{fig:attn_weights_2}
\end{figure}

\clearpage


==================================================
Filename: 997_odd_bounds.tex
==================================================
\section{Informal Discussion on Out-of-distribution Generalization Bounds and future work}
\label{app:ood}
The topic of OOD generalization has been studied extensively from a theoretical perspective. Researchers often focus on bounding the risk on the test distribution using a mixture of terms that go to zero as the number of samples increases, plus a divergence term that does not necessarily go to zero depending on the distributions and the hypothesis class. For an extensive survey, we refer the reader to \cite{redko2022surveydomainadaptationtheory}. Although, in general, such bounds offer valuable intuition, they might not be tight for our particular setting. In particular, we examine a popular type of bound found in \cite{35391} which can be large even if the difference in the support of the train and test distributions is the smallest it can be. Note that there are more types of bounds in \cite{redko2022surveydomainadaptationtheory} than the one found in \cite{35391}. Although we have not conducted an in-depth analysis of all cases, we note that all of them depend in a worst-case on the hypothesis class. We believe that to improve upon such generic bounds, one must consider the dynamics of the training procedure for deep positional attention architectures. We find this topic extremely interesting, but we leave it for future work, as it is highly non-trivial.

In what follows, we use a popular example of one of these bounds \cite{35391} and illustrate why it is not tight for a simple task of interest. Briefly, the main issue with this particular OOD bound is that it depends in a worst-case manner on the hypothesis class. We demonstrate this issue using the task of computing the minimum over an array of length $n$. We assume that $n$ is even. For simplicity, we do not work with the cumulative version of the minimum problem, as we did in the main paper. Therefore, the ground truth is simply the minimum of the input array. $\mathcal{D}_\text{train}$ is the train distribution over arrays of length $n$, where each component of the array is sampled independently and uniformly at random from integers in the range $0$ to $L_{\mathcal{D}\text{train}}$, where $L_{\mathcal{D}\text{train}}$ is a constant. $\mathcal{D}_\text{test}$ is the test distribution over arrays of length $n$, where each component of the array is sampled independently and uniformly at random from integers in the range $0$ to $L_{\mathcal{D}_\text{train}} + z$, where $z \geq 0$ is a constant integer. We use an equal number of samples from the train and test distributions, denoted by $m$. The loss function is $\ell(h(x), y) = |h(x) - y|$, where $h$ is a hypothesis, $x$ is an input array of length $n$, and $y = \text{min}(x)$, which is the minimum function over the array $x$. 

The hypothesis class $H$ is the architecture in \Cref{eq:arch} with $\log_2 n$ layers, $2$ heads per layer, and $W_O$ and $W_V$ as the identity matrices for all layers. For positional vectors in general position with dimension $n$, Lemma \ref{lem:hardmax} implies that there exist key and query matrices of size $n \times n$ that can represent any attention pattern at each layer. The MLP at each layer consists of $2$ layers with a hidden dimension equal to $4$. We use the ReLU activation function in all MLPs. This allows the MLP to represent the minimum and maximum functions on two input values exactly. This is because the minimum and maximum functions can be written using ReLUs and linear operations:
\begin{align*}
    \mbox{min}(x_1, x_2) & = \frac{1}{2} (\mbox{ReLU}(x_1+x_2) - \mbox{ReLU}(-x_1-x_2) 
                        - \mbox{ReLU}(x_1-x_2) - \mbox{ReLU}(x_2-x_1))
\end{align*}
and 
\begin{align*}
    \mbox{max}(x_1, x_2) & = \frac{1}{2} (\mbox{ReLU}(x_1+x_2) - \mbox{ReLU}(-x_1-x_2) 
                         + \mbox{ReLU}(x_1-x_2) + \mbox{ReLU}(x_2-x_1)).
\end{align*}
Note that the MLP's ability to represent the minimum function for two inputs exactly is also the reason it can represent the maximum function. In other words, the exact representation of the minimum function comes with the consequence that the MLP can also represent the maximum function for two inputs. This observation is crucial later when we show that an existing popular bound from \cite{35391} is not tight for this particular task.

Furthermore, we assume that $|h(x)|$ is constant, which further implies that the magnitude of the loss is bounded above by a constant. Observe that this hypothesis class can represent a binary tree reduction algorithm for the minimum and maximum functions. This is possible because positional attention can represent any attention pattern, and the MLPs can represent the minimum and maximum functions for two inputs exactly. Specifically, the first layer of the positional attention architecture can represent the connections between layers $0$ (leaf nodes) and $1$ in the binary tree computational graph, the second layer can represent the connections between layers $1$ and $2$, and so on, up to the $\log_2 n$-th layer. The MLPs are used locally at each node of the binary tree to compute the minimum between two input values. Therefore, the minimum and maximum functions over an array of $n$ elements are in the hypothesis class.

Let us now discuss one of the most popular OOD generalization bound results for regression. We will use the third case of Theorem 8 in \cite{35391}, which states the following.
\begin{theorem}[Theorem 8 in \cite{35391}, repurposed for the minimum function task]\label{thm:app:ood_bound}
    Assume that the loss function $\ell$ is symmetric, it obeys the triangle inequality, and it is bounded above by a constant. If the minimum function is in the hypothesis class $H$, then, for any hypothesis $h\in H$, the following holds:
    \begin{align*}
        R_{\mathcal{D}_\text{test}}(h) \le R_{\mathcal{D}_\text{train}}(h) + \mbox{disc}_{\ell}(\mathcal{D}_\text{test},\mathcal{D}_\text{train})
    \end{align*}
    where 
    $$
    \mbox{disc}_{\ell}(\mathcal{D}_\text{test},\mathcal{D}_\text{train}) = \max_{h, h' \in H} \ \left| R_{\mathcal{D}_\text{test}}(h,h') - R_{\mathcal{D}_\text{train}}(h,h')\right|
    $$
    and
    $$
        R_{\mathcal{D}_\text{test}}(h,h') = \mathbb{E}_{(x,y)\sim \mathcal{D}_\text{test}}[\ell(h(x), h'(x))].
    $$
    and $R_{\mathcal{D}_\text{train}}(h,h')$ is defined similarly.
\end{theorem}
The above theorem states that the difference in risk between the test and train distributions is bounded only by the discrepancy term between the two distributions. It is important to note that the discrepancy term depends on the hypothesis class and measures the worst-case difference in the train and test risks within that class. The fact that the discrepancy considers the worst-case difference is why this bound is not tight for our task of computing the minimum function.

Let us now focus on the discrepancy term. Corollary 7 in \cite{35391} states that 
$$
\mbox{disc}_{\ell}(\mathcal{D}_\text{test},\mathcal{D}_\text{train}) \le \mbox{disc}_{\ell}(\hat{\mathcal{D}}_\text{test},\hat{\mathcal{D}}_\text{train}) + 4 (\hat{\mathcal{R}}_{\mathcal{S}_\text{test}}(H) + \hat{\mathcal{R}}_{\mathcal{S}_\text{train}}(H)) + \mathcal{O}\left(\sqrt{\frac{1}{m}} \right),
$$
where $\hat{\mathcal{D}}_\text{test}$ and $\hat{\mathcal{D}}_\text{train}$ are the empirical versions of the test and train distributions, repsectively. We will use the common assumption that these empirical distributions are uniform over the samples. $\mathcal{S}_\text{test}$ and $\mathcal{S}_\text{train}$ are the sample sets for the train and test distributions, respectively. Moreover, $\hat{\mathcal{R}}_{\mathcal{S}_\text{test}}(H)$ and $\hat{\mathcal{R}}_{\mathcal{S}_\text{train}}(H)$ are the empirical Rademacher complexities for the train and test distributions, respectively. It is well-known that the part of the bound corresponding to the empirical Rademacher complexities goes to zero as the number of samples increases. The same holds for the square-root term in the bound. Therefore, the only term left to understand is the discrepancy between the empirical distributions. Let us try to understand how this term behaves. Its definition is:
$$
\mbox{disc}_{\ell}(\hat{\mathcal{D}}_\text{test},\hat{\mathcal{D}}_\text{train}) = \max_{h, h' \in H} \left| \frac{1}{m} \sum_{x \in \mathcal{S}_\text{test}} \ell(h(x), h'(x)) - \frac{1}{m} \sum_{x \in \mathcal{S}_\text{train}} \ell(h(x), h'(x))\right|.
$$
A lower bound of $\mbox{disc}_{\ell}(\hat{\mathcal{D}}_\text{test},\hat{\mathcal{D}}_\text{train})$ is given by setting $h$ to be the minimum function and $h'$ to the maximum function:
$$
\mbox{disc}_{\ell}(\hat{\mathcal{D}}_\text{test},\hat{\mathcal{D}}_\text{train}) \ge \left| \frac{1}{m} \sum_{x \in \mathcal{S}_\text{test}} (\max(x) - \min(x)) - \frac{1}{m} \sum_{x \in \mathcal{S}_\text{train}} (\max(x) - \min(x))\right|
$$
We claim that for polynomial number of samples $m$, e.g., $m=n^c$, where $c$ is a positive integer, there exists $n_0$, such that for all  $n \ge n_0$ we have that $\min(x)=0$ and $\max(x) = L_{\mathcal{D}_\text{train}}$ for all $x \in \mathcal{S}_\text{train}$, and $\min(x)=0$ and $\max(x) = L_{\mathcal{D}_\text{train}} + z$ for all $x \in \mathcal{S}_\text{test}$ with probability at least $0.8$. The proof of this claim is trivial and we provide it below. For now, let us discuss the implications of this claim. We have that 
$\mbox{disc}_{\ell}(\hat{\mathcal{D}}_\text{test},\hat{\mathcal{D}}_\text{train}) \ge 
L_{\mathcal{D}_\text{train}} + z - L_{\mathcal{D}_\text{train}} = z$ with probability at least $0.8$. Therefore, even if $z$ is the smallest it can be such that there is a difference between the train and test distributions in this particular setting, i.e., $z=1$, then the empirical discrepancy is going to be at least $1$. This means that the upper bound of Theorem \ref{thm:app:ood_bound} is at least one, and it is not going to zero as the number of samples increases. This is because the discrepancy definition considers the worst-case scenario without considering the training procedure. In practice, the training procedure may help to discover a hypothesis that is close to the minimum function since the minimum function is part of the hypothesis class. If the hypothesis discovered by the training procedure is close enough to the minimum function, the OOD generalization error may be much smaller than 1. Therefore, depending on the learning task and the hypothesis class, the bound in Theorem \ref{thm:app:ood_bound} can be loose.

Consider the following example, $L_{\mathcal{D}_\text{train}} = 3$ and $z=1$. Therefore, the largest value in the test distribution is $4$. For $m$ and $n$ as noted above, the bound implies that the loss might be up to $1$ for any hypothesis $h$. This further implies that for any hypothesis $h$ the relative error might be up to $25\%$ with probability at least $0.8$, despite the fact that the hypothesis class includes the true function and the training procedure could converge to a good approximation of it.

Let us prove the above probability claim. For $x\in \mathcal{S}_\text{train}$ we have 
\begin{align*}
    \mathbb{P}(x_i \neq 0 \, \text{ for all } i \in [n]) & = \prod_{j=1}^n \mathbb{P}(x_j \neq 0) \\ 
    & = \prod_{j=1}^n \mathbb{P}(x_j \in \{1,2,\dots, L_{\mathcal{D}_\text{train}}\}) \\ 
    & = \prod_{j=1}^n \frac{L_{\mathcal{D}_\text{train}}}{L_{\mathcal{D}_\text{train}}+1} \\ 
    & = \left( \frac{L_{\mathcal{D}_\text{train}}}{L_{\mathcal{D}_\text{train}}+1} \right)^n.
\end{align*}
Similarly, we have that 
$$
\mathbb{P}(x_i \neq L_{\mathcal{D}_\text{train}} \mbox{ for all } i \in [n]) = \left( \frac{L_{\mathcal{D}_\text{train}}}{L_{\mathcal{D}_\text{train}}+1} \right)^n.
$$
Furthermore, we have that 
\begin{align*}
    \mathbb{P}(\exists\, i,j\in [n] : x_i=0 \mbox{ and } x_j = L_{\mathcal{D}_\text{train}}) & = 1 - \mathbb{P}(x_i \neq 0 \,\forall\, i\in [n] \mbox{ or } x_i \neq L_{\mathcal{D}_\text{train}} \, \forall\, i\in [n]) \\ 
    & \ge 1 - \mathbb{P}(x_i \neq 0 \,\forall\, i\in[n]) - \mathbb{P}(x_i \neq L_{\mathcal{D}_\text{train}}\,\forall\, i\in [n]) \\ 
    & = 1 - 2\left( \frac{L_{\mathcal{D}_\text{train}}}{L_{\mathcal{D}_\text{train}}+1} \right)^n.
\end{align*}
Therefore, we conclude that 
$$
\mathbb{P}(\mbox{all samples } x\in \mathcal{S}_\text{train}  \mbox{ have at least one } 0 \mbox{ or } L_{\mathcal{D}_\text{train}}) \ge \left(1 - 2\left( \frac{L_{\mathcal{D}_\text{train}}}{L_{\mathcal{D}_\text{train}}+1} \right)^n\right)^m.
$$
and, similarly, we conclude that
$$
\mathbb{P}(\mbox{all samples } x\in \mathcal{S}_\text{test}  \mbox{ have at least one } 0 \mbox{ or } L_{\mathcal{D}_\text{train}}+z) \ge \left(1 - 2\left( \frac{L_{\mathcal{D}_\text{train}} + z}{L_{\mathcal{D}_\text{train}}+z+1} \right)^n\right)^m.
$$

For a polynomial number of samples $m$, e.g., $m=n^c$, where $c$ is a positive integer, there exists some $n_0 \in \mathbb{N}$, such that for all $n \ge n_0$ we have that the latter two probabilities are at least $0.9$ (since for $m=n^c$ both lower bounds tend to $1$ as $n$ tends to infinity). Therefore, our claim about the minimum and maximum over the sampled arrays holds with probability at least $0.81$.

%
%

%

%
%
%
%
%
%
%
%

%

%
%
%
%
%
%
%
%
%
%
%
%
%
%
%
%
%
%
%

%

%
%
%

%
%
%
%
%
%

%
%
%

%
%
%
%
%

%
%
%
%
%
%

%
%
%
%
%
%
%
%

%

%
%
%
%
%
%
%

%
%
%
%
%
%
%
%
%
%
%

%
%
%
%

%

%
%
%
%
%
%
%

%
%
%
%
%
%
%
%
%
%

%
%
%
%



==================================================
Filename: math_commands.tex
==================================================
%

\usepackage{amsmath,amsfonts,bm}

\usepackage{amsthm}
\theoremstyle{definition}
\newtheorem*{definition*}{Definition}

%
\newcommand{\figleft}{{\em (Left)}}
\newcommand{\figcenter}{{\em (Center)}}
\newcommand{\figright}{{\em (Right)}}
\newcommand{\figtop}{{\em (Top)}}
\newcommand{\figbottom}{{\em (Bottom)}}
\newcommand{\captiona}{{\em (a)}}
\newcommand{\captionb}{{\em (b)}}
\newcommand{\captionc}{{\em (c)}}
\newcommand{\captiond}{{\em (d)}}

%
\newcommand{\newterm}[1]{{\bf #1}}


%
\def\figref#1{figure~\ref{#1}}
%
\def\Figref#1{Figure~\ref{#1}}
\def\twofigref#1#2{figures \ref{#1} and \ref{#2}}
\def\quadfigref#1#2#3#4{figures \ref{#1}, \ref{#2}, \ref{#3} and \ref{#4}}
%
\def\secref#1{section~\ref{#1}}
%
\def\Secref#1{Section~\ref{#1}}
%
\def\twosecrefs#1#2{sections \ref{#1} and \ref{#2}}
%
\def\secrefs#1#2#3{sections \ref{#1}, \ref{#2} and \ref{#3}}
%
\def\eqref#1{equation~\ref{#1}}
%
\def\Eqref#1{Equation~\ref{#1}}
%
\def\plaineqref#1{\ref{#1}}
%
\def\chapref#1{chapter~\ref{#1}}
%
\def\Chapref#1{Chapter~\ref{#1}}
%
\def\rangechapref#1#2{chapters\ref{#1}--\ref{#2}}
%
\def\algref#1{algorithm~\ref{#1}}
%
\def\Algref#1{Algorithm~\ref{#1}}
\def\twoalgref#1#2{algorithms \ref{#1} and \ref{#2}}
\def\Twoalgref#1#2{Algorithms \ref{#1} and \ref{#2}}
%
\def\partref#1{part~\ref{#1}}
%
\def\Partref#1{Part~\ref{#1}}
\def\twopartref#1#2{parts \ref{#1} and \ref{#2}}

\def\ceil#1{\lceil #1 \rceil}
\def\floor#1{\lfloor #1 \rfloor}
\def\1{\bm{1}}
\newcommand{\train}{\mathcal{D}}
\newcommand{\valid}{\mathcal{D_{\mathrm{valid}}}}
\newcommand{\test}{\mathcal{D_{\mathrm{test}}}}

\def\eps{{\epsilon}}


%
\def\reta{{\textnormal{$\eta$}}}
\def\ra{{\textnormal{a}}}
\def\rb{{\textnormal{b}}}
\def\rc{{\textnormal{c}}}
\def\rd{{\textnormal{d}}}
\def\re{{\textnormal{e}}}
\def\rf{{\textnormal{f}}}
\def\rg{{\textnormal{g}}}
\def\rh{{\textnormal{h}}}
\def\ri{{\textnormal{i}}}
\def\rj{{\textnormal{j}}}
\def\rk{{\textnormal{k}}}
\def\rl{{\textnormal{l}}}
%
\def\rn{{\textnormal{n}}}
\def\ro{{\textnormal{o}}}
\def\rp{{\textnormal{p}}}
\def\rq{{\textnormal{q}}}
\def\rr{{\textnormal{r}}}
\def\rs{{\textnormal{s}}}
\def\rt{{\textnormal{t}}}
\def\ru{{\textnormal{u}}}
\def\rv{{\textnormal{v}}}
\def\rw{{\textnormal{w}}}
\def\rx{{\textnormal{x}}}
\def\ry{{\textnormal{y}}}
\def\rz{{\textnormal{z}}}

%
\def\rvepsilon{{\mathbf{\epsilon}}}
\def\rvtheta{{\mathbf{\theta}}}
\def\rva{{\mathbf{a}}}
\def\rvb{{\mathbf{b}}}
\def\rvc{{\mathbf{c}}}
\def\rvd{{\mathbf{d}}}
\def\rve{{\mathbf{e}}}
\def\rvf{{\mathbf{f}}}
\def\rvg{{\mathbf{g}}}
\def\rvh{{\mathbf{h}}}
\def\rvu{{\mathbf{i}}}
\def\rvj{{\mathbf{j}}}
\def\rvk{{\mathbf{k}}}
\def\rvl{{\mathbf{l}}}
\def\rvm{{\mathbf{m}}}
\def\rvn{{\mathbf{n}}}
\def\rvo{{\mathbf{o}}}
\def\rvp{{\mathbf{p}}}
\def\rvq{{\mathbf{q}}}
\def\rvr{{\mathbf{r}}}
\def\rvs{{\mathbf{s}}}
\def\rvt{{\mathbf{t}}}
\def\rvu{{\mathbf{u}}}
\def\rvv{{\mathbf{v}}}
\def\rvw{{\mathbf{w}}}
\def\rvx{{\mathbf{x}}}
\def\rvy{{\mathbf{y}}}
\def\rvz{{\mathbf{z}}}

%
\def\erva{{\textnormal{a}}}
\def\ervb{{\textnormal{b}}}
\def\ervc{{\textnormal{c}}}
\def\ervd{{\textnormal{d}}}
\def\erve{{\textnormal{e}}}
\def\ervf{{\textnormal{f}}}
\def\ervg{{\textnormal{g}}}
\def\ervh{{\textnormal{h}}}
\def\ervi{{\textnormal{i}}}
\def\ervj{{\textnormal{j}}}
\def\ervk{{\textnormal{k}}}
\def\ervl{{\textnormal{l}}}
\def\ervm{{\textnormal{m}}}
\def\ervn{{\textnormal{n}}}
\def\ervo{{\textnormal{o}}}
\def\ervp{{\textnormal{p}}}
\def\ervq{{\textnormal{q}}}
\def\ervr{{\textnormal{r}}}
\def\ervs{{\textnormal{s}}}
\def\ervt{{\textnormal{t}}}
\def\ervu{{\textnormal{u}}}
\def\ervv{{\textnormal{v}}}
\def\ervw{{\textnormal{w}}}
\def\ervx{{\textnormal{x}}}
\def\ervy{{\textnormal{y}}}
\def\ervz{{\textnormal{z}}}

%
\def\rmA{{\mathbf{A}}}
\def\rmB{{\mathbf{B}}}
\def\rmC{{\mathbf{C}}}
\def\rmD{{\mathbf{D}}}
\def\rmE{{\mathbf{E}}}
\def\rmF{{\mathbf{F}}}
\def\rmG{{\mathbf{G}}}
\def\rmH{{\mathbf{H}}}
\def\rmI{{\mathbf{I}}}
\def\rmJ{{\mathbf{J}}}
\def\rmK{{\mathbf{K}}}
\def\rmL{{\mathbf{L}}}
\def\rmM{{\mathbf{M}}}
\def\rmN{{\mathbf{N}}}
\def\rmO{{\mathbf{O}}}
\def\rmP{{\mathbf{P}}}
\def\rmQ{{\mathbf{Q}}}
\def\rmR{{\mathbf{R}}}
\def\rmS{{\mathbf{S}}}
\def\rmT{{\mathbf{T}}}
\def\rmU{{\mathbf{U}}}
\def\rmV{{\mathbf{V}}}
\def\rmW{{\mathbf{W}}}
\def\rmX{{\mathbf{X}}}
\def\rmY{{\mathbf{Y}}}
\def\rmZ{{\mathbf{Z}}}

%
\def\ermA{{\textnormal{A}}}
\def\ermB{{\textnormal{B}}}
\def\ermC{{\textnormal{C}}}
\def\ermD{{\textnormal{D}}}
\def\ermE{{\textnormal{E}}}
\def\ermF{{\textnormal{F}}}
\def\ermG{{\textnormal{G}}}
\def\ermH{{\textnormal{H}}}
\def\ermI{{\textnormal{I}}}
\def\ermJ{{\textnormal{J}}}
\def\ermK{{\textnormal{K}}}
\def\ermL{{\textnormal{L}}}
\def\ermM{{\textnormal{M}}}
\def\ermN{{\textnormal{N}}}
\def\ermO{{\textnormal{O}}}
\def\ermP{{\textnormal{P}}}
\def\ermQ{{\textnormal{Q}}}
\def\ermR{{\textnormal{R}}}
\def\ermS{{\textnormal{S}}}
\def\ermT{{\textnormal{T}}}
\def\ermU{{\textnormal{U}}}
\def\ermV{{\textnormal{V}}}
\def\ermW{{\textnormal{W}}}
\def\ermX{{\textnormal{X}}}
\def\ermY{{\textnormal{Y}}}
\def\ermZ{{\textnormal{Z}}}

%
\def\vzero{{\bm{0}}}
\def\vone{{\bm{1}}}
\def\vmu{{\bm{\mu}}}
\def\vtheta{{\bm{\theta}}}
\def\va{{\bm{a}}}
\def\vb{{\bm{b}}}
\def\vc{{\bm{c}}}
\def\vd{{\bm{d}}}
\def\ve{{\bm{e}}}
\def\vf{{\bm{f}}}
\def\vg{{\bm{g}}}
\def\vh{{\bm{h}}}
\def\vi{{\bm{i}}}
\def\vj{{\bm{j}}}
\def\vk{{\bm{k}}}
\def\vl{{\bm{l}}}
\def\vm{{\bm{m}}}
\def\vn{{\bm{n}}}
\def\vo{{\bm{o}}}
\def\vp{{\bm{p}}}
\def\vq{{\bm{q}}}
\def\vr{{\bm{r}}}
\def\vs{{\bm{s}}}
\def\vt{{\bm{t}}}
\def\vu{{\bm{u}}}
\def\vv{{\bm{v}}}
\def\vw{{\bm{w}}}
\def\vx{{\bm{x}}}
\def\vy{{\bm{y}}}
\def\vz{{\bm{z}}}

%
\def\evalpha{{\alpha}}
\def\evbeta{{\beta}}
\def\evepsilon{{\epsilon}}
\def\evlambda{{\lambda}}
\def\evomega{{\omega}}
\def\evmu{{\mu}}
\def\evpsi{{\psi}}
\def\evsigma{{\sigma}}
\def\evtheta{{\theta}}
\def\eva{{a}}
\def\evb{{b}}
\def\evc{{c}}
\def\evd{{d}}
\def\eve{{e}}
\def\evf{{f}}
\def\evg{{g}}
\def\evh{{h}}
\def\evi{{i}}
\def\evj{{j}}
\def\evk{{k}}
\def\evl{{l}}
\def\evm{{m}}
\def\evn{{n}}
\def\evo{{o}}
\def\evp{{p}}
\def\evq{{q}}
\def\evr{{r}}
\def\evs{{s}}
\def\evt{{t}}
\def\evu{{u}}
\def\evv{{v}}
\def\evw{{w}}
\def\evx{{x}}
\def\evy{{y}}
\def\evz{{z}}

%
\def\mA{{\bm{A}}}
\def\mB{{\bm{B}}}
\def\mC{{\bm{C}}}
\def\mD{{\bm{D}}}
\def\mE{{\bm{E}}}
\def\mF{{\bm{F}}}
\def\mG{{\bm{G}}}
\def\mH{{\bm{H}}}
\def\mI{{\bm{I}}}
\def\mJ{{\bm{J}}}
\def\mK{{\bm{K}}}
\def\mL{{\bm{L}}}
\def\mM{{\bm{M}}}
\def\mN{{\bm{N}}}
\def\mO{{\bm{O}}}
\def\mP{{\bm{P}}}
\def\mQ{{\bm{Q}}}
\def\mR{{\bm{R}}}
\def\mS{{\bm{S}}}
\def\mT{{\bm{T}}}
\def\mU{{\bm{U}}}
\def\mV{{\bm{V}}}
\def\mW{{\bm{W}}}
\def\mX{{\bm{X}}}
\def\mY{{\bm{Y}}}
\def\mZ{{\bm{Z}}}
\def\mBeta{{\bm{\beta}}}
\def\mPhi{{\bm{\Phi}}}
\def\mLambda{{\bm{\Lambda}}}
\def\mSigma{{\bm{\Sigma}}}

%
\DeclareMathAlphabet{\mathsfit}{\encodingdefault}{\sfdefault}{m}{sl}
\SetMathAlphabet{\mathsfit}{bold}{\encodingdefault}{\sfdefault}{bx}{n}
\newcommand{\tens}[1]{\bm{\mathsfit{#1}}}
\def\tA{{\tens{A}}}
\def\tB{{\tens{B}}}
\def\tC{{\tens{C}}}
\def\tD{{\tens{D}}}
\def\tE{{\tens{E}}}
\def\tF{{\tens{F}}}
\def\tG{{\tens{G}}}
\def\tH{{\tens{H}}}
\def\tI{{\tens{I}}}
\def\tJ{{\tens{J}}}
\def\tK{{\tens{K}}}
\def\tL{{\tens{L}}}
\def\tM{{\tens{M}}}
\def\tN{{\tens{N}}}
\def\tO{{\tens{O}}}
\def\tP{{\tens{P}}}
\def\tQ{{\tens{Q}}}
\def\tR{{\tens{R}}}
\def\tS{{\tens{S}}}
\def\tT{{\tens{T}}}
\def\tU{{\tens{U}}}
\def\tV{{\tens{V}}}
\def\tW{{\tens{W}}}
\def\tX{{\tens{X}}}
\def\tY{{\tens{Y}}}
\def\tZ{{\tens{Z}}}


%
\def\gA{{\mathcal{A}}}
\def\gB{{\mathcal{B}}}
\def\gC{{\mathcal{C}}}
\def\gD{{\mathcal{D}}}
\def\gE{{\mathcal{E}}}
\def\gF{{\mathcal{F}}}
\def\gG{{\mathcal{G}}}
\def\gH{{\mathcal{H}}}
\def\gI{{\mathcal{I}}}
\def\gJ{{\mathcal{J}}}
\def\gK{{\mathcal{K}}}
\def\gL{{\mathcal{L}}}
\def\gM{{\mathcal{M}}}
\def\gN{{\mathcal{N}}}
\def\gO{{\mathcal{O}}}
\def\gP{{\mathcal{P}}}
\def\gQ{{\mathcal{Q}}}
\def\gR{{\mathcal{R}}}
\def\gS{{\mathcal{S}}}
\def\gT{{\mathcal{T}}}
\def\gU{{\mathcal{U}}}
\def\gV{{\mathcal{V}}}
\def\gW{{\mathcal{W}}}
\def\gX{{\mathcal{X}}}
\def\gY{{\mathcal{Y}}}
\def\gZ{{\mathcal{Z}}}

%
\def\sA{{\mathbb{A}}}
\def\sB{{\mathbb{B}}}
\def\sC{{\mathbb{C}}}
\def\sD{{\mathbb{D}}}
%
%
\def\sF{{\mathbb{F}}}
\def\sG{{\mathbb{G}}}
\def\sH{{\mathbb{H}}}
\def\sI{{\mathbb{I}}}
\def\sJ{{\mathbb{J}}}
\def\sK{{\mathbb{K}}}
\def\sL{{\mathbb{L}}}
\def\sM{{\mathbb{M}}}
\def\sN{{\mathbb{N}}}
\def\sO{{\mathbb{O}}}
\def\sP{{\mathbb{P}}}
\def\sQ{{\mathbb{Q}}}
\def\sR{{\mathbb{R}}}
\def\sS{{\mathbb{S}}}
\def\sT{{\mathbb{T}}}
\def\sU{{\mathbb{U}}}
\def\sV{{\mathbb{V}}}
\def\sW{{\mathbb{W}}}
\def\sX{{\mathbb{X}}}
\def\sY{{\mathbb{Y}}}
\def\sZ{{\mathbb{Z}}}

%
\def\emLambda{{\Lambda}}
\def\emA{{A}}
\def\emB{{B}}
\def\emC{{C}}
\def\emD{{D}}
\def\emE{{E}}
\def\emF{{F}}
\def\emG{{G}}
\def\emH{{H}}
\def\emI{{I}}
\def\emJ{{J}}
\def\emK{{K}}
\def\emL{{L}}
\def\emM{{M}}
\def\emN{{N}}
\def\emO{{O}}
\def\emP{{P}}
\def\emQ{{Q}}
\def\emR{{R}}
\def\emS{{S}}
\def\emT{{T}}
\def\emU{{U}}
\def\emV{{V}}
\def\emW{{W}}
\def\emX{{X}}
\def\emY{{Y}}
\def\emZ{{Z}}
\def\emSigma{{\Sigma}}

%
%
\newcommand{\etens}[1]{\mathsfit{#1}}
\def\etLambda{{\etens{\Lambda}}}
\def\etA{{\etens{A}}}
\def\etB{{\etens{B}}}
\def\etC{{\etens{C}}}
\def\etD{{\etens{D}}}
\def\etE{{\etens{E}}}
\def\etF{{\etens{F}}}
\def\etG{{\etens{G}}}
\def\etH{{\etens{H}}}
\def\etI{{\etens{I}}}
\def\etJ{{\etens{J}}}
\def\etK{{\etens{K}}}
\def\etL{{\etens{L}}}
\def\etM{{\etens{M}}}
\def\etN{{\etens{N}}}
\def\etO{{\etens{O}}}
\def\etP{{\etens{P}}}
\def\etQ{{\etens{Q}}}
\def\etR{{\etens{R}}}
\def\etS{{\etens{S}}}
\def\etT{{\etens{T}}}
\def\etU{{\etens{U}}}
\def\etV{{\etens{V}}}
\def\etW{{\etens{W}}}
\def\etX{{\etens{X}}}
\def\etY{{\etens{Y}}}
\def\etZ{{\etens{Z}}}

%
\newcommand{\pdata}{p_{\rm{data}}}
%
\newcommand{\ptrain}{\hat{p}_{\rm{data}}}
\newcommand{\Ptrain}{\hat{P}_{\rm{data}}}
%
\newcommand{\pmodel}{p_{\rm{model}}}
\newcommand{\Pmodel}{P_{\rm{model}}}
\newcommand{\ptildemodel}{\tilde{p}_{\rm{model}}}
%
\newcommand{\pencode}{p_{\rm{encoder}}}
\newcommand{\pdecode}{p_{\rm{decoder}}}
\newcommand{\precons}{p_{\rm{reconstruct}}}

\newcommand{\laplace}{\mathrm{Laplace}} %

\newcommand{\E}{\mathbb{E}}
\newcommand{\Ls}{\mathcal{L}}
\newcommand{\R}{\mathbb{R}}
\newcommand{\emp}{\tilde{p}}
\newcommand{\lr}{\alpha}
\newcommand{\reg}{\lambda}
\newcommand{\rect}{\mathrm{rectifier}}
\newcommand{\softmax}{\mathrm{softmax}}
\newcommand{\sigmoid}{\sigma}
\newcommand{\softplus}{\zeta}
\newcommand{\KL}{D_{\mathrm{KL}}}
\newcommand{\Var}{\mathrm{Var}}
\newcommand{\standarderror}{\mathrm{SE}}
\newcommand{\Cov}{\mathrm{Cov}}
%
%
%
\newcommand{\normlzero}{L^0}
\newcommand{\normlone}{L^1}
\newcommand{\normltwo}{L^2}
\newcommand{\normlp}{L^p}
\newcommand{\normmax}{L^\infty}

\newcommand{\parents}{Pa} %

\DeclareMathOperator*{\argmax}{arg\,max}
\DeclareMathOperator*{\argmin}{arg\,min}

\DeclareMathOperator{\sign}{sign}
\DeclareMathOperator{\Tr}{Tr}
\let\ab\allowbreak


